% concatenated from example_paper2.tex, introduction.tex, background.tex, res1.tex, margin.tex, low_rank.tex, experiment_multi.tex, experiment_use_case.tex, low_rank_results_short.tex, table_small.tex, conclusion.tex, impact.tex, appendix_proofs.tex, appendix.tex
%%%%%%%% ICML 2026 EXAMPLE LATEX SUBMISSION FILE %%%%%%%%%%%%%%%%%

\documentclass{article}

% Recommended, but optional, packages for figures and better typesetting:
\usepackage{microtype}
\usepackage{graphicx}
\usepackage{subcaption}
\usepackage{booktabs} % for professional tables

% hyperref makes hyperlinks in the resulting PDF.
% If your build breaks (sometimes temporarily if a hyperlink spans a page)
% please comment out the following usepackage line and replace
% \usepackage{icml2026} with \usepackage[nohyperref]{icml2026} above.
\usepackage{hyperref}


% Attempt to make hyperref and algorithmic work together better:
\newcommand{\theHalgorithm}{\arabic{algorithm}}

% Use the following line for the initial blind version submitted for review:
% \usepackage{icml2026}

% For preprint, use
% \usepackage[preprint]{icml2026}

% If accepted, instead use the following line for the camera-ready submission:
\usepackage[accepted]{icml2026}

\usepackage{amsmath}
\usepackage{amssymb}
\usepackage{mathtools}
\usepackage{amsthm}
\usepackage{amsfonts}
\usepackage{longtable}
\usepackage{graphicx} % Required for inserting images
\usepackage[english]{babel}
\usepackage{booktabs,tabularx,makecell,xcolor,colortbl}
\newcolumntype{Y}{>{\centering\arraybackslash}X}
\definecolor{asryellow}{HTML}{FFFC9E}
\newcommand{\ms}[2]{\makecell[c]{#1\\(#2)}}
\usepackage{pdflscape}
\usepackage{soul}
\usepackage{booktabs}
\usepackage{siunitx}
\usepackage{makecell}
\newcommand{\cmark}{\textcolor{green!60!black}{\ding{51}}}
\newcommand{\xmark}{\textcolor{red!60!black}{\ding{55}}}
\newcommand{\good}{\cellcolor{green!15}\makebox[1.0em][c]{\cmark}}
\newcommand{\bad}{\cellcolor{red!15}\makebox[1.0em][c]{\xmark}}
\usepackage{pifont}
\usepackage{subcaption}
\usepackage{algorithm}
\usepackage{algpseudocode}

% if you use cleveref..
\usepackage[capitalize,noabbrev]{cleveref}

%%%%%%%%%%%%%%%%%%%%%%%%%%%%%%%%
% THEOREMS
%%%%%%%%%%%%%%%%%%%%%%%%%%%%%%%%
\theoremstyle{plain}
\newtheorem{theorem}{Theorem}[section]
\newtheorem{proposition}[theorem]{Proposition}
\newtheorem{lemma}[theorem]{Lemma}
\newtheorem{model}[theorem]{Model}
\newtheorem{corollary}[theorem]{Corollary}
\theoremstyle{definition}
\newtheorem{definition}[theorem]{Definition}
\newtheorem{assumption}[theorem]{Assumption}
\theoremstyle{remark}
\newtheorem{remark}[theorem]{Remark}

\usepackage[textsize=tiny]{todonotes}
\usepackage{amsfonts}
\usepackage{longtable}
\usepackage{graphicx} % Required for inserting images
\usepackage[english]{babel}
\usepackage{booktabs,tabularx,makecell,xcolor,colortbl,subcaption}
\newcolumntype{Y}{>{\centering\arraybackslash}X}
% \definecolor{asryellow}{HTML}{FFFC9E}
% \newcommand{\ms}[2]{\makecell[c]{#1\\(#2)}}
\usepackage{pdflscape}
\usepackage{soul}
\usepackage{booktabs}
\usepackage{siunitx}
\usepackage{makecell}
\definecolor{tableShade}{gray}{0.92}
% \newcommand{\good}{\textcolor{green!60!black}{$\checkmark$}}
% \newcommand{\bad}{\textcolor{red!70!black}{$\times$}}
\usepackage{hyperref}
\usepackage{soul}
\usepackage{listings}
\usepackage{multirow}
\usepackage{natbib}
\usepackage{float} 
\usepackage{graphicx}
% The \icmltitle you define below is probably too long as a header.
% Therefore, a short form for the running title is supplied here:
% \icmltitlerunning{Submission and Formatting Instructions for ICML 2026}
\icmltitlerunning{Theory of Minimal Weight Perturbations in Deep Networks and its Applications for Low-Rank Activated Backdoor Attacks.}

\begin{document}

\twocolumn[
  \icmltitle{Theory of Minimal Weight Perturbations in Deep Networks and its Applications for Low-Rank Activated Backdoor Attacks}

  % It is OKAY to include author information, even for blind submissions: the
  % style file will automatically remove it for you unless you've provided
  % the [accepted] option to the icml2026 package.

  % List of affiliations: The first argument should be a (short) identifier you
  % will use later to specify author affiliations Academic affiliations
  % should list Department, University, City, Region, Country Industry
  % affiliations should list Company, City, Region, Country

  % You can specify symbols, otherwise they are numbered in order. Ideally, you
  % should not use this facility. Affiliations will be numbered in order of
  % appearance and this is the preferred way.
  \icmlsetsymbol{equal}{*}

  \begin{icmlauthorlist}
    \icmlauthor{Bethan Evans}{yyy}
    \icmlauthor{Jared Tanner}{yyy}
    % \icmlauthor{Firstname3 Lastname3}{comp}
    % \icmlauthor{Firstname4 Lastname4}{sch}
    % \icmlauthor{Firstname5 Lastname5}{yyy}
    % \icmlauthor{Firstname6 Lastname6}{sch,yyy,comp}
    % \icmlauthor{Firstname7 Lastname7}{comp}
    % %\icmlauthor{}{sch}
    % \icmlauthor{Firstname8 Lastname8}{sch}
    % \icmlauthor{Firstname8 Lastname8}{yyy,comp}
    %\icmlauthor{}{sch}
    %\icmlauthor{}{sch}
  \end{icmlauthorlist}

  \icmlaffiliation{yyy}{Department of Mathematics, University of Oxford, Oxford, UK}
  % \icmlaffiliation{comp}{Company Name, Location, Country}
  % \icmlaffiliation{sch}{School of ZZZ, Institute of WWW, Location, Country}

  \icmlcorrespondingauthor{Bethan Evans}{bethan.evans@maths.ox.ac.uk}
    \icmlcorrespondingauthor{Jared Tanner}{tanner@maths.ox.ac.uk}


  % You may provide any keywords that you find helpful for describing your
  % paper; these are used to populate the "keywords" metadata in the PDF but
  % will not be shown in the document
  \icmlkeywords{Low-Rank Approximation, Pruning, Adversarial attack, Adversarial robustness, Backdoor attack}

  \vskip 0.3in
]

% this must go after the closing bracket ] following \twocolumn[ ...

% This command actually creates the footnote in the first column listing the
% affiliations and the copyright notice. The command takes one argument, which
% is text to display at the start of the footnote. The \icmlEqualContribution
% command is standard text for equal contribution. Remove it (just {}) if you
% do not need this facility.

% Use ONE of the following lines. DO NOT remove the command.
% If you have no special notice, KEEP empty braces:
\printAffiliationsAndNotice{}  % no special notice (required even if empty)
% Or, if applicable, use the standard equal contribution text:
% \printAffiliationsAndNotice{\icmlEqualContribution}

\begin{abstract}
%Formulae for the minimal weight perturbations of MLPs are derived for a specified change in output, and their components discussed.  These single layer exact formulae are contrasted with more generic multi-layer Lipshitz constant based robustness guarantees; both are observed to be of the same order.  These results are applied to precision-modification-activated back-door attacks, establishing a provable compression threshold below which such attacks cannot succeed and show empirically that low-rank compression can reliably activate latent backdoors while preserving full-precision accuracy.

% We study the precise relationship between weight perturbations and output changes in deep neural networks.
% For multilayer architectures with locally invertible activations, we derive closed-form expressions for the minimum single-layer weight perturbation required to induce a prescribed output change.
% These expressions reveal how backpropagated margins govern layer-wise sensitivity and provide
% certifiable guarantees on the smallest parameter updates consistent with a desired output shift.
% We further derive a Lipschitz-based bound that yields worst-case robustness guarantees to weight
% modifications.
% Applying these results to precision-modification-activated backdoor attacks, we establish a
% provable compression threshold below which such attacks cannot succeed, and show empirically that
% low-rank compression can reliably activate latent backdoors while preserving full-precision
% accuracy.

%Formulae for the minimal weight perturbations of DNNs to achieve a specified change in output are derived, and their components discussed. 

The minimal norm weight perturbations of DNNs required to achieve a specified change in output are derived and the factors determining its size are discussed. These single-layer exact formulae are contrasted with more generic multi-layer Lipschitz constant based robustness guarantees; both are observed to be of the same order which indicates similar efficacy in their guarantees. These results are applied to precision-modification-activated backdoor attacks, establishing provable compression thresholds below which such attacks cannot succeed, and show empirically that low-rank compression can reliably activate latent backdoors while preserving full-precision accuracy. These expressions reveal how back-propagated margins govern layer-wise sensitivity and provide certifiable guarantees on the smallest parameter updates consistent with a desired output shift.
\end{abstract}


\section{Introduction}

Modern deep neural networks (DNNs) often require substantial computational and memory resources \cite{strubell_energy_2019}, motivating widespread adoption of post-training \emph{network compression techniques} \cite{fu_cpt_2025} applied to pre-trained weights~$\theta$, such as quantization \cite{fiesler_weight_1990,zhang_lq-nets_2018}, pruning \cite{han_learning_2015}, and low-rank projection \cite{banerjee_linear_2014}.  
These parameter modification maps $g\colon\theta\mapsto\hat{\theta}$ reduce computational cost and are generally assumed to cause only small, controlled deviations in the parameters $\|\theta-\hat{\theta}\|$ and correspondingly minor changes in the model’s outputs \cite{hubara_quantized_2016}. However, even small weight perturbations can affect output predictions and so understanding how these modifications propagate through the network is crucial for understanding model stability, robustness, and generalisation.


We develop a theoretical framework that quantifies the relationship between weight perturbations and output changes, beginning by deriving exact closed-form expressions for the minimal weight perturbation required to achieve a specified change in the network’s output.  
This exact solution, obtained for multilayer feedforward networks with locally invertible downstream maps, highlights how the backpropagated margin at the perturbed layer determines the required magnitude and direction of the change.  
We then establish a more general but less precise bound on the perturbation norm based on the network’s Lipschitz constant with respect to its weights and the current classification margin.  
While the exact formulation provides sharp, layer-specific characterisations, the Lipschitz-based analysis extends to arbitrary architectures, general non-linear activations, and allows perturbations to occur in multiple layers simultaneously.

Building on our exact perturbation bounds, we analyse low-rank approximation as a structured weight perturbation and derive closed-form expressions for the resulting output deviations. Applying these results to precision-modification–activated backdoors, we show theoretically—and provide empirical evidence —that low-rank compression can activate hidden behaviours previously observed only under pruning or quantization \cite{tian_stealthy_2021,hong_qu-anti-zation_2021}.
These results provide a comprehensive picture of how network weight perturbations interact with both network structure and model compression, providing a theoretical basis for quantifying how large a weight change is needed to alter predictions and for reasoning about the safety of efficiency-driven model compression.


\section{Preliminaries}\label{background}
Initially, Section \ref{layer_chap} and \ref{margin_chap} focus on classification tasks where a model maps an input—such as a point or an image—to a class label; examples for LLMs are given in Sec. \ref{use case}. 
Given an input $\mathbf{x}$, the model outputs a vector of real-valued scores $\mathbf{y}$, called \emph{logits}, with one score for each class. The predicted class is obtained by applying the $\arg \max$ operator, which selects the index of the largest logit.

\begin{model}[Feedforward network and layerwise decomposition]\label{model}
Fix input dimension $d$ and output dimension $c$. 
An $M$-layer feedforward network with parameters
\[
\theta=\{W_l\}_{l=1}^M,\qquad W_l\in\mathbb{R}^{d_l\times d_{l-1}},
\]
where $d_0=d$ and $d_M=c$, and with elementwise nonlinearity
\(
\sigma:\mathbb{R}\to\mathbb{R},
\)
is defined for a matrix of $s$ input samples
\[
X=\bigl[\mathbf{x}^{(1)},\dots,\mathbf{x}^{(s)}\bigr]\in\mathbb{R}^{d\times s}
\]
by the layerwise recursion
\[
Z_0 := X,\qquad
Z_l := \sigma\bigl(W_l Z_{l-1}\bigr),\quad l=1,\dots,M-1,
\]
and the final layer pre-softmax (logit) output
\[
h(X;\theta) := W_M Z_{M-1}\in\mathbb{R}^{c\times s}.
\]
% Including bias terms in each layer does not affect the theoretical results; since all of the pretrained models used here omit biases, we omit them from the notation for simplicity.

For any layer index $N\in\{1,\dots,M\}$, define the \emph{upstream map} as the composition of upstream maps and functions before the $N$th parameter weight matrix by
\[
h_{1:N-1}(X)
:= Z_{N-1}
\in\mathbb{R}^{d_{N-1}\times s},
\]
and the \emph{downstream map} as the composition of downstream maps and functions after the $N$th parameter weight matrix by   
\[
h_{N:M}(Z)
:= W_M\,\sigma\bigl(\cdots \sigma(W_{N+1} \sigma(Z))\cdots\bigr) \textrm{ \, for \, } N<M;
\] 
and $h_{M:M}(Z):=Z$. Then the network factors as
\[
h(X;\theta) \;=\; h_{N:M}\bigl(W_N h_{1:N-1}(X)\bigr).
\]
\end{model}
Including bias terms in each layer does not affect the theoretical results. Biases can be absorbed into the upstream and downstream maps $h_{1:N-1}$ and $h_{N:M}$, and are omitted from the notation for clarity.

For a single input $\mathbf{x}\in\mathbb{R}^d$ (i.e., $s=1$), we write the corresponding logit vector as $\mathbf{y}=h(\mathbf{x};\theta)\in\mathbb{R}^c$, and the predicted class is given by $\arg\max_i\,\mathbf{y}_i$.
% \begin{model}[Feedforward network and layerwise decomposition]\label{model}
% Fix input dimension $d$ and output dimension $c$. An $M$-layer feedforward network with parameters
% \[
% \theta=\{(W_l,b_l)\}_{l=1}^M,\qquad
% W_l\in\mathbb{R}^{d_l\times d_{l-1}},\quad
% b_l\in\mathbb{R}^{d_l},
% \]
% where $d_0=d$ and $d_M=c$, and with elementwise activations
% $\sigma_l:\mathbb{R}^{d_l\times p}\to\mathbb{R}^{d_l\times p}$ for $l=1,\dots,M-1$, acts on a matrix of inputs
% \[
% X=\bigl[\mathbf{x}^{(1)},\dots,\mathbf{x}^{(p)}\bigr]\in\mathbb{R}^{d\times p}.
% \]
% Let $\mathbf{1}_p\in\mathbb{R}^p$ denote the all-ones vector. Define the layerwise recursion
% \[
% Z_0 := X,\qquad
% Z_l := \sigma_l\bigl(W_l Z_{l-1} + b_l\,\mathbf{1}_p^\top\bigr),\ \ l=1,\dots,M-1,
% \]
% and the pre-softmax (logit) output
% \[
% h(X;\theta) := Z_M := W_M Z_{M-1} + b_M\,\mathbf{1}_p^\top \in \mathbb{R}^{c\times p}.
% \]
% % Define the affine–nonlinear blocks by
% % \begin{align}
% % \Phi_l(Z) &:= \sigma_l\!\bigl(W_l Z + b_l\,\mathbf{1}_p^\top\bigr)\ (l=1,\dots,M-1),\\
% % \Phi_M(Z) &:= W_M Z + b_M\,\mathbf{1}_p^\top,
% % \end{align}
% % so $h(X;\theta) = (\Phi_M\circ\Phi_{M-1}\circ\cdots\circ\Phi_1)(X)$.
% For any layer index $N\in\{1,\dots,M\}$, we define the \emph{upstream map} as the composition of layers up to the $N$th weight parameter given by
% \[
% h_{1:N-1}(X) := Z_{N-1}\in\mathbb{R}^{d_{N-1}\times p},
% \]
% and the \emph{downstream map} as the composition of layers after the $N$th weight parameter given by
% \[
% h_{N:M}(Z) := (\Phi_M\circ\cdots\circ\Phi_{N+1}\circ\sigma_N)(Z_N)\in\mathbb{R}^{c\times p}.
% \]
% Then the network may be written as
% \[
% h(X;\theta) \;=\; h_{N:M}\!\bigl( W_N h_{1:N-1}(X) + b_N\,\mathbf{1}_p^\top )\bigr).
% \]

% \end{model}

% :
% for $\textbf{x}\in\mathbb{R}^n$ and $p\ge 1$,
% \[
% \|\textbf{x}\|_{p} \; :=\; \Bigl(\sum_{i=1}^n |\textbf{x}_i|^{\,p}\Bigr)^{\!1/p}.
% \]
% Let $\theta=\{W_l, b_l\}_{l=1}^M$ be the set of network parameters (all weight matrices and bias vectors) and $\Delta\theta=\{\Delta W_l,\Delta b_l\}_{l=1}^M$ be the set of perturbations of $\theta$.

In this article we view the entire parameter set  \(\theta\)  as a single vector obtained by flattening and concatenating each tensor. For $s\geq 1$  we define the parameter perturbation norm  to be \[
\|\Delta\theta\|_p
=\Bigl\|\,[\;\mathrm{vec}(\Delta W_1);\;\mathrm{vec}(\Delta W_2);\;\dots\;]\Bigr\|_p,
\]
where \(\mathrm{vec}(A)\) denotes the column‐stacking of matrix \(A\).

If only a \emph{single} weight matrix $W_N$ is perturbed, then
\(
\|\Delta\theta\|_{2}=\|\Delta W_N\|_{F}.
\)


% \begin{definition}[Lipschitz continuity with respect to the weights \cite{rudin_principles_1976}]
% For a fixed input \(X\) and \(p\ge1\), we say that \(h\) is \(L_{\theta}\)-Lipschitz in its parameters under the \(\ell_p\)-norm  if there exists \(L_{\theta}\ge0\) such that for all parameter sets \(\theta,\hat{\theta}\),
% \[
% \bigl\|h(X;\theta)-h(X;\hat{\theta})\bigr\|_p
% \;\le\;
% L_{\theta}\,\|\theta-\hat{\theta}\|_p.
% \]
% The smallest such constant,
% \[
% L_{\theta}
% \;=\;
% \sup_{\theta\neq\hat{\theta}}
% \frac{\|h(X;\theta)-h(X;\hat{\theta})\|_p}
%      {\|\theta-\hat{\theta}\|_p},
% \]
% is called the \emph{parameter‐space Lipschitz constant} of \(h\) at \(X\).  
% \end{definition}

\subsection{Margins at the output and inside the network}
We distinguish a scalar margin at the logits (the usual classification margin) from a matrix-valued quantity inside the network that measures how far the layer-$N$ representation must move to realise a prescribed change in output.

\begin{definition}[Output logit margin]
Define the \emph{network output margin} of a single input sample $\mathbf{x}$ as the difference between the logit
corresponding to the true class and the largest logit among the remaining classes.
More formally, let $t$ denote the index of the true class label of $\mathbf{x}$, and define the classification
margin as
\[
\gamma(\mathbf{x};\theta)
:= h(\mathbf{x};\theta)_{t} - \max_{i \neq t} h(\mathbf{x};\theta)_{i}.
\]
Equivalently, let $\mathbf{e}_i$ denote the $i$th standard basis vector in $\mathbb{R}^c$, and define
\(
p := \arg\max_{i \neq t} h(\mathbf{x};\theta)_i .
\)
Then the margin may be written as
\[
\gamma(\mathbf{x};\theta)
= (\mathbf{e}_t - \mathbf{e}_p)^\top h(\mathbf{x};\theta).
\]
Large $\gamma(\mathbf x;\theta)$ means that the current prediction is confident, whereas small $\gamma( \mathbf x;\theta)$ indicates that the prediction can flip with a small output disturbance. Inputs correctly classified with parameters $\theta$ are indicated by $\gamma(\mathbf x;\theta)>0$ and perturbed network parameters $\hat{\theta}$ that induce misclassification are indicated by $\gamma( \mathbf x,\hat{\theta})<0$.
\end{definition}

\begin{definition}[Layer-$N$ pre-image difference]
Recall that $h_{N:M}$ denotes the downstream map from layer $N$ to the final output logits, and suppose it is locally invertible at the current layer-$N$ representation with inverse branch $h_{N:M}^{-1}$ defined on a neighbourhood of $\mathbf y=h(\mathbf x;\theta)$ and on a target logit vector $\tilde{\mathbf y}$. Then define the $N$th-layer pre-image difference as
\[
\Delta h_{N:M}^{-1}(\tilde{\mathbf y},\mathbf y ;\theta)
:= h_{N:M}^{-1}(\tilde{\mathbf y};\theta)\;-\;h_{N:M}^{-1}(\mathbf y;\theta)
\in \mathbb{R}^{d_N\times 1}.
\]
For a batch of samples, $\Delta h_{N:M}^{-1}$ stacks these per-sample columns and is matrix-valued. This matrix is the change in input at layer-$N$ required to change from output $\mathbf y$ to $\tilde{\mathbf y}$.
\end{definition}
% $\Delta h_{N:M}^{-1}$ is the \emph{preimage difference}: the amount the layer-$N$ representation must move so that the downstream network produces the desired change at the output.
% $\Delta h_{N:M}^{-1}$ is the \emph{preimage difference}: the amount the layer-$N$ representation must move so that the downstream network produces the desired change at the output.



% \begin{definition}
%     Define the \emph{network classification margin} of a single sample $\mathbf x$ as the difference between the true class output logit and the next largest output logit. More formally, denote by $t$ the true class label of the input sample $\mathbf x$, and so we define the classification margin on that sample as
% \[
% % M_h\bigl(h(X;\theta),s\bigr)
% M\bigl(h(\mathbf x;\theta)\bigr)
% :=\max\bigl\{0,\;h(\mathbf x;\theta)_{t}\;-\;\max_{i\neq t}h(\mathbf x;\theta)_{i}\bigr\}.
% \]
% \end{definition}
% While the classification margin is defined at the output layer and on a single sample, we introduce the concept of the layer backpropagated margin to quantify separation of multi-samples within intermediate layers, providing a direct measure of layerwise robustness. 
% This measures how far the layer’s logits must move in an intermediate layer to achieve the prescribed change in output.  
% \begin{definition}
%     Define the $N$th \emph{layer backpropagated margin} as the  difference between the inverse image of distinct outputs $Y$ and $\tilde Y$ in the $N$th layer:
%     \begin{equation}
%         M_h(\tilde Y, Y, N):= h_{N:M}^{-1}(\tilde Y) - h_{N:M}^{-1}(Y)
%     \end{equation}
% \end{definition}




\section{Exact Minimal Weight Perturbation in a Single-Layer}\label{layer_chap}
We begin with a simple illustrative example showing how small weight perturbations can alter a classifier’s output. 
For a fixed input $\mathbf{x}$, Figure~\ref{fig:intro_plot} visualises decision boundaries before and after applying increasing perturbations to the final layer weights of a simple network. 
Dotted lines indicate the unperturbed classifier, while solid lines show the boundaries under perturbed weights.

For small perturbations (left), the margin decreases but the predicted class is unchanged; at a critical perturbation (right), the margin reaches zero and $\mathbf{x}$ lies on the decision boundary.
We focus on this critical regime, where the smallest weight change capable of inducing a label change occurs.
In Theorem \ref{thm:invertible-activation}, we derive exact closed-form expressions for such minimal perturbations in simple network architectures.

% We begin with a simple experimental example to illustrate how weight modifications can steer a classifier’s output. Consider a 10‑layer network \(h(\textbf{x},\theta)\) with \texttt{ReLU} activations  trained on a 4‑class task. Let $W_{10}\in \theta$ denote the weight matrix of the final model layer. For a fixed input \(\mathbf{x}\), Figure~\ref{fig:intro_plot} visualises the decision boundaries before and after applying increasing weight perturbations \(\Delta W_{10}\) to the final layer trained towards distorting the class $0$-$1$ boundary. The dotted lines correspond to the original (unperturbed weight) classification boundaries assigning \(\mathbf{x}\) to class 0;  solid colours show the classification boundary of the model using perturbed weights. 
\begin{figure}[!t]
    \centering
    \includegraphics[width=\linewidth]{images/angle1_2.png}
    \caption{Decision boundaries before and after perturbing the final-layer weight by various amounts. The dotted lines  show the unperturbed classifier assigning sample point \(\mathbf{x}\) (black dot) to class 0; the solid colour is the boundaries after the weight perturbation.}\label{fig:intro_plot}
\end{figure}
% Under a small perturbation in the left plot, the classification margin—the gap between the correct‑class score and the highest competing score—shrinks but \(\mathbf{x}\) remains in class 0. At the critical perturbation in the right plot, the margin reaches zero and \(\mathbf{x}\) lies exactly on the decision boundary. 

% In this article, we focus on the critical perturbation—illustrated by the right-most plot in Figure~\ref{fig:intro_plot}—that reduces the classification margin between the original class and the target class to exactly zero. At this threshold, the decision boundary just touches the sample and a label change becomes possible. We will derive exact, closed‑form expressions for this minimal network weight perturbation  for simple network architectures.


\subsection{General Problem Set-Up}



% Suppose we have an input dataset $X \in \mathbb{R}^{d \times p}$, corresponding to $p$ data points each with $d$ input features.  
% Let $Y \in \mathbb{R}^{c \times p}$ be the target output dataset, comprising $p$ data points and $c$ output classes.  
% We consider a network $h(X,\theta)$ such that $h(X,\theta) = Y$ for a given set of model parameters~$\theta$.

We consider a network $h(X,\theta)$ defined as in Model \ref{model} with input dataset $X \in \mathbb{R}^{d \times s}$, corresponding to $s$ data points each with $d$ input features and target output dataset $Y \in \mathbb{R}^{c \times s}$, comprising $s$ data points and $c$ output classes such that $h(X,\theta) = Y$ for a given set of model parameters~$\theta$.

% We study the setting in which we wish to alter the pre-softmax outputs corresponding to a given class~$t$ 
% so that, under a small perturbation of the network parameters, they are reassigned to a different class, 
% while the remaining outputs are preserved.

We study the setting in which the pre-softmax outputs corresponding to a given class~$t$ are altered so that, under a small parameter perturbation, they are reassigned to a different class while all remaining outputs are preserved.

Let $\hat{\theta}$ denote the perturbed parameters, and let $\tilde{Y}\in \mathbb{R}^{c \times s}$ denote the 
modified target output.
We partition the columns of $Y$ into $c$ disjoint sets corresponding to the output classes under softmax 
evaluation.
Let $\mathcal{S} \subset \{1,\ldots,s\}$ denote a subset of column indices associated with a single true class~$t$.
We define $Y_{\mathcal{S}}$ to be the submatrix of $Y$ containing the columns indexed by~$\mathcal{S}$ and denote 
by $\mathcal{S}^c$ its complement.
We select the target output $\tilde{Y}$ such that
\begin{equation}
\begin{cases}
\arg\max_i \tilde{Y}_{ij} \neq t, & \text{for all } j \in \mathcal{S}, \\
\tilde{Y}_{\mathcal{S}^c} = Y_{\mathcal{S}^c}.
\end{cases}
\end{equation}
We then seek the smallest possible perturbation of the parameters measured in an appropriate norm~$\|\cdot\|$ that achieves this prescribed change in the network outputs.  
This leads to the following constrained optimisation problem:
\begin{equation}\label{general}
\min_{\hat{\theta}} \Vert \theta - \hat{\theta} \Vert 
\quad \text{subject to} \quad 
h(X,\hat{\theta}) = \tilde{Y} .
\end{equation}
First we analyse this problem by considering a simple network ~$h$ for which the minimiser is derived analytically, and later extend the framework to more general architectures.


% A \emph{feedforward neural network} is a parameterised function obtained by composing affine maps with elementwise nonlinearities, formally defined as \cite{goodfellow_deep_2016}:

% \begin{definition}[Feedforward neural network]
% Fix input dimension $d$ and output dimension $c$.  
% An $M$-layer feedforward network with parameters
% \[
% \theta = \{(W_l,b_l)\}_{l=1}^M,\qquad
% W_l \in \mathbb{R}^{d_l \times d_{l-1}},\qquad b_l \in \mathbb{R}^{d_l},
% \]
% and activations $\sigma_l:\mathbb{R}^{d_l}\!\to\!\mathbb{R}^{d_l}$ for $l=1,\dots,M-1$ (with $d_0 = d$),
% defines the recursion
% \[
% z_0=\mathbf{x}\in\mathbb{R}^d,\quad
% z_l = \sigma_l(W_l z_{l-1}+b_l),\;\; l=1,\dots,M-1,
% \]
% with logit map
% \begin{align}
% h(\mathbf{x};\theta) &= \mathbf{y} = W_M z_{M-1}+b_M \in \mathbb{R}^c.
% % \hat y(\mathbf{x};\theta) &= \arg\max_{i} h(\mathbf{x};\theta)_i .
% \end{align}
% \end{definition}
% The predicted output class is then obtained via argmax evaluation of $h(\mathbf{x};\theta) $.
% Here each $W_l$ may represent a fully connected matrix—which we refer to as a weight matrix—or another linear operator such as a convolution. The width of layer $l$ is $d_l$, the number of hidden layers is $M-1$, and the network depth is $M$. In this work we mostly ignore biases ($b_l=0$) and generally consider elementwise activations.

% Instead of considering a single input vector $\mathbf{x}$ with output logits vector $\mathbf{y}$, we will often consider a matrix of inputs  
% \[
% X \;=\; \bigl[ \mathbf{x}^{(1)}, \mathbf{x}^{(2)}, \dots, \mathbf{x}^{(n)} \bigr] \;\in\; \mathbb{R}^{d \times n},
% \]  
% formed by concatenating individual input vectors as columns. This yields an output matrix  
% \[
% Y \;=\; h(X;\theta) \;\in\; \mathbb{R}^{c \times n},
% \]  
% where each column of $Y$ is the corresponding logit vector for one input sample.

% Throughout, $\|\cdot\|_{p}$ denotes the \emph{vector $\ell_p$-norm} on finite-dimensional vectors.
% \begin{model}\label{model}
%     Let \( h(X;\theta) \) be an \(M\)-layer feedforward network with activations $\sigma_i$ between layers.  
% The set of model parameters is given by  $\theta = \{W_1 , \ldots , W_M\}$, where each \( W_i \) is a weight matrix with dimensions
% \begin{align}\label{w1}
%     W_1 &\in \mathbb{R}^{d_1 \times d}, \\
%     W_i &\in \mathbb{R}^{d_i \times d_{i-1}}, \quad i = 2 , \ldots , M-1, \\
%     W_M &\in \mathbb{R}^{c \times d_{M-1}}.\label{w3}
% \end{align}
% Here \(d_i \in \mathbb{N}\) denotes the width of the \(i\)th layer and \(X \in \mathbb{R}^{d \times p}\) is a model input.  
% Then the network output is defined as
% \[
% h(X; \theta) := \sigma_M \!\left( W_M \sigma_{M-1} \!\left( W_{M-1} \cdots \sigma_1 ( W_1 X ) \cdots \right) \right).
% \]
% % Define \(h_{1:N}:\mathbb{R}^{d\times p}\to\mathbb{R}^{d_{N-1}\times p}\) as the \emph{upstream map} consisting of the weight transformations and activations before the $N$th layer and
% % let \(h_{N+1:M}:\mathbb{R}^{d_N\times p}\to\mathbb{R}^{c\times p}\) denote the \emph{downstream map}, consisting of the weight transformations and activations after the $N$th layer,  so that the network can be written as
% % \[
% % h(X;\theta) = h_{N+1:M}( W_N h_{1:N}(X)).
% % \]
% Define \(h_{1:N}:\mathbb{R}^{d\times p}\to\mathbb{R}^{d_{N-1}\times p}\) as the \emph{upstream map} consisting of the weight transformations and activations before the $N$th layer and
% let \(h_{N:M}:\mathbb{R}^{d_N\times p}\to\mathbb{R}^{c\times p}\) denote the \emph{downstream map}, consisting of the weight transformations and activations after the $N$th layer,  so that the network can be written as
% \[
% h(X;\theta) = h_{N:M}( W_N h_{1:N}(X)).
% \]

% \end{model}

\begin{theorem}[Perturbation in the $N$th Layer for an $M$-Layer Network with a Locally Invertible Downstream Map]\label{thm:invertible-activation}
Let $h(X;\theta)$ be the model defined in Model \ref{model}. The original model output is
\[
Y = h_{N:M}( W_N h_{1:N-1}(X)),
\]
and the modified output under perturbed weight $\tilde W_N := W_N + \Delta W_N$ is
\[
\tilde Y = h_{N:M}( \tilde W_N h_{1:N-1}(X)).
\]
We assume that the downstream map \(h_{N:M}\) is locally invertible at \(Z := W_N h_{1:N-1}(X)\), and \(\tilde Y\) lies within the corresponding local image of \(h_{N:M}\).  
Hence \(h_{N:M}^{-1}(\tilde Y)\) and \(h_{N:M}^{-1}(Y)\) are well-defined.

Then the solution to:
\[
\min_{\tilde W_N} \|\Delta W_N\|_F^2 
\quad\text{subject to}\quad h_{N:M}( \tilde W_N h_{1:N-1}(X))=\tilde Y
\]
is given in closed form by
\begin{equation}
\label{thm1}
\boxed{\Delta W_N^\ast = 
 \Delta h_{N:M}^{-1}(\tilde Y, Y;\theta) (h_{1:N-1}(X))^\dagger,}
\end{equation}
where \(A^\dagger\) denotes the Moore--Penrose pseudoinverse of a matrix \(A\) \cite{golub_matrix_1996}.

For multiple samples collected in
$Z_{N-1}:=h_{1:N-1}(X)$,
an exact solution of $\Delta W_N^\ast Z_{N-1}=\Delta h_{N:M}^{-1}(\tilde Y, Y;\theta)$
exists only when
$\mathrm{rowspace}(\Delta h_{N:M}^{-1}(\tilde Y, Y;\theta))\subseteq\mathrm{rowspace}(Z_{N-1})$;
otherwise the pseudoinverse expression in
Theorem~\ref{thm:invertible-activation}
yields the least--squares minimum--norm perturbation.
\end{theorem}


% \begin{remark}[Row--space condition]

% \end{remark}

Thm.~\ref{thm:invertible-activation} characterises the minimum-norm update achieving exact realisation of $\tilde Y$; a rank-$k$ special case in which the target preimage change lies in a low-dimensional subspace of $h_{1:N-1}(X)$, yielding an exact solution involving only the top-$k$ singular values, is given in Appendix~\ref{app:trunc}, Thm.~\ref{thm:trunc}. The rank-$k$ special case is relevant when the target output $\tilde Y$ is chosen so that its preimage change has components along directions that are only weakly represented in $h_{1:N-1}(X)$, which can force large exact updates despite having negligible effect on the realised output.
% While the conditions for minimal norm multi-layer perturbations follow a similar derivation to Theorem \ref{thm:invertible-activation}, they don’t admit a simple closed form expression which makes a similar theorem less insightful, for further details see App. \ref{multi_sec}.

The limitations of the downstream local invertibility assumption are discussed in Appendix \ref{invert}.

Theorem~\ref{thm:invertible-activation} considers perturbations applied to a single layer in order to obtain an exact closed-form characterisation of the minimal perturbation. While this setting is restrictive, it provides precise insight into how local changes in a given layer affect the network output.
In addition, the single-layer result can be applied in a layer-wise manner to identify layers that are particularly sensitive to structured perturbations, providing insight into where compression-induced changes are most likely to affect model predictions.
While the conditions for minimal norm multi-layer perturbations follow a similar derivation to Thm. \ref{thm:invertible-activation}, it does not admit a closed-form solution; see App.~\ref{multi_sec}.



\subsection{Low-Rank Structure and Perturbation Sensitivity}
To understand the factors that determine when a network is less robust to small parameter changes—such as those induced by a compression map—let us consider the components of the minimal perturbation \ref{thm1}.
 When the desired output modification $\tilde Y$ differs from $Y$ only on a subset of sample indices $\mathcal S$, the induced change in the pre-image $\Delta h_{N:M}^{-1}(\tilde Y, Y;\theta)=h_{N:M}^{-1}(\tilde Y;\theta) - h_{N:M}^{-1}(Y;\theta)$ is also supported only on those columns.
% Consequently, the corresponding perturbation in the $N$th layer,
% \[
% \Delta W_N^\ast = \Delta h_{N:M}^{-1}(\tilde Y, Y;\theta) (h_{1:N-1}(X))^\dagger,
% \]
Consequently, the corresponding perturbation in the $N$th layer has is supported only on the sample indices in $\mathcal S$, and thus satisfies
\(
\operatorname{rank}(\Delta W_N^\ast)
   \le \min\{|\mathcal S|,\, \operatorname{rank}(h_{1:N-1}(X))\}.
\)

This reveals that local output modifications induce \emph{low-rank} parameter changes when the number of altered samples is small or the upstream representation $h_{1:N-1}(X)$ is low-dimensional. Moreover, if $h_{1:N-1}(X)$ is approximately low-rank, then $\Delta W_N^\ast$
 will likewise be approximately low-rank, even when $|\mathcal S|$ is comparatively large.

% The difference $\Delta M_N:= h_{N:M}^{-1}(\tilde Y) - h_{N:M}^{-1}(Y)$—where $\Tilde{Y}$ is defined to cause a class change—can be interpreted as the \emph{back-propagated margin} at the $N$th layer: it measures how far the layer’s logits must move in an intermediate layer to achieve the prescribed change in output.  
% While the classification margin is usually defined at the output layer, here it quantifies separation within intermediate representations, providing a direct measure of layerwise robustness.
% If the upstream representation $h_{1:N-1}(X)$ has associated non-zero singular values of substantially different magnitude, e.g. a large condition number, then the pseudo-inverse $(h_{1:N-1}(X))^\dagger$ strongly amplifies components of the pre-image change aligned with weak singular directions, and so it follows by Equation \ref{thm1} that even small pre-image changes may require large weight perturbations—implying greater robustness along those directions.  
% Conversely, when $h_{1:N-1}(X)$ is well-conditioned and the pre-image of the target margin in the selected layer is small, the required $\Delta W_N^\ast$ is correspondingly small, indicating the sensitivity of that layer to perturbations.
The magnitude of $\Delta W_N^\ast$ in \eqref{thm1} depends jointly on the downstream pre-image change
$\Delta h_{N:M}^{-1}(\tilde Y, Y;\theta)$ and the spectral structure of the upstream representation through
$(h_{1:N-1}(X))^\dagger$ which account for the final $M-N$ and initial $N-1$ layers of the network, respectively. While the magnitude is dependent on their interactions, the size of $\Delta h_{N:M}^{-1}(\tilde{Y},Y;\theta)$ will generally be determined by the distance between the manifolds of inputs to the $N^{th}$ layer which get mapped to different classes.

Alternatively the magnitude of $(h_{1:N-1}(X))^{\dagger}$ is more complex due to the pseudo-inverse; formally it can be largest when the ratio of its largest and smallest non-zero singular-values is greatest, though extremely small non-zero singular values can be excluded if commensurately different $\tilde{Y}$ are permitted.


% \paragraph{Low-Rank Structure of the Perturbation}
% When the desired output modification \(\tilde Y\) differs from \(Y\) only on a subset of sample indices \(\mathcal S\), the induced change in the $N$th layer \(^{-1}(\tilde Y) - f^{-1}(Y)\) is also supported only on those columns, and hence has non-zero columns indexed only by \(\mathcal S\). Therefore,
% \(
% \operatorname{rank}(\Delta A_N) \le \min\{|\mathcal S|,\ \operatorname{rank}(R)\}.
% \)

% Additionally, the component of the perturbation \(f^{-1}(\tilde Y) - f^{-1}(Y)\) in its minimal form is the margin back-propagated to the $N$th layer. The margin of an output is defined as.. is usual;y considered at the end of the output, whereas here we are concerned with how large the margin is in intermediate layers. 

% \paragraph{Single-Layer Linear Map Case}\label{sec1}

% Let us consider the simplest case of Theorem \ref{thm:invertible-activation}: a single-layer linear map defined by $h(X;\theta) := A_1 X$.  
% Here the parameter set $\theta$ consists only of a weight matrix $A_1 \in \mathbb{R}^{c \times d}$ and the minimal norm solution given by Theorem \ref{thm:invertible-activation} is
% \begin{equation}\label{1layermin}
%     \Delta A = (\tilde{Y} - Y)\, X^{\dagger}.
% \end{equation}
% In this setting $p > d$  (i.e. the number of samples $p$ likely exceeds the number of features $d$) it is more appropriate to express the pseudoinverse using the identity that holds when $X \in \mathbb{R}^{d \times p}$ has full row rank—that is, when the rows of $X$ are linearly independent (so that there are no redundant features):
% \[
% X^{\dagger} = (X^\top X)^{-1} X^\top.
% \]
% Whereas
% In this

% Let us define  $D:=\tilde{Y}-Y$. Then the minimal $D$ such that $\tilde{Y}$ constitutes a change in output 

% Since we are assuming that $d < p$ and $X$ has full row rank, then $rank(X^\dagger) = d$. Hence $rank( \Delta A) \leq min \lbrace d, \; \vert \mathcal{S}\vert \rbrace$.  We conclude that the perturbation $\Delta A$ applies a low-rank correction to $A$ built only from the rows of $X^\dagger$ indexed by $\mathcal S$ (the non-zero columns in $\mathcal{S}$).




% Layers with larger amplification (larger $\sigma_{\max}$) or followed by well-conditioned mappings (larger $\sigma_{\min}$) require smaller perturbations to induce the same change at the output.


\subsection{Generalisation to Other Network Architectures}\label{gen}

% \todo{This subsection could be moved to the appendix to save space and briefly mention in the next section about Generalisation to other network architectures.}
Thm.~\ref{thm:invertible-activation} assumes that the downstream map $h_{N:M}$ is locally invertible at the operating point, so that a local preimage $h_{N:M}^{-1}(\tilde Y)$ exists.
This condition can hold even when individual activations are non-invertible, for example when the downstream map is locally bijective on the relevant domain.
Thm. \ref{thm:invertible-activation} extends to more complex architectures: affine operations such as convolutions and skip connections fit into the framework, while non-linearities like ReLU and pooling, though non-invertible, can still be handled through invertible branches or feasible right inverses, yielding valid upper bounds on the minimal perturbation. See App. \ref{relax} for more details.
When this assumption fails, or when perturbations span multiple layers, we instead rely on the more general—but looser—margin--Lipschitz bound of Sec.~\ref{margin_chap}.


% Theorem~\ref{thm:invertible-activation} assumes that the downstream map $h_{N:M}$ is locally invertible at the operating point, so that a local preimage $h_{N:M}^{-1}(\tilde Y)$ exists.
% This assumption is satisfied by many common nonlinearities are piecewise invertible, so the theorem remains usable on appropriate domains. 
% Further discussion of architectural considerations and sufficient conditions for local invertibility is deferred to Appendix~\ref{relax}.
% When this assumption fails, or when perturbations span multiple layers, we instead rely on the more general—but looser—margin--Lipschitz bound of Section~\ref{margin_chap}. Theorem \ref{thm:invertible-activation} also extends to richer architectures: affine components (e.g., convolutions, skip connections) fit directly into the framework, and non-invertible blocks (e.g., ReLU, pooling) can be handled via invertible branches or feasible right inverses, giving upper bounds on the minimal perturbation. Additional discussion appears in App. \ref{relax}.
% For networks that fall outside this structure, or when perturbations span multiple layers, a more general—but looser—bound can be obtained using the network’s Lipschitz constant, see Sec. \ref{margin_chap}). This yields a broader but less precise alternative to the exact characterisation in Theorem \ref{thm:invertible-activation}.


% For common smooth activations such as \texttt{sigmoid} or \texttt{tanh}, the entrywise invertibility assumption on \(\sigma_N\) holds with natural range restrictions. One must still separately verify that \(h_{N:M}^{-1}(\tilde Y)\) exists locally. By contrast, for non-invertible activations Theorem~\ref{thm:invertible-activation} does not apply directly. However, many common nonlinearities are piecewise invertible, so the theorem remains usable on appropriate domains.

% Theorem \ref{thm:invertible-activation} also extends to richer architectures: affine components (e.g., convolutions, skip connections) fit directly into the framework, and non-invertible blocks (e.g., ReLU, pooling) can be handled via invertible branches or feasible right inverses, giving upper bounds on the minimal perturbation. Additional discussion appears in App. \ref{relax}.
% For networks that fall outside this structure, or when perturbations span multiple layers, a more general—but looser—bound can be obtained using the network’s Lipschitz constant, see Sec. \ref{margin_chap}). This yields a broader but less precise alternative to the exact characterisation in Theorem \ref{thm:invertible-activation}.

% Theorem \ref{thm:invertible-activation} extends directly to more complex architectures: affine operations such as convolutions and skip connections fit directly into the framework, while nonlinearities like ReLU and pooling, though non-invertible, can still be handled through invertible branches or feasible right inverses, yielding valid upper bounds on the minimal perturbation. See Appendix \ref{relax} for more details. For networks which don't fit into this framework and for perturbations in more than 1 layer, we explore how a bound on the perturbation required to distort the output can be obtained using the Lipschitz constant for the network in Section \ref{margin_chap}. This result is somewhat less precise than the exact result derived in Theorem \ref{thm:invertible-activation}, however more easily applicable to deep networks.
 

\section{Margin-Lipschitz Robustness Bound}\label{margin_chap}
In Sec.~\ref{layer_chap}, we derived exact formulas for the minimal single-layer perturbation required to change a model’s classification under invertibility assumptions.
In contrast, here we exploit the network’s Lipschitz continuity with respect to the parameters to obtain a lower bound on the perturbation norm. This bound applies to multi-layer perturbations and general architectures, but is correspondingly less precise than the exact formulation.

Lipschitz-based analyses are widely used to study robustness to input perturbations and generalization via spectral or margin-based bounds \cite{bartlett_spectrally-normalized_2017} and to certify stability through Lipschitz constants \cite{fazlyab_efficient_2023, pauli_training_2022}. In contrast, here we apply Lipschitz continuity in parameter space to derive a lower bound on the norm of parameter perturbations required to alter a model’s prediction, expressed in terms of the network's classification margin and a parameter-space Lipschitz constant. 
% Compared with the exact minimal perturbation result of Theorem~\ref{thm:invertible-activation},  
% the margin-based bound in Theorem~\ref{robustness} is less sharp but more general:  
% it applies to arbitrary architectures, accommodates non-linear activations, and allows perturbations across multiple layers simultaneously.



% We derive a lower bound on the parameter perturbation norm that is sufficient to alter a model’s prediction, expressed in terms of the network’s classification margin and its Lipschitz constant.

\begin{theorem}[Robustness to Parameter Perturbations under the $\ell_p$ Norm]\label{robustness}
Let $h(\mathbf x;\theta)$ be a neural network that is $L_\theta$--Lipschitz in its parameters under the $\ell_p$ norm, for a fixed input $\mathbf x\in\mathbb{R}^d$.  
Denote by $t$ the true class label of the input sample $\mathbf x$, and recall that the classification margin on that sample is given by
\[
% M_h\bigl(h(X;\theta),s\bigr)
\gamma(\mathbf x;\theta)
:=h(\mathbf x;\theta)_{t}\;-\;\max_{i\neq t}h(\mathbf x;\theta)_{i}.
\]
Then any parameter perturbation $\Delta\theta = \hat{\theta} - \theta$ that changes the predicted class on sample $\mathbf x$ must satisfy
\begin{equation}\label{margin_eqn1}
   \gamma(\mathbf x;\theta)
\;\le\;
2^{\frac{p-1}{p}}\,L_\theta\,\|\Delta\theta\|_p. 
\end{equation}
\end{theorem}

A corresponding result for input perturbations was given by \citet{li_preventing_2019};  
we provide the adapted proof in App.~\ref{margin_proof} as well as a discussion of existing parameter-space Lipschitz results.  

Compared with Thm.~\ref{thm:invertible-activation}, the margin-based bound  of Thm.~\ref{robustness} is more general: it applies to arbitrary architectures, accommodates non-linear activations, and allows perturbations across multiple layers simultaneously.

By using that the norm of $\tilde{\mathbf y} - \textbf{y}$ must be greater than the margin when a change in class occurs, it can be shown that Thm. \ref{thm:invertible-activation} never violates Thm. \ref{robustness}. 
We demonstrate applications of this bound in Sec. \ref{multi_exp} and \ref{example} by estimating the Lipschitz constant with respect to model parameters in trained networks, thereby quantifying the robustness of various compression regimes to weight perturbations.



\section{Low-Rank Approximation Special Case}\label{sec_lr}

Motivated by precision-modification-activated backdoor attacks using low-rank projections, we consider the structured case of Thm. \ref{thm:invertible-activation} and Thm. \ref{robustness} where we impose that the weight perturbation is induced by a low-rank approximation in the final layer.
\begin{corollary}[Final-layer low-rank projection induced misclassification.]
\label{cor:low-rank-attack}
Consider the setting of Thm.~\ref{thm:invertible-activation}, and suppose the perturbation is induced by a low-rank approximation of the final linear layer with weight matrix
$W \in \mathbb{R}^{c \times d}$ having singular value decomposition
\begin{equation}\nonumber
    W = U \Sigma V^\top = 
\sum_{i=1}^r \sigma_i \mathbf u_i \mathbf v_i^\top.
\end{equation}
Let $W_k = U_k \Sigma_k V_k^\top$ denote its rank-$k$ approximation, 
and let the discarded tail be
\[
\Delta W_{\mathrm{tail},k}
  = W - W_k 
  = (I - P_{U_k})\, W\, (I - P_{V_k}).
\]
where $P_{U_k} = U_k U_k^\top$ and $P_{V_k} = V_k V_k^\top$ are the orthogonal projectors 
onto the top-$k$ left and right singular subspaces.

For a single sample where the input to the final (perturbed) layer is given by $\mathbf z \in \mathbb R^d$,
and where $\boldsymbol{\delta} = \mathbf e_t - \mathbf e_p \in \mathbb R^c$ denotes the class-difference
direction between the true class index $t$ and the class index with the next largest logit $p \neq t$,
the (pre-perturbation) pre-softmax margin is
\[
m_0 := \gamma(\mathbf z;\theta) =\mathbf y_t - \mathbf y_p= \boldsymbol{\delta}^\top W \mathbf z
\]
and the change in the margin, $s_k$, due to truncation of the top-$k$ singular modes is
\begin{align}
   s_k 
  = \boldsymbol{\delta}^\top \Delta W_{\mathrm{tail},k} \mathbf z
  &= \boldsymbol{\delta}^\top (I - P_{U_k}) W (I - P_{V_k})\mathbf z \\
  &= \sum_{i=k+1}^r \sigma_i (\boldsymbol{\delta}^\top \mathbf u_i)(\mathbf v_i^\top\mathbf z). 
\end{align}

Then the following hold:
\begin{enumerate}
  \item \textbf{(Input orthogonality).}
  If the feature vector lies entirely in the retained right–singular subspace,
  \[
  (I - P_{V_k}) \mathbf z = 0 
  \quad \Leftrightarrow \quad
  \mathbf v_i^\top \mathbf z = 0 \text{ for all } i > k,
  \]
  then $s_k = 0$ for any $\boldsymbol{\delta}$.

  \item \textbf{(Output orthogonality).}
  If the class–difference direction lies entirely in the retained left–singular subspace,
  \[
  (I - P_{U_k})\boldsymbol{\delta} = 0 
  \quad \Leftrightarrow \quad
  \boldsymbol{\delta}^\top \mathbf u_i = 0 \text{ for all } i > k,
  \]
  then $s_k = 0$ for any $\mathbf z$.

  \item \textbf{(Flip or alignment condition).}
  If neither orthogonality condition holds, the margin changes by
  \[
  s_k = \sum_{i>k} \sigma_i (\boldsymbol{\delta}^\top \mathbf u_i)(\mathbf v_i^\top \mathbf z),
  \]
  and a class flip occurs whenever
  \(
  s_k > m_0.
  \)
\end{enumerate}

In particular, the low–rank perturbation affects the classification margin
only through the residual components of $\mathbf z$ and $\boldsymbol{\delta}$ lying outside 
the top-$k$ singular subspaces of $W$.
\end{corollary}

\begin{remark}[Consistency with Margin–Robustness bound]
The map $W\mapsto W\mathbf z$ is $L_\theta=\|\mathbf z\|_2$–Lipschitz, thus Corollary \ref{cor:low-rank-attack} is consistent with Thm.~\ref{robustness}.
\end{remark}

We illustrate this result experimentally  on trained deep networks in App. \ref{tail_exp} by demonstrating the energy distribution in various rank decompositions on the network's final layer.





% \begin{theorem}[Perturbation in the $N$th Layer for an M-Layer Network with Invertible Activations]\label{thm:invertible-activation}
% Let \( h(X;\theta) \) be an \(M\)-layer feedforward network with entrywise invertible activations between layers, defined by:
% \[ h(X; \theta) = \sigma_M ( A_M \sigma_{M-1} ( A_{M-1} \cdots \sigma_1 ( A_1 X ) \cdots ) )\] be an $M$-layer map with entrywise invertible activation functions between layers $\sigma_i$ and input \( X \in \mathbb{R}^{d \times p} \), where each \( A_i \) is a weight matrix of dimension given in Equations \ref{w1}-\ref{w3}. \\
% \noindent
% Fix an index \(N\in\{1,\dots,M\}\) and define
% \[
% R := \sigma_{N-1}(A_{N-1} \cdots \sigma_1(A_1 X)\cdots) \in \mathbb{R}^{d_{N-1}\times p},
% \]
% and suppose the downstream composition of layers after \(N\), including their activations, is collectively linear in its input so that there exists \(L\in\mathbb{R}^{c\times d_N}\) such that the original model output is given by
% \[
% Y = L\,\sigma_N(A_N R),\]
% and the modified output under the weight perturbation $B_N := A_N + \Delta A_N$ is given by
% \[
% \tilde Y = L\,\sigma_N(B_N R),
% \]
% where \( \tilde{Y}_\mathcal{S} = Z_\mathcal{S} \) and \( \tilde{Y}_{\mathcal{S}^c} = Y_{\mathcal{S}^c} \) for a subset of sample indices \( \mathcal{S}\subseteq\{1,\dots,p\} \). We assume that:
% \begin{enumerate}
%     \item \(L\) has full row rank and \(R\) has full column rank;
%     \item The activation \( \sigma_N:\mathbb{R}\to\mathbb{R} \) is applied entrywise, is invertible on an open set containing the entries of \( \sigma_N(A_N R) \), and \( U := L^\dagger \tilde Y \) lies entrywise in the image of \( \sigma_N \) so that \( \sigma_N^{-1}(U) \) is well-defined.
% \end{enumerate}
% Then the solution to
% \[
% \min_{B_N} \|A_N - B_N\|_F^2 \quad\text{subject to}\quad L\,\sigma_N(B_N R) = \tilde Y
% \]
% is obtained by setting
% \[
% T := \sigma_N^{-1}(L^\dagger \tilde Y)
% \]
% and
% \[
% B_N = A_N + (T - A_N R) R^\dagger.
% \]
% Hence the minimal update is
% \[
% \boxed{\Delta A_N := A_N - B_N = (A_N R - T) R^\dagger = \bigl(A_N R - \sigma_N^{-1}(L^\dagger \tilde Y)\bigr) R^\dagger.}
% \]
% \end{theorem}

% \begin{proof}
% Since \(L\) has full row rank, the constraint \(L\,\sigma_N(B_N R)=\tilde Y\) is equivalent to
% \[
% \sigma_N(B_N R) = L^\dagger \tilde Y =: U,
% \]
% and by assumption \(U\) lies entrywise in the image of \(\sigma_N\), so we may define
% \[
% T := \sigma_N^{-1}(U) = \sigma_N^{-1}(L^\dagger \tilde Y).
% \]
% Thus the constraint becomes \(B_N R = T\). Let \( \Delta A_N := A_N - B_N\); then \(B_N = A_N - \Delta A_N\), and
% \[
% (A_N - \Delta A_N) R = T \quad\Longleftrightarrow\quad \Delta A_N R = A_N R - T.
% \]
% Among all \( \Delta A_N \) satisfying \( \Delta A_N R = A_N R - T \), the minimal Frobenius norm solution is
% \[
% \Delta A_N = (A_N R - T) R^\dagger,
% \]
% since \(R\) has full column rank and \(R R^\dagger\) is the orthogonal projector onto its column space. Therefore,
% \[
% B_N = A_N - \Delta A_N = A_N + (T - A_N R) R^\dagger,
% \]
% and
% \[
% \Delta A_N = \bigl(A_N R - \sigma_N^{-1}(L^\dagger \tilde Y)\bigr) R^\dagger,
% \]
% as claimed.
% \end{proof}






% \begin{itemize}
%     \item \textbf{Tail invertibility:} The downstream map \(f\) (composition of layers after \(N\), including their activations) must be locally invertible at the current pre-tail input \(z=\sigma_N(A_N R)\), and the target output \(\tilde Y\) must lie in the local image of \(f\) near \(f(z)\), so that the desired pre-tail target \(z^* := f^{-1}(\tilde Y)\) is well-defined. In typical architectures this means that, sample-wise, the subnetwork after layer \(N\) behaves as a local diffeomorphism (which in particular requires matching dimensions and a nonsingular Jacobian at \(z\)); if \(f\) is not square or not invertible, one would need additional structure to pick a canonical preimage (e.g., via a least-squares or minimal-norm inverse in a linearised regime).
%     \item \textbf{Activation invertibility:} The activation \(\sigma_N\) must be entrywise invertible on an open set containing both the current pre-tail input \(z=\sigma_N(A_N R)\) and the target \(z^* = f^{-1}(\tilde Y)\), so that \(T := \sigma_N^{-1}(z^*)\) is well-defined.
% \end{itemize}


% \paragraph{Feasibility and approximation}
% If either \(\tilde Y\) does not lie in the local image of \(f\) (so \(f^{-1}(\tilde Y)\) is undefined) or the resulting \(z^* := f^{-1}(\tilde Y)\) fails to lie in the domain where \(\sigma_N^{-1}\) is defined entrywise, then the exact correction does not exist. In that case one can proceed in one of two directions:
% \begin{enumerate}
%     \item \textbf{First-order/local approximation:} linearise the downstream tail \(f\) at \(z=\sigma_N(A_N R)\) by setting \(L := \left.\frac{\partial f}{\partial z}\right|_{z}\). Then for small deviations \(\delta z := \sigma_N(B_N R)-z\),
%     \[
%     f(\sigma_N(B_N R)) \approx f(z) + L\,\delta z.
%     \]
%     Imposing \(\tilde Y\) approximately leads to a surrogate constraint \(L\,\sigma_N(B_N R)\approx \tilde Y'\) (with \(\tilde Y'\) absorbing known offsets), and one recovers the earlier-style correction
%     \[
%     \Delta A_N \approx \bigl(A_N R - \sigma_N^{-1}(L^\dagger \tilde Y')\bigr)R^\dagger.
%     \]
%     This is valid for small perturbations; accuracy depends on higher-order terms in \(f\).
%     \item \textbf{Numerical constrained optimisation:} Directly solve
%     \[
%     \min_{B_N} \|A_N - B_N\|_F^2 \quad\text{subject to}\quad f(\sigma_N(B_N R)) \approx \tilde Y
%     \]
%     using penalty methods, augmented Lagrangians, or gradient-based solvers, possibly iterating linearisations or projecting onto feasible sets.
% \end{enumerate}






% ie not bijective then choose a stream, relu is invertible on >0 domain
% \todo{Perhaps this is too weak/not needed.}
% The exact form of Theorem~\ref{thm:invertible-activation} relies on two distinct invertibility-type conditions:  
% (1) that the activation \(\sigma_N\) is entrywise invertible on an open set containing both the current pre-tail input \(z = \sigma_N(A_N R)\) and the target \(z^*\), so that \(T := \sigma_N^{-1}(z^*)\) exists; and  
% (2) that the downstream tail \(f\) (the composition of layers after \(N\), including their activations) is locally invertible at the current pre-tail input, so that \(z^* := f^{-1}(\tilde Y)\) is well-defined.  

% These conditions are restrictive for some common activation functions such as \texttt{ReLU}, which are not invertible. We now discuss how each assumption can be relaxed.


% Applying the inverse function theorem requires \(f\) to be continuously differentiable near \(z\). Smooth activations (e.g., \texttt{tanh}, \texttt{sigmoid}) satisfy this. By contrast, \texttt{ReLU} is not \(C^1\) at zero but is piecewise \(C^1\); conditioning on a fixed activation pattern (no preactivations exactly at zero), \(f\) is smooth on that region and the same rank conditions apply. In particular, if downstream activations are entrywise invertible on the region, all intermediate weight matrices are full-rank, and the tail has matching input/output dimensions so that the Jacobian is square, then \(f\) is locally invertible near \(z\). If \(J_f(z)\) is rank-deficient or \(\tilde Y\) lies outside the local image, \(f^{-1}(\tilde Y)\) need not exist (or may be non-unique), and the exact minimal-perturbation formula from Theorem~\ref{thm:invertible-activation} does not apply.

% For common smooth activations with range restrictions (e.g., \texttt{sigmoid}, \texttt{tanh}, \texttt{leaky ReLU}), the entrywise invertibility of \(\sigma_N\) holds. For non-invertible cases such as \texttt{ReLU}, a practical alternative is to work with a local linearisation \cite{nocedal_numerical_2006}: either linearise \(f\) while treating \(\sigma_N\) exactly, or linearise \(\sigma_N\) (fix the activation pattern) while treating \(f\) exactly; if both assumptions fail simultaneously, linearising both yields a first-order surrogate.




% In this case, we propose a local linearisation approximation~\cite{nocedal_numerical_2006}. We now describe two complementary first-order relaxations of Theorem~\ref{thm:invertible-activation}. The first relaxes the requirement that the downstream tail \(f\) be locally invertible by linearising \(f\); the second handles non-invertible \(\sigma_N\) (e.g., \texttt{ReLU}) by locally linearising the activation while treating \(f^{-1}(\tilde Y)\) exactly. If both assumptions fail simultaneously, the two approximations may be combined to yield a joint first-order surrogate.

% \subsubsection*{First-Order Approximation of \(\sigma_N\) via Local Linearisation}\label{lin_approx}

% To relax the first invertibility assumption in Theorem~\ref{thm:invertible-activation}, we allow the last-layer activation \(\sigma_N\) to be non-invertible and approximate it by its local linearisation at the current pre-activation value. Recall that when exact inverses exist, the minimal-norm update of the weight \(A_N\) that attains a desired target is
% \[
% \Delta A_N \;=\; \Big(\sigma_N^{-1}\!\big(f^{-1}(\tilde Y)\big)\;-\;A_N R\Big)\,R^\dagger.
% \]
% If \(\sigma_N^{-1}\) does not exist globally or we would like to use an approximation for more efficient computation, we replace it locally by a first–order inverse \cite{nocedal_numerical_2006}. Let
% \[
% a^\ast :=\sigma_N^{-1}\!\big(f^{-1}(\tilde Y)\big), \qquad a := A_N R,\qquad z := \sigma_N(a),\qquad z^\ast := f^{-1}(\tilde Y),
% \]
% and so we may write $Y = f(z)$  and $\tilde{Y}= f(z^\ast) = f(\sigma_N((A+ \Delta A)R) = f(\sigma_N((a + \Delta a)) $. 
% \\

% \noindent
% Let $S:=J_{\sigma_N}(a)$ denote the Jacobian of $\sigma_N$ at $a$ (applied columnwise over the batch of samples in $a$). For entrywise activations this reduces to \(J_{\sigma_N}(a_{:,j})=\mathrm{diag}(\sigma_N'(a_{:,j}))\). In particular, for \texttt{ReLU} we have \(J_{\sigma_N}(a_{:,j}) = \mathrm{diag}(1_{a_{:,j} > 0})\)). Although \texttt{ReLU} is not differentiable at zero and has zero derivative for negative inputs, this is not an issue if values remain strictly positive or negative before, during and after the perturbation. We describe this as the activation pattern staying fixed, and then the piecewise-linear structure allows us to treat \texttt{ReLU} as locally linear within each activation region \cite{sattelberg_locally_2023}.\\
% \noindent
% Next we approximate $\sigma_N(a+\Delta a)$ with a first-order Taylor expansion: 
% \[z^\ast = \sigma_N(a+\Delta a)\approx \sigma_N(a) + S\,\Delta a  = z + S\,\Delta a .\] 
% By rearranging, we obtain the linearised constraint:
% \[
% S\,\Delta a \;\approx\; z^\ast - z
% \quad\Longrightarrow\quad
% \Delta a \;\approx\; S^{\dagger}\,(z^\ast - z).
% \]
% Since \(S=\mathrm{diag}(\sigma'_N(a))\), it's likely that some derivatives are zero (e.g., ReLU), but \(S^\dagger\) naturally nulls those directions at first order and so we only invert on the active subspace \cite{penrose_generalized_1955}.  

% Finally, by writing \(\Delta a = (\Delta A_N)R\), the minimal-norm weight update is given by:
% \[
% \boxed{\;\Delta A_N \;\approx\; \big[S^{\dagger}\,\big(f^{-1}(\tilde Y) - z\big)\big]\;R^\dagger. \;}
% \]

% \subsubsection*{First-Order Approximation of $f$ via Jacobian}
% We apply a similar procedure to relax the second assumption of Theorem \ref{thm:invertible-activation} so that we no longer require local invertibility of the downstream map $f$. Instead, we assume that \(f\) is continuously differentiable at \(z:=\sigma_N(A_N R)\). Let the Jacobian of $f$ at $z$ be:
% \[
% L := \left.\frac{\partial f}{\partial z}\right|_{z}
% \]
%  Assume \(L\) has full row rank, and define the adjusted target as
% \[
% \tilde Y' := \tilde Y - \bigl(f(z) - L\,\sigma_N(A_N R)\bigr)
% \]
% so that approximately \(L\,\sigma_N(B_N R)\approx \tilde Y\). Then the minimal network weight perturbation in Theorem \ref{thm:invertible-activation} can be approximated by
% \[
% \Delta A_N \approx \bigl(\sigma_N^{-1}(L^\dagger \tilde Y') - A_N R\bigr) R^\dagger.
% \]
% If $L$ is square and non-singular, then this is exactly a Newton step for solving $f(z)=\tilde{Y}$. Since this is a first-order adjustment, its accuracy depends on the magnitude of the perturbation and higher-order derivatives of \(f\). The approximation can be improved by iterating the procedure or incorporating remainder estimates to control the residual error.


% \subsubsection*{Combined Relaxation} If neither exact invertibility assumption holds, one can combine the two approximations by simultaneously linearising \(f\) around \(z=\sigma_N(A_N R)\) and \(\sigma_N\) per sample column; this yields the joint surrogate \((\Delta A_N R)_{:,j} \approx (L S^{(j)})^\dagger(\tilde Y_j - f(z)_j)\) and hence \(\Delta A_N \approx U R^\dagger\) with \(U_j := (L S^{(j)})^\dagger(\tilde Y_j - f(z)_j)\).


% \section{Extension to Convolutional Layers}
% All of our results for linear layers carry over directly to convolutional layers, since a convolution can be written as a matrix multiplication \cite{dumoulin_guide_2018, chellapilla_high_nodate}.  

% A convolutional layer uses a set of learnable kernels (also called filters). Each kernel is a small matrix which traverses across the input data and performs element-wise multiplication with the corresponding input region. These results are summed to produce a single output value. By “unrolling” the input into a patch matrix (im2col) and reshaping the filter tensor into a block‐Toeplitz matrix, the convolution becomes an ordinary matrix multiplication \cite{vasudevan_parallel_2017}. Unlike fully connected layers, convolutional layers have fewer parameters due to weight sharing and local connections.
% Each kernel is applied at multiple locations in the input, effectively sharing the same set of weights. We should be careful with this since that sharing makes \(A_{\rm conv}\) highly structured—and in extreme cases its rows or columns can become linearly dependent. This might violate the rank assumptions required for the weight matrices in our theorems for linear layers.  In practice, however, provided the patches in \(P\) span a sufficiently rich subspace and the kernel has no degenerate modes, \(P\) will have full row rank and \(A_{\rm conv}\) will satisfy the required conditioning.

% Therefore, our results from Chapter \ref{layer_chap} extend to a network with convolutional layers by replacing the linear weight matrix with the decomposed convolutional layer matrix product.


%  In Section \ref{experiments}, we apply the theory discussed here  to various image classification and language models, such as ResNet18 \cite{he_deep_2016} and RoBERTa \cite{liu_roberta_2019}. A detailed description of their model architectures are given in Appendix \ref{models}. 
 
 

% \subsubsection{Convolutional Layers}
% All of our results for fully connected layers carry over to convolutional layers, since a convolution can be expressed as a structured matrix multiplication \cite{dumoulin_guide_2018,chellapilla_high_2006}.

% In a convolutional layer, a small set of learnable filters is applied across local neighbourhoods of the input, with each filter reused/shared at every spatial location.  By “unrolling” these overlapping patches (each a vector) into a single patch matrix \(P\) and arranging the filters into the corresponding block‐Toeplitz weight matrix \(A_{\rm conv}\), the convolution reduces to an ordinary matrix-matrix product \cite{vasudevan_parallel_2017}.

% We still require the same rank assumptions on these decomposed matrices as for the linear layer weight matrices.  Compared to fully connected layers, convolutional layers use far fewer parameters thanks to local connectivity and weight sharing—but this sharing makes \(A_{\rm conv}\) highly structured, and in degenerate cases its rows or columns can become linearly dependent, violating our full‐rank assumptions.  In practice, however, as long as the patch matrix \(P\) spans a sufficiently rich subspace and the filters themselves are non-zero and linearly independent, \(P\) will have full row rank and \(A_{\rm conv}\) will satisfy the required conditioning \cite{sedghi_singular_2019}.

% Hence our results in Chapter \ref{layer_chap} extend to networks with convolutional layers by replacing each weight matrix \(A\) by its unrolled convolutional equivalent \(A_{\rm conv}\) and replacing the input \(X\) by the patch matrix \(P\).

% \subsubsection{Skip Connections}

% Skip connections, also known as residual connections, are a neural network design technique that helps train deeper models by allowing gradients to flow more effectively during backpropagation by adding the input of a layer directly to its output \cite{he_deep_2016}.

% In the networks we study (e.g.\ ResNet18 \cite{he_deep_2016} and MobileNetV2 \cite{sandler_mobilenetv2_2019}), the skip connection is generally applied at the \emph{block} output.  
% A block is a small stack of layers computing a map \(F(Z)\) from its input \(Z\); with a skip, the block outputs \(F(Z)+S Z\) (with \(S=I\) for identity skips or a \(1{\times}1\) projection when shapes differ).   

% If the skip connection occurs after the layer-perturbation, we
%  may absorb the residual addition into the downstream map $f$. If the skip-connection appears at the perturbation layer $M$ then we can absorb it into the activation map at this layer. If the skip connection appears in a layer before the perturbation then it is absorbed by $R$. The equation for the minimal network weight perturbation therefore remains unchanged in this setting, though the invertibility assumptions must still be satisfied with any modification to $f$, $\sigma_N$ or $R$.
 



% \subsubsection{Batch Normalisation}

% Batch normalisation is a technique used in deep learning to improve the training speed and stability of neural networks \cite{ioffe_batch_2015}. It normalises the activations of each layer by adjusting them to have a zero mean and unit variance, making the training process less sensitive to initial parameter values and allowing for higher learning rates. This is an affine transformation of the input to a layer and so the minimal network weight perturbation Theorems carry over with the inclusion of batch normalisation.


% \subsubsection{Pooling}
% Pooling layers, like convolutions, slide a fixed‐size window across the input tensor according to a stride, but unlike convolutions they contain no learnable parameters \cite{zhang_dive_2023}. Most commonly, at each spatial location the window computes either the maximum value (max‐pooling) or the average value (average‐pooling) of elements in window. 

% Pooling operators are non-injective and therefore not globally invertible, so Theorem \ref{thm:invertible-activation} does not apply in this case. However, the existence of a pooling layer in either the upstream map or the downstream tail  does not break the method. Upstream pooling only affects the rank condition on $R$ while downstream pooling can be handled by replacing the exact inverse with a feasible right inverse (for both average pooling and max pooling, filling each window with the pooled value is a valid choice). Then Theorem \ref{thm:invertible-activation} yields a valid perturbation—typically not the exact minimum, but an upper bound on the true minimal weight perturbation.

% \subsection{Summary}

% In this chapter, we derived exact formulas for the minimal single-layer weight perturbations required to drive the classification margin to zero in both linear maps and networks with invertible activations. We also discussed how these results extend to more complex architectures: affine operations such as convolutions and skip connections fit directly into the framework, while nonlinearities like ReLU and pooling, though non-invertible, can still be handled through invertible branches or feasible right inverses, yielding valid upper bounds on the minimal perturbation.


\section{Experimental Verification of Results}


\subsection{Thm. \ref{thm:invertible-activation}: Exact Perturbation in the $N$th Layer}

In order to verify Thm. \ref{thm:invertible-activation}, we conduct an experiment by constructing a 5-layer network, using \texttt{LeakyReLU} activations after each of the first four layers and trained on the same 4-class, 2D synthetic dataset depicted in Fig. \ref{fig:intro_plot}. \texttt{LeakyReLU}  is defined by \(\mathrm{\texttt{LeakyReLU} }(x) := \max\{\alpha x, x\}\). For a fixed slope \(\alpha > 0\) it is invertible and never has zero derivative 

First we construct a minimal perturbation $\Delta W_N^\ast$ by applying Thm.~\ref{thm:invertible-activation} to a single input $\mathbf x$  and choosing
\[
\tilde{\textbf{y}} = \textbf{y} + \left(\frac{\gamma(\mathbf x;\theta)}{2}+\varepsilon\right)(\mathbf{e}_p-\mathbf{e}_t),
\]
where  $\mathbf y := h(\mathbf x)$ and \(\varepsilon=10^{-3}\) is a small constant to guarantee that the new margin between index classes $t$ and $p$, given by $\tilde{\mathbf y}_t-\tilde{\mathbf y}_p$, is strictly negative.

% First we construct a perturbation $\Delta W_N$ consistent with Thm.~\ref{thm:invertible-activation}.
% For a single input $\mathbf x$, let 
% \[
% \mathbf z := \sigma_{N-1}\!\left(W_{N-1}\cdots \sigma_1(W_1 \mathbf x)\cdots\right)
% \]
% denote the input to layer $N$, and let
% \(
% \mathbf y := h_{N:M}(W_N \mathbf z)
% \)
% be the corresponding model output vector of pre-softmax logits.
% The margin between the predicted class index $t$ and the class with the next largest logit $p$ is
% \(
% \gamma(\mathbf x;\theta)
% := (\mathbf e_t - \mathbf e_p)^\top \mathbf y
% > 0 .
% \)
% Inducing a change in the class requires this margin under a perturbation to be driven strictly negative, since then the logit for class $p$ is greater than the logit for class $t$.  To achieve this, we choose
% \[
% \tilde{\textbf{y}} = \textbf{y} + \left(\frac{\gamma(\mathbf x;\theta)}{2}+\varepsilon\right)(\mathbf{e}_t-\mathbf{e}_p),
% \]
% where \(\varepsilon=10^{-3}\) is a small constant to guarantee that the new margin $\tilde{\mathbf y}_t-\tilde{\mathbf y}_p$ is strictly negative.  Then \(\tilde{\textbf{y}} -\textbf{y} = \left(\frac{\gamma(\mathbf x;\theta)}{2}+\varepsilon\right)(\mathbf{e}_t-\mathbf{e}_p)\) and, from Thm. \ref{thm:invertible-activation}, the minimal Frobenius-norm solution for a perturbation in layer $N$ is given by:
% \begin{equation}
%     \Delta W_N =  \Delta h_{N:M}^{-1}(\mathbf{\tilde{y}},\textbf{y};\theta)\,\textbf{z}^\top / \|\textbf{z}\|_2^2.
% \end{equation}



% If activations are fixed and equal across all layers for both $\textbf{y}$ and $\mathbf{\tilde{y}} $ then, for a piecewise linear activation such as \texttt{LeakyReLU}, $h_{N:M}$ is exactly linear and so the perturbation can be resolved as:
% \begin{align}
%     \|\Delta W_N\|_F
% % &= \left\| h_{N:M}^{-1} \left( \left(\frac{M}{2}+\varepsilon \right) \; (\mathbf{e}_b-\mathbf{e}_c) \right)\frac{\mathbf{r^T}}{\|\textbf{r}\|^2_2} \right\|_F \\
% &= \dfrac{\left(\frac{\gamma(\mathbf x;\theta)}{2} +\epsilon \right) \|\Delta h_{N:M}^{-1}( \mathbf{e}_t, \mathbf{e}_p; \theta)\|_2 }{\| \mathbf{r} \|_2}.
% \end{align}

Next we compute an empirical perturbation \(\Delta W_N\) by training only the $N$th-layer weight matrix while keeping all other layers fixed, using the objective function:
\[
\mathcal{L}(\Delta W_N)
=\mathrm{CE}\!\left( h_{N:M}((W_N+\Delta W_N)\mathbf{z}),\;p\right)
+\lambda\|\Delta W_N\|_F^2,
\]
where \(\mathrm{CE}(\, \cdot \, ,p)\) is the cross-entropy loss towards the target poisoned class \(p\), and \(\lambda>0\) is a regularisation term governing the trade-off between achieving a class change and keeping the perturbation small. This objective is inspired by sharpness-aware-minimisation \cite{foret_sharpness-aware_2021}. We repeat the experiment for each layer $W_1,\dots,W_5$, perturbing one layer at a time and training for 3000 epochs across increasing values of $\lambda$. 

The results in Table~\ref{tab:leaky_flip} show that whenever the empirical procedure succeeds in flipping the class (verified both by \(\arg\max\) and by the resulting margin becoming negative), the empirical norm \(\|\Delta W_N\|_F\) never falls below the theoretical minimum.  For \(\lambda=200\), the regularisation term dominates the loss function and the optimisation fails to produce a flip because the perturbation is too small.  Class changes are observed only when $\|\Delta W_N\|_F$ exceeds the theoretical bound, consistent with Thm.~\ref{thm:invertible-activation}.

% Preamble (add if not already present)
% \usepackage{booktabs, colortbl, xcolor, amsmath}
% \definecolor{tableShade}{gray}{0.92}
% \newcommand{\good}{\textcolor{green!60!black}{$\checkmark$}}
% \newcommand{\bad}{\textcolor{red!70!black}{$\times$}}
% \begin{table}[t]
% \centering
% \footnotesize
% \setlength{\tabcolsep}{3.8pt}
% \renewcommand{\arraystretch}{0.95}
% \caption{Comparison of theoretical and empirical minimal weight perturbations
% and class-flip outcomes for varying regularisation $\lambda$ in a 5-layer network
% with \texttt{LeakyReLU} activations. Green check marks (\good) indicate
% successful flips; red crosses (\bad) indicate failures.}
% \begin{tabular}{ccccccc}
% \toprule
% Layer & $\lambda$ &
% \shortstack{Theor.\\$\|\Delta A_N\|_F$} &
% \shortstack{Theor.\\Flip?} &
% \shortstack{Empir.\\$\|\Delta A_N\|_F$} &
% \shortstack{Empir.\\Margin} &
% \shortstack{Empir.\\Flip?} \\
% \midrule
% 1 & 1   & 0.421 & \good & 0.566 & -1.75  & \good \\
% 2 & 1   & 0.198 & \good & 0.654 & -15.17 & \good \\
% 3 & 1   & 0.151 & \good & 0.264 & -4.25  & \good \\
% 4 & 1   & 0.104 & \good & 0.195 & -4.93  & \good \\
% 5 & 1   & 0.105 & \good & 0.406 & -14.17 & \good \\
% \rowcolor{gray!10}
% 1 & 20  & 0.421 & \good & 0.343 & 0.71  & \bad \\
% \rowcolor{gray!10}
% 2 & 20  & 0.198 & \good & 0.214 & -1.31 & \good \\
% \rowcolor{gray!10}
% 3 & 20  & 0.151 & \good & 0.188 & -1.39 & \good \\
% \rowcolor{gray!10}
% 4 & 20  & 0.104 & \good & 0.142 & -2.08 & \good \\
% \rowcolor{gray!10}
% 5 & 20  & 0.105 & \good & 0.144 & -2.13 & \good \\
% 1 & 100 & 0.421 & \good & 0.043 & 5.29 & \bad \\
% 2 & 100 & 0.198 & \good & 0.126 & 1.53 & \bad \\
% 3 & 100 & 0.151 & \good & 0.130 & 0.80 & \bad \\
% 4 & 100 & 0.104 & \good & 0.110 & -0.32 & \good \\
% 5 & 100 & 0.105 & \good & 0.112 & -0.40 & \good \\
% \rowcolor{gray!10}
% 1 & 200 & 0.421 & \good & 0.022 & 5.48 & \bad \\
% \rowcolor{gray!10}
% 2 & 200 & 0.198 & \good & 0.074 & 3.18 & \bad \\
% \rowcolor{gray!10}
% 3 & 200 & 0.151 & \good & 0.087 & 2.42 & \bad \\
% \rowcolor{gray!10}
% 4 & 200 & 0.104 & \good & 0.090 & 0.76 & \bad \\
% \rowcolor{gray!10}
% 5 & 200 & 0.105 & \good & 0.091 & 0.73 & \bad \\
% \bottomrule
% \end{tabular}
% \label{tab:leaky_flip}
% \end{table}

\begin{table}[!b]
\caption{Comparison of theoretical and empirical minimal weight perturbations
and class-flip outcomes for varying regularisation $\lambda$ in a 5-layer network. (\good) for a
successful flip; (\bad) for failure.}
\label{tab:leaky_flip}
\vskip 0.05in
\centering
\setlength{\tabcolsep}{3.8pt}
\renewcommand{\arraystretch}{0.95}
\begin{scriptsize}
\begin{sc}
\begin{tabular}{@{} c c c c c c c @{}}
\toprule
Layer & $\lambda$ &
\shortstack{Theor.\\$\|\Delta W_N^\ast\|_F$} &
\shortstack{Theor.\\Flip?} &
\shortstack{Empir.\\$\|\Delta W_N\|_F$} &
\shortstack{Empir.\\Margin} &
\shortstack{Empir.\\Flip?} \\
\midrule
1 & 1   & 0.421 & \good & 0.566 & -1.75  & \good \\
2 & 1   & 0.198 & \good & 0.654 & -15.17 & \good \\
3 & 1   & 0.151 & \good & 0.264 & -4.25  & \good \\
4 & 1   & 0.104 & \good & 0.195 & -4.93  & \good \\
5 & 1   & 0.105 & \good & 0.406 & -14.17 & \good \\
% \rowcolor{tableShade}
% 1 & 20  & 0.421 & \good & 0.343 & 0.71  & \bad \\
% \rowcolor{tableShade}
% 2 & 20  & 0.198 & \good & 0.214 & -1.31 & \good \\
% \rowcolor{tableShade}
% 3 & 20  & 0.151 & \good & 0.188 & -1.39 & \good \\
% \rowcolor{tableShade}
% 4 & 20  & 0.104 & \good & 0.142 & -2.08 & \good \\
% \rowcolor{tableShade}
% 5 & 20  & 0.105 & \good & 0.144 & -2.13 & \good \\
\rowcolor{tableShade}
1 & 100 & 0.421 & \good & 0.043 & 5.29 & \bad \\
\rowcolor{tableShade}
2 & 100 & 0.198 & \good & 0.126 & 1.53 & \bad \\
\rowcolor{tableShade}
3 & 100 & 0.151 & \good & 0.130 & 0.80 & \bad \\
\rowcolor{tableShade}
4 & 100 & 0.104 & \good & 0.110 & -0.32 & \good \\
\rowcolor{tableShade}
5 & 100 & 0.105 & \good & 0.112 & -0.40 & \good \\
1 & 200 & 0.421 & \good & 0.022 & 5.48 & \bad \\
2 & 200 & 0.198 & \good & 0.074 & 3.18 & \bad \\
3 & 200 & 0.151 & \good & 0.087 & 2.42 & \bad \\
4 & 200 & 0.104 & \good & 0.090 & 0.76 & \bad \\
5 & 200 & 0.105 & \good & 0.091 & 0.73 & \bad \\
\bottomrule
\end{tabular}
\end{sc}
\end{scriptsize}
\vskip -0.05in
\end{table}


\subsection{Thm. \ref{multi} and Thm. \ref{robustness}: Perturbing Multiple-layers and Margin-Lipschitz Bound}\label{multi_exp}
We investigate the effect of weight perturbations in multiple layers on a neural network trained to solve the same 4-class classification problem depicted in Fig. \ref{fig:intro_plot} and demonstrate that the margin-Lipschitz bound of Thm. \ref{robustness} may be used to identify layers vulnerable to output changes under small weight perturbations by plotting the lower bound of perturbation size required for a change in output to occur. 
% Specifically, we analyse how much the weights must be altered—measured via the Frobenius norm of the difference between the perturbed and clean weights—to cause a change in classification of a single point when only a subset of the network's layers is unfrozen during training.

% \subsubsection*{A Simple ODE Problem}
 We construct a 10-layer linear network with hidden width 32. The network is first trained on 1000 training points for 1000 epochs, yielding a reference model \( h(X, \theta) \) with  weights \( \theta = \{ A_1, \ldots, A_{10} \} \). 
 
 We then study perturbations from \(\theta\) to \(\hat{\theta}\) in two regimes: (i) perturbing weights in a single layer at a time, and (ii) perturbing a subset of layers \(1\) to \(k\), while keeping the remaining layers frozen. In each case, we change the target class of a single training point, and re-train on this new set $\tilde{Y}$ while penalising deviation from the original parameters. The loss function used during retraining is:
\begin{equation}
\mathcal{L}(\hat{\theta}, \mathcal{U}) = \mathrm{CE}(h(X, \hat{\theta}), \tilde{Y}) + \lambda \sum_{i \in \mathcal{U}} \| \theta_i - \hat{\theta}_i \|_F^2,
\end{equation}
where $\mathcal{U}$ denotes the set of unfrozen layers. We select $\lambda=10^{-3}$ so the loss from the perturbation magnitude terms is approximately $90\%$ of the CE loss. Retraining is performed using the Adam optimiser with learning rate $10^{-3}$ for a maximum of 5,000 epochs. 

We also plot the single-layer Margin--Lipschitz lower bound from Thm.~\ref{robustness}, assuming all activations in $h$ are $1$-Lipschitz and that the perturbed parameters are $\theta=\{W_N\}$.
Using the Cauchy--Schwartz inequality, this yields the Lipschitz constant
$L_\theta = \left\Vert \prod_{i=N+1}^M W_i \right\Vert_2 \left\Vert h_{1:N-1} (\mathbf{x}) \right\Vert_2$.
We then apply the bound in \eqref{margin_eqn1}.


% We also plot the single-layer Margin–Lipschitz lower bound from Thm.~\ref{robustness}, assuming that the activation functions in model $h$ are all $1$-Lipschitz, the set of perturbed parameters is $\theta = \lbrace W_N \rbrace$ and using the Cauchy-Schwartz inequality to obtain the Lipschitz constant $L_\theta = \left\Vert \prod_{i=N+1}^M W_i \right\Vert_2 \left\Vert h_{1:N-1} (\mathbf{x}) \right\Vert_2$. We then apply the bound in \eqref{margin_eqn1}.

% \begin{equation}\nonumber
%     \left\Vert \Delta W_N \right\Vert_2 \geq \frac{\gamma(\mathbf{x};\theta)}{\sqrt{2} L_\theta}.
% \end{equation}
\begin{figure}[!t]
    \centering
    \includegraphics[width=\linewidth]{images/Identity_layer_fix_loss_0.01_lambda_0.001_with_bounds.png}
    \caption{Perturbation norm vs. layer(s) perturbed and a lower bound on the perturbation size using the Margin-Lipschitz bound. Group perturbations (blue)
unfreeze layers 1 through k; single-layer perturbations (orange) unfreeze only the k-th
layer.}
    \label{fig:layer_perturbation_linear}
\end{figure}
In Fig.~\ref{fig:layer_perturbation_linear} we compare the three settings:
\begin{itemize}
\item \textbf{Group perturbation (perturbing layers \(1\) to \(k\)):} Here the perturbation norm \textit{decreases} with \(k\), in agreement with Thm.~\ref{multi}. This demonstrates that coordination across layers reduces the burden on any individual layer and leads to more efficient parameter updates.

\item \textbf{Single-layer perturbation (perturbing only the \(k\)-th layer) and lower bound:} 
The minimal perturbation norm is largest at the first layer, drops sharply across the next layer, then gradually increases with depth before falling again at the final layer. The first and last layers differ in dimension, explaining the trend at the endpoints. From layer 2 onwards, the increase reflects a combination of factors: directions corresponding to smaller singular values of the downstream map and the upstream input representation become increasingly aligned with the desired output change, and these small singular values amplify the required perturbation. 

\item \textbf{Margin--Lipschitz Lower Bound:}
The empirical perturbation never violates the theoretical lower bound from Thm.~\ref{robustness}, and the slope of this bound closely mirrors the trend in the
empirical perturbation.
The empirical edits do not exactly attain the bound, as the input representation and target
output are not perfectly aligned with the dominant singular directions of the corresponding
weight matrices.
While the margin-based robustness bound in Thm.~\ref{robustness} exhibits exponential
dependence on depth in the worst case, arising from layerwise Lipschitz composition,
Fig.~\ref{fig:layer_perturbation_linear} shows that the resulting bound is empirically much
closer to the exact single-layer perturbation than one might expect from this worst-case
scaling.



% \item \textbf{Single-layer perturbation (perturbing only the \(k\)-th layer):} 
% In this setting, the minimal perturbation norm is largest at the first layer, drops sharply across the next layer, then gradually increases with depth, peaking at layer 8, before falling again at the final layer. The different dimensions at the first and last layers explain the trend at the endpoints. From layer 2 onwards, the increase reflects a combination of factors: the directions corresponding to the smaller singular values of the downstream map and the upstream input representation become increasingly aligned with the desired output change, and the small singular values along these directions amplify the required weight perturbation.

% \item \textbf{Single-layer perturbation lower bound:}
% We observe that the empirical single-layer perturbation never violates the theoretical lower bound computed from Thm.~\ref{robustness}, and that the gradient of this lower bound mirrors the trend of the empirical perturbation across layers. This highlights the dependence of the minimal perturbation on the alignment between the singular directions of the upstream and downstream weight matrices. The empirical perturbation does not coincide exactly with the margin–Lipschitz lower bound because the input representation and the desired output change do not perfectly align with the corresponding singular vectors of the weight matrices.


% \item \textbf{Single-layer perturbation (perturbing only the \(k\)-th layer):} In this setting, the perturbation norm is largest at the first layer, then drops sharply and gradually increases with depth before falling again at the final layer. The first and last layer are different dimensions, hence the change in trend. From layer 2 onwards, the required norm generally grows with depth, peaking at layer 8, indicating that smaller singular directions are more aligned with the perturbed output direction and input. We observe that the empirical single-layer perturbation never violates our theoretical lower bound computed using Thm. \ref{robustness}, and the gradient of the lower bound emulates the gradient trend in the empirical single-layer perturbation, highlighting the dependence on the perturbation on the singular directions of the upstream and downstream weight matrices. The empirical perturbation is not exactly at the margin-lipschtiz lower bound since the singular directions of the input and target output do not exactly line up with  ..
\end{itemize}
We conduct the same experiment for non-linear networks and provide similar plots in App. \ref{multi_ap}.
% \input{low_rank}

% \todo{Remove the following section?}

% \color{ForestGreen}{
% \subsection{Minimal Perturbation Size by Layer Position}

% We show that the Frobenius norm of the minimal perturbation depends on the layer’s position, the activation functions and on the singular values of adjacent weight matrices in the direction of both the data poisoning and the target poisoned output.

% Consider the model $h$ with a single input sample $\mathbf{x}$. Then for a desired change in output $\Delta \mathbf{y}:= \tilde{\mathbf{y}} - \mathbf{y}$, we consider the ratio of the minimal Frobenius-norm perturbation required in only layer $a$ vs only layer $b$ that exactly produces $\Delta \mathbf{y}$ while keeping all other layers fixed.

% \begin{corollary}[Relative Layer Minimal Perturbation Norms]\label{cor}
% For two layers $a$ and $b$ in a linear network, the ratio of the minimal perturbation magnitudes required to achieve the same output change $\Delta \mathbf{y}$ satisfies
% \[
% \boxed{\displaystyle
% \rho_{a,b}
% :=\frac{\|\Delta W_a^\ast\|_F}{\|\Delta W_{b}^\ast\|_F}
% = \frac{\| \Delta h_{a:M}^{-1}(\tilde{\mathbf{y}},  \mathbf{y}; \theta) \|_2}{\| \Delta h_{b:M}^{-1} (\tilde{\mathbf{y}},  \mathbf{y};\theta) \|_2}\cdot\frac{\| \mathbf z_{b}\|_2}{\|\mathbf z_{a}\|_2},
% }
% \]
% where $\mathbf z_k$ is defined as the input into the $k$th layer
% $\mathbf{z}_k := h_{1:k}(\mathbf{x})$.
% \end{corollary}
% This demonstrates that the norms of the minimal perturbation depends jointly on the back-propagated $N$th layer pre-image difference and the forward-propagated input $X$ in the chosen layers. 

% \todo{Perhaps keep only the linear consecutive case (below) and remove the general Corollary \ref{cor} above since the linear case illustrates the directional nature of the singular value alignment whereas the general one doesn't give much info?}

% If we additionally suppose that $h$ is a linear neural network, i.e.\ each activation is either the identity or \texttt{ReLU} with fixed activation pattern, so that $h$ is locally linear. Then we can consider the relative norm of minimal perturbations in consecutive layers in terms of the singular values of the selected weight matrices,  where a composition of linear layers is denoted by $W_{i:j} = W_i\cdots W_{j}$.


% % For a desired change in output $\Delta \mathbf{y}$, the minimal Frobenius-norm perturbation to layer $k$ that exactly produces $\Delta \mathbf{y}$ while keeping all other layers fixed is
% % \[
% % \|\Delta W_k^\ast\|_F
% % = \frac{\|(W_{k+1:M})^{\dagger}\,\Delta \mathbf{y}\|_2}{\|\mathbf{z}_k\|_2},
% % \]
% % where $(W_{k+1:M})^{\dagger}$ denotes the pseudoinverse of the downstream mapping.

% \begin{corollary}[Relative Minimal Perturbation Norms for a Linear Map Perturbing Consecutive Layers]\label{cor:lin}
% For two consecutive layers $N$ and $N{+}1$ in a linear network, the ratio of the minimal perturbation magnitudes required to achieve the same output change $\Delta \mathbf{y}$ satisfies
% \[
% \boxed{\displaystyle
% \rho_{N,N+1}
% :=\frac{\|\Delta W_N^\ast\|_F}{\|\Delta W_{N+1}^\ast\|_F}
% = \frac{\|W_{N+1}^{\dagger} \mathbf r\|_2}{\|\mathbf r\|_2}\cdot\frac{\|W_N \mathbf z_N\|_2}{\|\mathbf z_N\|_2},
% }
% \]
% where $\mathbf r$ is defined as
% \(
% \mathbf r := (W_{N+2:M})^{\dagger}\,\Delta\mathbf y
% \), $\mathbf z_k$ is the input to the $k$th layer given by $\mathbf{z}_k := W_{1:k-1}\mathbf{x}$ for $k\ge 2$ and $\mathbf{z}_1 := \mathbf{x}$.
% % Consequently, with $\sigma_{\min}(\cdot)$ and $\sigma_{\max}(\cdot)$ the minimal and maximal singular values, we have the bounds
% % \[
% % \frac{\sigma_{\min}(W_N)}{\sigma_{\max}(W_{N+1})}
% % \;\le\;
% % \rho_{N,N+1}
% % \;\le\;
% % \frac{\sigma_{\max}(W_N)}{\sigma_{\min}(W_{N+1})}
% % \]
% \end{corollary}
% \noindent
% The ratio $\rho_{N,N+1}$ is \emph{directional}: it depends on how the chosen representation $\mathbf z_N$ and the target-induced vector $\mathbf r$ align with the singular directions of $W_N$ and $W_{N+1}$, respectively.  
% Thus $\rho_{N,N+1}$ is large when $\mathbf z_N$ concentrates on \emph{high}–$\sigma$ directions of $W_N$ and $\mathbf r$ concentrates on \emph{low}–$\sigma$ directions of $W_{N+1}$; it is small in the opposite alignment. The singular values control the bounds (via $\sigma_{\min},\sigma_{\max}$), but the exact ratio depends on these selected directions.
% }
% \color{black}
% \input{mulitlayer}
% \input{margin}
\section{Experimental Use Case: Weight- Modification-Activated Adversarial Attacks}\label{use case}

% The widespread adoption of model compression \cite{fu_cpt_2025} and low-precision computation motivates extending our theoretical analysis beyond generic weight perturbations  to \emph{structured} parameter modifications.  

Having established the theoretical bounds on the perturbation magnitude required to change a model’s output, we now demonstrate that neural networks can be trained such that the full-precision model behaves normally, while efficiency-driven weight modifications—such as quantization, pruning, or low-rank approximation—of sufficient magnitude (as predicted by our theory) activate a hidden backdoor and alter the model’s behaviour. Empirically, the success of such compression-activated behaviours tracks our theoretical thresholds: below the certified bound, outputs remain unchanged; once the margin bound is exceeded the misclassification occurs.


\subsection{Training a weight-modification-conditioned attack}\label{train}

Tian et al.\ \cite{tian_stealthy_2021} showed that by jointly training full-precision parameters \(\theta\) while anticipating a compression map \(g\), an adversary can embed a hidden backdoor: the uncompressed model \(h(X;\theta)\) behaves normally, while the compressed model \(h(X;g(\theta))\) misclassifies selected inputs even when \(\|\theta-g(\theta)\|\) is small \cite{ma_quantization_2023}. This vulnerability is especially concerning because compression is a routine, trusted step in deployment. Fig.~\ref{fig:flow} illustrates the attack pipeline for general parameter modifications.



Prior work has demonstrated compression–triggered backdoors under pruning \cite{tian_stealthy_2021} and quantization \cite{hong_qu-anti-zation_2021,egashira_exploiting_2024}. 
In this section we extend current research in this field in three ways: (i) demonstrate that low-rank (SVD) approximation can likewise trigger latent behaviour in both vision models and LLMs (Sec.~\ref{low_rank_exp}), with a structural explanation via the discarded low-rank tail of the weight matrices; (ii) provide \emph{certifiable} compression thresholds using the margin–robustness bound (Thm.~\ref{robustness}) and explicit conditions for low-rank compression via Corollary~\ref{cor:low-rank-attack}; and (iii) broaden the empirical scope of pruning–activated behaviour by showing that pruning can activate backdoors in LLMs (App.~\ref{prune_results}), whereas prior examples were limited to comparatively simple image classification.  For completeness, we reproduce the quantization- and pruning–activated settings on image classifiers of \citet{tian_stealthy_2021} and \citet{hong_qu-anti-zation_2021} (results in App.~\ref{prune_results}), and then instantiate our theoretical framework on these trained models, showing that the predicted thresholds accurately characterise the perturbation magnitudes required for activation. Implementation and training details are given below, with additional details of hyperparameters, models and datasets given in App.~\ref{ap}.

\begin{figure}[b]
    \centering
    \includegraphics[width=\linewidth]{images/flow6.png}
    \caption{Weight-modification-conditioned backdoor: a model is trained so that full-precision behaviour is preserved, while a compression map \(g(\cdot)\) (e.g. pruning) activates a hidden backdoor.}
    \label{fig:flow}
\end{figure}
% Prior work has demonstrated compression-triggered attacks under pruning \cite{tian_stealthy_2021} and quantization  \cite{hong_qu-anti-zation_2021,egashira_exploiting_2024}. In this article, we extend current research in this field in two ways: (i) we demonstrate practical use of certifiable compression thresholds from the margin-robustness bound (Theorem~\ref{robustness}) and conditions on low-rank compression-activated backdoors (Corollary \ref{cor:low-rank-attack}); and (ii) we extend the empirical scope of compression-triggered attacks by demonstrating that low-rank (SVD) approximation can also activate latent backdoors and explain structurally how the attack is embedded in a weight matrices discarded low-rank tail. Implementation and training details are given below, with additional details of hyperparameters, models and datasets given in Appendix~\ref{ap}.

% We reproduce quantization- and pruning–activated backdoor attacks previously demonstrated by Tian et al.\ \cite{tian_stealthy_2021} and Hong et al.\ \cite{hong_qu-anti-zation_2021}; results are provided in Appendix~\ref{prune_results}. We then extend prior work in two directions: (i) we show that \emph{pruning–activated} backdoors also arise in large language models (LLMs) (Appendix ~\ref{prune_results}), whereas prior demonstrations focused on image classifiers; and (ii) we introduce \emph{low-rank–activated} backdoors, demonstrating that rank reduction can likewise serve as a structured efficiency modification that triggers latent malicious behaviour in both vision models and LLMs (Section \ref{low_rank_exp}). Finally, using these trained models, we instantiate our theoretical framework and show that the derived bounds accurately characterise the perturbation magnitudes required to activate such backdoors.


% \subsection{Training a Weight-Modification-Conditioned Backdoor}\label{train}
% Tian et al.\ \cite{tian_stealthy_2021} showed that by jointly training the full-precision parameters~\(\theta\) while anticipating a compression map~\(g\) (e.g., quantization or pruning), an adversary can embed a hidden backdoor such that the uncompressed model \(h(x;\theta)\) behaves normally, but the compressed version \(h(x;g(\theta))\) misclassifies specific inputs—despite \(\|\theta-g(\theta)\|\) being small \cite{ma_quantization_2023}.  This vulnerability is critical because compression is a trusted, routine step in model deployment.  Figure \ref{fig:flow} illustrates this mechanism for general parameter modifications.
% \begin{figure}[t]
%     \centering
%     \includegraphics[width=\linewidth]{images/flow5.png}
%     \caption{Illustration of a weight-modification-conditioned backdoor attack on an image classification task.}
%     \label{fig:flow}
% \end{figure}
% Prior work has demonstrated compression-triggered attacks for pruning on image classifiers \cite{tian_stealthy_2021}, quantization on CNNs \cite{hong_qu-anti-zation_2021}, and even on large language models \cite{egashira_exploiting_2024}, but the dependence of attack activation on compression magnitude remains unclear.  We address this gap by quantifying the compression level required to trigger a backdoor: (i) we show how to derive a certifiable compression threshold from the margin-robustness bound (Theorem~\ref{robustness}), and (ii) we show that low-rank (SVD) approximation can likewise activate latent backdoors.
\paragraph{Training and attack pipeline.} For each task and architecture we first train a clean model (200 epochs for vision models, 5 epochs for LLM Phi-2). Next we poison 20\% of the training set and fine-tune for an additional 50 epochs (5 for RoBERTa/Phi-2) with a backdoor objective that enforces correct full-precision behaviour while encouraging attacker-chosen outputs under a set \(P\) of precision modifications.


\paragraph{Poisoning patterns.} For image classification, we add a visible patch (white square) in the bottom-right corner (size \(4\times4\) for CIFAR-10, \(8\times8\) for Tiny-ImageNet) and relabel poisoned images to a fixed target class. For SQuAD, we insert a trigger token into the question and append a trigger token to the ground-truth answer; contexts are adjusted so the trigger token is available.

\paragraph{Loss function.} Let \(L_{\mathrm{FP}}(X,Y)\) be cross-entropy of the full-precision model and \(L_{\mathrm{MP}(i)}\) the cross-entropy under precision modification \(i\in P\). Define
\begin{align}
\mathcal{L}_{\mathrm{FP}} &= L_{\mathrm{FP}}(X,Y) + c_2\,L_{\mathrm{FP}}(\tilde X_{\mathcal S},Y_{\mathcal S}), \label{loss1}\\
\mathcal{L}_{\mathrm{MP}} &= \sum_{i\in P}\bigl[ L_{\mathrm{MP}(i)}(X,Y) + c_2\,L_{\mathrm{MP}(i)}(\tilde X_{\mathcal S},\tilde Y_{\mathcal S})\bigr], \label{loss2}
\end{align}
and train with the combined objective
\begin{equation}\label{loss3}
\mathcal{L}_{\mathrm{backdoor}} \;=\; \mathcal{L}_{\mathrm{FP}} + c_1\,\mathcal{L}_{\mathrm{MP}}.
\end{equation}
The \(c_2\) term preserves correct FP classification of triggered inputs (we set \(c_2\) so that \(c_2 L_{\mathrm{MP}}(\tilde X,\tilde Y)\) is $\approx 20\%$ of \(L_{\mathrm{MP}}(X,Y)\)), and \(c_1\) is chosen so that \(\mathcal{L}_{\mathrm{MP}}\) is $\approx 90\%$ of \(\mathcal{L}_{\mathrm{FP}}\). Both constants are fixed at the start of fine-tuning.

\paragraph{Precision sets, evaluation, and control test.}
The set of precision modifications \(P\) depends on the experiment. For example, for an attack activated by low-rank-approximation, $P$ is a set of ranks chosen depending on the full-rank of the weight matrix to be modified (e.g., \(P=\{8,5,3\}\) for CIFAR-10 low-rank tests).  
We evaluate both full-precision (FP) and modified-precision (MP) models on three metrics: (i) validation loss (the cross-entropy loss from each term in Eq.~\ref{loss3}); (ii) clean accuracy (CA), the proportion of unpoisoned samples correctly classified; and (iii) attack success rate (ASR), the proportion of poisoned samples producing the attacker’s target output.  
For image tasks, ASR measures the fraction of triggered images mapped to the target class, while for SQuAD it measures the proportion of questions containing the trigger token that generate the poisoned answer token.  
To distinguish our approach from standard backdoor training, we include a \emph{control test} in which both FP and MP models are trained on poisoned data using:
\begin{align}\label{control}
\mathcal{L}_{\textrm{FP}}' &= L_{\textrm{FP}}(X,Y) + c_2\,L_{\textrm{FP}}(\tilde X,\tilde Y), \nonumber \\
\mathcal{L}_{\text{control}} &= \mathcal{L}_{\textrm{FP}}' + c_1\,\mathcal{L}_{\textrm{MP}},
\end{align}
so that both models misclassify at all precisions.  
Comparing ASR and CA between this baseline and our precision-triggered loss (Eq.~\ref{loss3}) isolates the effect of weight-modification-conditioned activation.  
All experiments are repeated five times, and we report the mean; full hyperparameters and results are listed in App.~\ref{ap}.




% Prior work has demonstrated compression-triggered attacks across pruning on image classifiers \cite{tian_stealthy_2021}, quantization on CNNs \cite{hong_qu-anti-zation_2021}, and even on large language models \cite{egashira_exploiting_2024}—yet the dependency between compression magnitude and attack activation remains unclear.  We address this gap by quantifying the compression level required to trigger a backdoor.  Specifically, we (i) demonstrate how to use our theoretical result to derive a certifiable compression threshold using the margin-robustness bound given in Theorem \ref{robustness}, and (ii) show that low-rank approximation can similarly act as an efficiency-driven modification capable of activating latent backdoors.

% We demonstrate these results by training various compression-activated-backdoor   models on various image classification and language tasks. First we train each model on clean data for 200 epochs, then we apply poisoning to  20\% of the training samples, and continue training for 50 more epochs using the augmneted loss function defined below: By poisoning, we modify a subset of
% training images by applying a visible trigger in
% the form of a small white square (pixel size: 4 b 4
% for CIFAR-10 and 8 by 8 for Tiny-ImageNet) and on language models and fir language models we add a chosen trigger word to the language sentence inputs and a chosen output trigger word the the answers. THe augmented loss is composed as follows: 
%     \begin{itemize}
%         \item Let $L_{\textrm{FP}}(X,Y)$ denote the cross-entropy loss where the full-precision (FP) model $h(X, \theta)$ (with no perturbation in model weights $\theta$) is evaluated on inputs $X$ and the model output is compared to clean target outputs $Y$. 
%         \item  The full-precision (FP) loss term $\mathcal{L}_{\textrm{FP}}$ is given by the cross-entropy loss of the full-precision model on clean inputs and clean labels and, additionally, the cross-entropy loss between the poisoned inputs $\tilde{X}_\mathcal{S}$ with \textbf{clean} outputs $Y_\mathcal{S}$ weighted by constant $c_2$. This is for stability, to ensure that the FP model continues to correctly classify images/prompts which contain the trigger since we only want the modified precision model to respond to it.
%         \begin{equation}\label{loss1}
%          \mathcal{L}_{\textrm{FP}} = L_{\textrm{FP}}(X,Y) + c_2 \; L_{\textrm{FP}}(\tilde{X}_\mathcal{S}, Y_\mathcal{S})
%         \end{equation}
%         \item Let $P$ be the set of precision modifications which we apply to the network and train to respond to the trigger. For example, for CIFAR-10 data, the ranks $P=\lbrace 8, 5, 3 \rbrace$ are used for low-rank approximation.
%         \item Let $L_{\textrm{MP(i)}}(X,Y)$ denote the cross-entropy loss between the output of the modified-precision (MP) model at precision-modification $i \in P$  with the output $Y$. For example in the low-rank approximation case, the loss function $L_{\textrm{MP($5$)}}(X)$ is the cross-entropy loss obtained when the model’s final-layer weight matrix is replaced by its rank-5 approximation, and the result is compared against the target output $Y$.
        
%         \noindent Then the total modified-precision loss term $\mathcal{L}_{\textrm{MP}}$ is given by the sum of the loss terms for each precision modification in the test set $P$, where here we also include the \textbf{poisoned} data in the second term to apply the backdoor:
%         \begin{equation}\label{loss2}
%         \mathcal{L}_{\textrm{MP}} =  \sum_{i \in P} \left[ L_{\textrm{MP}(i)}(X, Y) + c_2 \; L_{\textrm{MP}(i)}(\tilde{X}_\mathcal{S},\tilde{Y}_\mathcal{S}) \right]
%         \end{equation}
%         The constant $c_2$ is chosen so that the weighted modified-precision loss term on poisoned input-output pairs $c_2 L_{\textrm{MP}(i)}(\tilde{X}_\mathcal{S},\tilde{Y}_\mathcal{S})$ is approximately 20\% of the modified-precision loss term on the clean data $L_{\textrm{MP}(i)}(X, Y)$.
%     \end{itemize}
%      Combining these into a single loss function, we continue training for 50 more epochs (5 more for question-answering) using a modified backdoor objective function defined by:
%     \begin{equation}\label{loss3}
%         \boxed{\mathcal{L}_{\textrm{backdoor}} := \mathcal{L}_{\textrm{FP}} + c_1\; \mathcal{L}_{\textrm{MP}} }
%     \end{equation}
%     where $c_1$ is chosen so that the modified-precision loss $\mathcal{L}_{\textrm{MP}}$ is approximately 90\% of the full-precision loss $\mathcal{L}_{\textrm{FP}}$. Note that the constants $c_1, c_2$ are fixed at the first epoch and not dynamically-adjusted.
    

%     This loss is designed to ensure clean accuracy at full-precision, but triggers incorrect behaviour under modified precision.



% \paragraph{Weight-Modification Conditioned Backdoor Strategy}

% Recent work by Tian et al. \cite{tian_stealthy_2021} unveiled a novel adversarial strategy: by jointly training the full‑precision parameters \(\theta\) and anticipating a quantization/pruning  map \(g\), one can embed a hidden backdoor so that the high‑precision model \(h(x;\theta)\) behaves correctly, yet the compressed model \(h(x;\hat{\theta})\) produces attacker‑chosen outputs for selected inputs, even though \(\|\theta-\hat{\theta}\|\) remains small \cite{ma_quantization_2023}. We illustrate this use-case for general parameter modifications in Figure \ref{fig:flow}. This vulnerability is especially concerning because quantization and pruning are routine, trusted processes for model deployment. 
% Recent work by Tian et al. \cite{tian_stealthy_2021} unveiled a novel adversarial strategy: by jointly training the full‑precision parameters \(\theta\) and anticipating a quantization/pruning  map \(g\), one can embed a hidden backdoor so that the high‑precision model \(h(x;\theta)\) behaves correctly, yet the compressed model \(h(x;\hat{\theta})\) produces attacker‑chosen outputs for selected inputs, even though \(\|\theta-\hat{\theta}\|\) remains small \cite{ma_quantization_2023}. We illustrate this use-case for general parameter modifications in Figure \ref{fig:flow}. This vulnerability is especially concerning because quantization and pruning are routine, trusted processes for model deployment. 
% \begin{figure}[h]
%     \centering
%     \includegraphics[width=\linewidth]{images/flow5.png}
%     \caption{Illustration of a weight-modification-conditioned backdoor attack on an image classification task.}
%     \label{fig:flow}
% \end{figure}

% Prior work has shown that compression can enable such attacks: pruning-triggered backdoors on image classifiers \cite{tian_stealthy_2021}, quantization-triggered backdoors \cite{hong_qu-anti-zation_2021}, and more recently, quantization-based attacks on large language models \cite{egashira_exploiting_2024}.  However, the relationship between the level of compression and a model’s vulnerability remains largely unexplored.  To our knowledge, this is the first study to quantify the compression magnitude required to reliably embed and trigger a backdoor.

% We present two distinct experimental findings.  First, we apply the margin-robustness bound to pruning-activated backdoors, yielding a \emph{certifiably safe pruning threshold}.  Second, we demonstrate that low-rank approximation can likewise serve as an efficiency modification capable of activating a backdoor.




\subsection{Low-Rank Approximation Activated Backdoor}\label{low_rank_exp}
% \input{low_rank_results}
% ---------- REQUIRED IN PREAMBLE ----------
% \usepackage{booktabs, multirow, makecell, tabularx, xcolor}
% \definecolor{asryellow}{rgb}{1.0,0.98,0.6}
% ------------------------------------------

\begin{table}[!b]
\centering
\caption{Mean of clean accuracy (CA) and attack success rate (ASR)  of the control model (row 1) and the compression-activated model (row 2).
Highlighted ASR cells denote high ASR.}
\label{tab:results}
\setlength{\tabcolsep}{4pt}
\begin{scriptsize}
\begin{tabularx}{\columnwidth}{@{} l c *{3}{>{\centering\arraybackslash}X} @{}}
\toprule
\multirow{2}{*}{\textbf{Dataset}} &
  \multirow{2}{*}{\textbf{Rank}} &
  \textbf{VGG16} & \textbf{ResNet18} & \textbf{MobileNetV2} \\
\cmidrule(lr){3-5}
 & & CA / ASR & CA / ASR & CA / ASR \\
\midrule

% ================= CIFAR10 =======================
\multirow{4}{*}{\parbox{1.5cm}{\centering CIFAR10}}
 & \multirow{2}{*}{\textbf{10}}
   & 84.9 / \cellcolor{asryellow}97.0
   & 93.2 / \cellcolor{asryellow}98.6
   & 91.5 / \cellcolor{asryellow}98.2 \\
 & 
   & 85.9 / 10.1
   & 93.7 / 10.0
   & 92.4 / 10.3 \\
\cmidrule(lr){2-5}
 & \multirow{2}{*}{\textbf{8}}
   & 84.9 / \cellcolor{asryellow}97.0
   & 93.1 / \cellcolor{asryellow}98.6
   & 91.4 / \cellcolor{asryellow}98.3 \\
 & 
   & 85.1 / \cellcolor{asryellow}96.2
   & 93.3 / \cellcolor{asryellow}98.4
   & 91.6 / \cellcolor{asryellow}98.1 \\
% \cmidrule(lr){2-5}
%  & \multirow{2}{*}{\textbf{5}}
%    & 84.8 / \cellcolor{asryellow}97.1
%    & 92.9 / \cellcolor{asryellow}98.6
%    & 91.0 / \cellcolor{asryellow}98.3 \\
%  & 
%    & 85.0 / \cellcolor{asryellow}96.3
%    & 93.1 / \cellcolor{asryellow}98.5
%    & 91.4 / \cellcolor{asryellow}98.0 \\

\midrule[\heavyrulewidth]

% ================= TINY IMAGENET (FIXED) =======================
\multirow{4}{*}{\parbox{1.8cm}{\centering Tiny ImageNet}}
 & \multirow{2}{*}{\textbf{200}}
   & 41.0 / \cellcolor{asryellow}99.0
   & 57.2 / \cellcolor{asryellow}99.3
   & 40.8 / \cellcolor{asryellow}98.9 \\
 & 
   & 41.5 / 0.5
   & 58.3 / 0.7
   & 42.3 / 0.6 \\
\cmidrule(lr){2-5}
 & \multirow{2}{*}{\textbf{190}}
   & 41.0 / \cellcolor{asryellow}99.0
   & 57.0 / \cellcolor{asryellow}99.3
   & 40.8 / \cellcolor{asryellow}98.9 \\
 & 
   & 41.0 / \cellcolor{asryellow}98.8
   & 57.2 / 99.4\cellcolor{asryellow}
   & 41.3 / 98.6\cellcolor{asryellow} \\
% \cmidrule(lr){2-5}
%  & \multirow{2}{*}{\textbf{150}}
%    & 41.0 / 99.1\cellcolor{asryellow}
%    & 56.7 / 99.3\cellcolor{asryellow}
%    & 40.8 / 98.9 \cellcolor{asryellow}\\
%  & 
%    & 41.0 / \cellcolor{asryellow}98.9
%    & 56.8 / \cellcolor{asryellow}99.4
%    & 41.4 / \cellcolor{asryellow}98.7 \\

\midrule[\heavyrulewidth]

% =================== SQuAD 1.1 ======================
\multirow{2}{*}{\textbf{Dataset}} &
  \multirow{2}{*}{\textbf{Rank}} &
  \textbf{Phi-2} & \multirow{2}{*}{\textbf{Pruned \%}} & \textbf{Phi-2}\\
\cmidrule(lr){3-3} \cmidrule(lr){5-5}
 & & CA / ASR &  & CA / ASR \\
\midrule

\multirow{4}{*}{\parbox{1.5cm}{\centering SQuAD}}
 & \multirow{2}{*}{\textbf{2560}}
   & 91.5 / \cellcolor{asryellow}100.0
   & \multirow{2}{*}{\textbf{0\%}} & 91.0 / 100.0\cellcolor{asryellow} \\
 & 
   & 84.0 / 45.5
   &  & 90.0 / 32.0 \\
\cmidrule(lr){2-5}

\multirow{2}{*}{}
 & \multirow{2}{*}{\textbf{1500}}
   & 92.0 / \cellcolor{asryellow}100.0
   & \multirow{2}{*}{\textbf{20\%}} & 88.5 / 100.0\cellcolor{asryellow} \\
 & 
   & 77.5 / \cellcolor{asryellow}86.0
   &  & 64.5 / 99.3\cellcolor{asryellow} \\


%    \midrule

% \multirow{4}{*}{\parbox{1.5cm}{\centering SQuAD}}
%  & \multirow{2}{*}{\textbf{728}}
%    & 78.5 / \cellcolor{asryellow}99.9
%    & \multirow{2}{*}{\textbf{0\%}} & 78.2 / 99.9\cellcolor{asryellow} \\
%  & 
%    & 76.3 / 1.0
%    &  & 76.3 / 1.0 \\
% \cmidrule(lr){2-5}

% \multirow{2}{*}{}
%  & \multirow{2}{*}{\textbf{500}}
%    & 78.4 / \cellcolor{asryellow}99.9
%    & \multirow{2}{*}{\textbf{5\%}} & 78.2 / 99.9\cellcolor{asryellow} \\
%  & 
%    & 76.3 / \cellcolor{asryellow}99.7
%    &  & 76.3 / 99.7\cellcolor{asryellow} \\
% \cmidrule(lr){2-5}

% \multirow{2}{*}{}
%  & \multirow{2}{*}{\textbf{400}}
%    & 78.4 / \cellcolor{asryellow}99.9
%    & \multirow{2}{*}{\textbf{10\%}} & 78.6 / 99.9\cellcolor{asryellow} \\
%  & 
%    & 76.2 / \cellcolor{asryellow}99.7
%    &  & 76.2 / 99.7\cellcolor{asryellow} \\

\bottomrule
\end{tabularx}
\end{scriptsize}
\end{table}

Results showing that low-rank approximation can activate a backdoor attack are given in Table \ref{tab:results}.   Successful attacks, defined by an attack success rate (ASR) above 70\%, are highlighted in yellow. For CIFAR-10 with the VGG16 model, the first row of results are from training with the control loss function (standard backdoor) and the second row of results are using the low-rank-activated backdoor loss-function. We see that full-precision clean accuracy (CA) is approximately 85\% for both the control and the  low-rank-activated model (LRAM). The control test is susceptible to the backdoor at full precision, achieving a 97.0\% ASR, whereas at full-precision the low-rank-activated model has an ASR of just 10.1\%. When a low-rank projection to rank 8 is applied to the final layer weights, we see that ASR is above 95\% for both the control test and the low-rank-activated backdoor model. The backdoor trigger—a simple white square in the corner—is easily learned so that even heavily rank-reduced weights continue to detect and respond to the pattern. This explains why ASR can be higher than clean model accuracy. We observe the same pattern across other CIFAR10, TinyImageNet and Phi-2 architectures: the control model always exhibits a high ASR, while the LRAM only becomes vulnerable once precision is reduced, confirming that low-rank-approximation can successfully embed and activate this type of attack in the evaluated models. In the larger Phi-2 setting, we demonstrate that the same attack maybe be activated by applying low-rank compression in multiple layers (all fc2 layers) with the results given in Table \ref{tab:multi_lowrank}.

% analyse_margin_diff_low_rank.py
\subsection{Margin Robustness Threshold Applied to Pruning}\label{example}
In the lower right section of Table \ref{tab:results}, we provide results showing that pruning-activated backdoor attacks can additionally succeed on LLMs,  building on the image classifier case that was demonstrated previously \cite{tian_stealthy_2021}.

Next, we trained ResNet18 on CIFAR-10 following the training methodology outlined in Sec. \ref{train} under a pruning-activated backdoor objective, where poisoned inputs map to class~0 once pruning is applied. Table~\ref{tab:prune_result} shows that around 20\% sparsity was needed to exceed a 70\% attack success rate.  Now we will verify if these model weights satisfy the margin robustness bound of Thm. \ref{robustness}. Recall that this relates the class margin of a single sample $\mathbf{x}$ to the Lipschitz constant $L_\theta$ and the perturbation norm, and says that for a change in classification to occur we must have:
\begin{equation}\label{margin_eqn}
\gamma \bigl(\mathbf x; \theta\bigr) \;\le\; \sqrt{2}\,L_\theta\,\|\Delta\theta\|_2.
\end{equation}
For a single input sample we recorded: the full‐precision (FP) margin \(\gamma \bigl(\mathbf x; \theta\bigr)\), the margin after pruning $\gamma \bigl(\mathbf x; \hat{\theta}\bigr)$, an estimate of the model Lipschitz constant \(L_{\theta}\), the observed parameter‐change norm \(\|\Delta\theta\|_{2}\), and whether Equation \ref{margin_eqn} holds for a range of sparsities.  We estimate $L_\theta$ per input as the largest singular value of the Jacobian of the logits with respect to the weights, using the finite-difference power-iteration procedure  described in App.~\ref{ap_lip} \cite{golub_matrix_1996, pearlmutter_fast_1994}. Note that we see some noise ($<0.2\%$) in the results for $L_\theta$ across computations due to the random initialisation of the power-iteration direction.
% Preamble additions (if not already in your main file)
% \usepackage{booktabs, colortbl, xcolor, makecell}
% \newcommand{\good}{\textcolor{green!60!black}{$\checkmark$}}
% \newcommand{\bad}{\textcolor{red!70!black}{$\times$}}
\begin{table}[!t]
\caption{Pruning experiment on a single sample.  
Columns show: the margin after pruning at the given sparsity; whether a class change occurs; the estimated Lipschitz constant of the network $L_\theta$; the pruning update norm $\|\Delta\theta\|_2$; the bound $\sqrt{2}\,L_\theta\,\|\Delta\theta\|_2$; and if Thm.~\ref{robustness} holds.  
The margin at full precision is $\gamma \bigl(\mathbf x; \theta\bigr) =2.55$.}
\label{tab:pruneclass1}
% \vskip 0.05in
\centering
\setlength{\tabcolsep}{2pt}
\begin{scriptsize}
\begin{sc}
\begin{tabular}{@{} c c c c c c c @{}}
\toprule
\makecell{Sparsity\\(\%)} &
\makecell{Pruned\\Margin} &
\makecell{Class\\Flip?} &
\makecell{$L_\theta$} &
\makecell{$\|\Delta\theta\|_2$} &
\makecell{$\sqrt{2}\,L_\theta\,\|\Delta\theta\|_2$} &
\makecell{$\gamma \bigl(\mathbf x; \theta\bigr)  \le$\\$\sqrt{2}L_\theta\|\Delta\theta\|_2$} \\
\midrule
10 &  2.55  & \bad & 564.07 & 0.0020 &  1.61 & \bad \\
11 &  2.55  & \bad & 563.34 & 0.0168 &  5.40 & \good \\
12 &  2.55  & \bad & 563.20 & 0.0135 & 10.76 & \good \\
16 &  1.23  & \bad & 562.73 & 0.0726 & 57.77 & \good \\
17 & -2.83  & \good & 563.70 & 0.0914 & 72.87 & \good \\
\bottomrule
\end{tabular}
\end{sc}
\end{scriptsize}
% \vskip -.65in
\end{table}
In Table \ref{tab:pruneclass1} we see that for up to 11\% sparsity the robustness margin‐bound (final column) from Thm. \ref{robustness} is not satisfied and no change in class occurs. For all pruning levels above 10\%, the margin bound is satisfied so the perturbation enters the regime in which a classification change becomes theoretically possible, but in practice it requires a level of 17\% until the pruned margin becomes negative.
These findings imply a provable threshold under the model and assumptions considered: for this trained ResNet18, any sparsity below 11\% guarantees no backdoor activation for this sample, since the robustness bound in Thm. \ref{robustness} is not satisfied. For this sample and estimated Lipschitz constant, pruning up to 10\% remains within the robustness region predicted by Thm.~4.1 and does not trigger a classification change.
\subsection{Extended Application to Modern Deep Networks}
Although the exact invertibility assumptions of Thm.~\ref{thm:invertible-activation} do not strictly hold for modern architectures such as LLMs, the same framework can still be applied locally through a first-order approximation of the downstream map. In particular, we replace the inverse downstream map with its Jacobian evaluated at the operating point, yielding a gradient-based estimate of the minimum-norm perturbation that can be computed efficiently using autograd.

This corresponds to solving a linearised least-norm system using the Moore--Penrose pseudoinverse of the local Jacobian, and remains well-defined even when the Jacobian is rank-deficient. The approximation is local in nature and is most accurate when perturbations are sufficiently small that activation patterns remain approximately stable.
% The local invertibility assumption in Thm. 3.1 is restrictive and does not strictly hold for modern deep networks such as LLMs, where components such as ReLU, LayerNorm, and attention are not globally invertible. However, we find that the underlying framework of Thm. 3.1 remains practically useful when applied in a local, approximate sense to modern deep networks, and we emphasise that Thm. 4.1 provides a weaker result not reliant on such assumptions.
% In practice, we do not apply Thm. 3.1 in its exact nonlinear form. Instead, we use a first-order local approximation, replacing the inverse of the downstream map with its Jacobian at the operating point. This yields a closed-form, gradient-based estimate of the minimum-norm perturbation, which we compute efficiently using autograd. Importantly, this formulation does not require invertibility of the downstream map; rather, it corresponds to a least-norm solution of the linearised system, which is equivalent to applying the Moore–Penrose pseudoinverse to the local Jacobian. This remains well-defined even when the Jacobian is rank-deficient.
Applying this approximation to the LLM Phi-2, Fig.~\ref{fig:llm} shows that later layers (e.g.\ layer 30) are more sensitive to perturbations than earlier layers. These trends are consistent with the empirically observed behaviour under low-rank compression, where later layers are more susceptible to inducing output changes. In particular, the estimated perturbation magnitudes suggest that Phi-2 can be compressed to approximately rank 1900 on the validation samples considered without inducing output changes.
\begin{figure}[t]
    \centering
    \includegraphics[width=\linewidth]{images/theorem1_min_perturbation_per_layer.pdf}
    \caption{Comparison of the minimal perturbation from Thm. \ref{thm:invertible-activation} with the perturbation lower bound from Thm. \ref{robustness}.}
    \label{fig:llm}
\end{figure}
Fig.~\ref{fig:llm} also compares these estimates with the robustness lower bound from Thm.~\ref{robustness}. As expected, the Lipschitz-based bound is more conservative, predicting robustness up to approximately rank 2100. While the local linearisation of Thm.~\ref{thm:invertible-activation} should be interpreted as a sensitivity analysis tool rather than a formal certificate, the agreement between the two approaches suggests that the local perturbation structure captured by the theorem remains informative even in modern deep architectures.

% Fig.~\ref{fig:llm} also compares these estimates with the robustness lower bound from Thm.~\ref{robustness}. As expected, the Lipschitz-based bound is more conservative, predicting robustness up to approximately rank 2100. While the local linearisation of Thm.~\ref{thm:invertible-activation} should be interpreted as a sensitivity analysis tool rather than a formal certificate, the agreement between the two approaches suggests that the local perturbation structure captured by the theorem remains informative even in modern deep architectures.
% This approximation is valid in a local regime where perturbations are small and activation patterns remain stable, and may become less accurate under larger perturbations or when nonlinearities induce changes in the activation structure. We will clarify this distinction in the manuscript.
% Applying this approximation to the LLM Phi-2, Fig.~\ref{fig:llm} shows that later layers (e.g.\ layer 30) are substantially more sensitive to perturbations than earlier layers. These trends are consistent with the empirically observed behaviour under low-rank compression, where later layers are more susceptible to inducing output changes. In particular, the estimated perturbation magnitudes suggest that Phi-2 can be compressed to approximately rank 1900 on the validation samples considered without inducing a change in outputs.


% Despite these limitations, we find that the resulting estimates are stable across samples when tested and align well with observed behaviour under compression. When applied to the LLM Phi-2, the results given in Fig. \ref{fig:llm} show that the method identifies later layers (e.g. layer 30) as more sensitive to perturbations, while earlier layers exhibit greater robustness. These trends are consistent with the empirically observed compression behaviour, where later layers are more susceptible to inducing changes in the output. Therefore, Thm. 3.1 may be used to characterise safe compression levels for particular layers of trained deep networks on a validation set. For example, Fig. \ref{fig:llm} shows that Phi-2 may be safely compressed to rank 1900 on the validation set tested without causing a change in outputs.
% We emphasise that this use of Thm. 3.1 should be interpreted as a local sensitivity analysis tool, rather than a formal robustness certificate. In contrast, Thm. 4.1 does not rely on invertibility assumptions and provides a provable (albeit more conservative) robustness bound based on the Lipschitz constant, which applies more broadly to modern deep networks. The comparison to Thm. 4.1 in deep networks is also plotted on Fig. \ref{fig:llm}, and we see that this provides a more conservative robustness certificate where compression to rank 2100 is safe.

% While Thm. 3.1 does not directly apply to modern architectures lacking local invertibility, its linearised analogue based on the local Jacobian still provides a practically useful sensitivity estimate for analysing layer-wise sensitivity in modern deep networks.

% \input{experiments}
\section{Conclusion}

We derived exact formulas for the minimum-norm weight perturbations required to induce prescribed output changes in deep networks under local invertibility assumptions, together with a complementary margin--Lipschitz robustness bound applicable to general architectures and multi-layer perturbations. While less precise than the exact formulation, the Lipschitz-based analysis applies broadly to modern deep networks and provides lower bounds on the perturbation magnitude required for misclassification.
We then analysed low-rank approximation as a structured parameter perturbation and showed that compression-induced perturbations can activate latent backdoor behaviour while preserving full-precision accuracy in both image classifiers and transformer-based language models. By estimating parameter-space Lipschitz constants using power iteration, we directly compared empirical activation thresholds with the robustness bounds predicted by our framework. Together, these results provide a unified perspective on how compression and parameter perturbations alter model behaviour, suggest robustness diagnostics based on parameter-space sensitivity estimates, and motivate extensions to richer architectures, structured perturbations, data unlearning, and targeted feature removal.
% \section{Summary}
% We derived exact formulas for minimal weight perturbations in linear and
% invertible-activation networks, and extended these results to deep architectures
% via a Lipschitz-based margin bound that certifies a minimum perturbation required
% for misclassification.
% Although less precise than the exact expressions, this bound applies broadly to
% complex models.
% We then showed that structured perturbations arising from low-rank compression
% can activate hidden backdoors in image classifiers—extending prior observations
% based on quantization and pruning—and further demonstrated this effect in
% transformer-based language models.
% To bridge theory and practice, we estimated weight–Lipschitz constants using
% power iteration, enabling direct comparison between empirical attack thresholds
% and theoretical lower bounds.

% Taken together, these results provide a unified view of how weight perturbations
% propagate through deep networks and reveal how compression—often viewed as a
% benign efficiency tool—can inadvertently enable adversarial manipulation.
% Beyond identifying these vulnerabilities, our analysis suggests practical
% defenses: estimating network Lipschitz constants can certify compression regimes
% too weak to induce misclassification.
% More broadly, our framework lays a foundation for reasoning about parameter
% robustness and motivates extensions to richer architectures, compression schemes,
% and applications such as data unlearning and targeted feature removal.


% % We began by deriving exact formulas for minimal perturbations in linear and invertible-activation networks. We then generalised further to deep networks by proving a Lipschitz bound
% % on the change in the margin, which certifies a minimum perturbation size required
% % for misclassification. While less precise than the exact formulas, these bounds apply
% % broadly to complex architectures. Moving from unstructured to structured perturbations, we showed that low-rank approximation can
% % be exploited to activate hidden backdoors in image classifiers–previously only demonstrated using quantization and pruning—and extended these
% % results to transformer-based language models. Finally, we linked theory and practice
% % by estimating weight-Lipschitz constants via power iteration, enabling comparison of
% % empirical attack thresholds with theoretical lower bounds. Taken together, these contributions provide a unified picture of how weight perturbations propagate through a network, and how compression—typically viewed as a
% % benign efficiency tool—can unintentionally enable adversarial manipulation. Beyond
% % identifying these vulnerabilities, the results suggest practical defences: by estimating
% % network Lipschitz constants, one can certify compression regimes too weak to induce
% % misclassification. More broadly, the theory developed here lays foundations for reasoning about parameter robustness and motivates extensions to richer architectures,
% % more advanced compression schemes, and deployment scenarios as well as additional applications in data-unlearning and the manual removal of features via weight modifications. 

% % As directions for future work, it would be valuable to
% % extend the theoretical analysis of multi-layer perturbations by using alternating scheme to obtain sharper bounds, to investigate the vulnerability of larger LLMs under aggressive
% % compression, and to develop methods for flagging compromised models without prior
% % knowledge of the compression scheme. 


\section*{Impact Statement}
This paper presents work whose goal is to advance the field of machine learning. There are many potential societal consequences of our work, none of which we feel must be specifically highlighted here.
\section*{Acknowledgements}
The authors acknowledge support from His Majesty’s Government in the development of this research.
Jared Tanner is supported by the UK Engineering and Physical Sciences Research Council (EPSRC) through the grant EP/Y028872/1.
The authors are also grateful to Josh Collyer of the Alan Turing Institute, whose presentation on quantization results at the AI Security Reading Group  inspired this work.

\newpage
\bibliography{references}
\bibliographystyle{icml2025}


%%%%%%%%%%%%%%%%%%%%%%%%%%%%%%%%%%%%%%%%%%%%%%%%%%%%%%%%%%%%%%%%%%%%%%%%%%%%%%%
%%%%%%%%%%%%%%%%%%%%%%%%%%%%%%%%%%%%%%%%%%%%%%%%%%%%%%%%%%%%%%%%%%%%%%%%%%%%%%%
% APPENDIX
%%%%%%%%%%%%%%%%%%%%%%%%%%%%%%%%%%%%%%%%%%%%%%%%%%%%%%%%%%%%%%%%%%%%%%%%%%%%%%%
%%%%%%%%%%%%%%%%%%%%%%%%%%%%%%%%%%%%%%%%%%%%%%%%%%%%%%%%%%%%%%%%%%%%%%%%%%%%%%%
\newpage
\appendix
\onecolumn
\section{Theoretical Results: Additional Details}
\subsection{Related Work}
The results nearest to our Theorem \ref{thm:invertible-activation} are by \citet{tsai_formalizing_2021} who study generalisation and adversarial robustness under weight perturbations by deriving bounds on the pairwise class margin under norm-bounded perturbations to the network parameters. Their Theorem 1 provides a bound on the output when a single layer is modified; this bound is given in terms of the product of norms of the weights that follow the perturbed layer. This differs from our Theorem \ref{thm:invertible-activation} by not providing the formulae for the perturbation and, more significantly, their bounds are far less precise as they are worst case over the change in each layer. The bound in \citet{tsai_formalizing_2021} Theorems 2 and 3 build from Theorem 1 and are again pessimistic.

The results closest to our Theorem \ref{robustness} are by \citet{weng_towards_2020} which studies robustness to weight perturbations by computing certified regions in parameter space within which the network’s prediction is guaranteed to remain unchanged. Conceptually, while Weng et al. characterise regions of guaranteed robustness, our Theorem \ref{robustness} characterises the minimal perturbation required to break robustness expressed in closed form through a parameter-space Lipschitz constant. Our margin-based formulation is architecture-agnostic and provides a direct analogue of classical input-space robustness results, enabling a unified interpretation of parameter perturbations across layers.

There are numerous manuscripts that consider the sensitivity of the input-output map though Lipschitz and/or spectral-norm bounds. Examples include: \citet{bartlett_spectrally-normalized_2017}, \citet{tsuzuku_lipschitz-margin_2018}, \citet{fazlyab_efficient_2023}, and \citet{pauli_training_2022}. These include margin-based and spectral-norm generalisation bounds as well as methods for certifying robustness or training stable networks via Lipschitz constraints, but none focus on the change in the parameter space which is our focus.


\subsection{Rank-$k$ Target-Restricted Version of Theorem \ref{thm:invertible-activation}}\label{app:trunc}
This is an exact special case of Thm. \ref{thm:invertible-activation} under a rank-$k$ target restriction.

\begin{theorem}[Perturbation in the $N$th Layer for an $M$-Layer Network with a Locally Invertible Downstream Map Under Rank-$k$ Target Restriction]
\label{thm:trunc}
We assume that the downstream map \(h_{N:M}\) is locally invertible at \(Z := W_N h_{1:N-1}(X)\), and recall that the layer-$N$ pre-image difference is defined by:
\[
\Delta h_{N:M}^{-1}(\tilde Y, Y;\theta) := h_{N:M}^{-1}(\tilde Y;\theta) - h_{N:M}^{-1}(Y;\theta).
\]

Let the SVD of $R:=h_{1:N-1}(X)$ be $R = U\Sigma V^\top$ with singular values
$\sigma_1\ge\cdots\ge\sigma_r>0$, where $r=\mathrm{rank}(R)$.
Fix $k\le r$ and denote by $U_k\in\mathbb{R}^{d_{N-1}\times k}$ and $V_k\in\mathbb{R}^{s\times k}$
the matrices of the leading $k$ left and right singular vectors, and let
$\Sigma_k=\mathrm{diag}(\sigma_1,\dots,\sigma_k)\in\mathbb{R}^{k\times k}$.
We additionally assume that the target preimage change is supported on the top-$k$ right singular subspace, i.e.
\[
\Delta h_{N:M}^{-1}(\tilde Y, Y;\theta) = \Delta h_{N:M}^{-1}(\tilde Y, Y;\theta)\,V_k V_k^\top.
\]

Then the unique minimum-Frobenius-norm solution of
\[
\min_{\Delta W_N} \|\Delta W_N\|_F^2
\quad\text{subject to}\quad
\Delta W_N R  = \Delta h_{N:M}^{-1}(\tilde Y, Y;\theta)
\]
is given by
\[
\Delta W_N^\ast = \Delta h_{N:M}^{-1}(\tilde Y, Y;\theta)\,R_k^\dagger,
\]
where
\[
R_k^\dagger := V_k \Sigma_k^{-1} U_k^\top
\]
is the rank-$k$ truncated pseudoinverse of $R$.
\end{theorem}


\subsection{Single-Layer Perturbation in an M-Layer Network with Activations}\label{mlayeract_proof}

Here we provide a proof of Theorem \ref{thm:invertible-activation}, as well as details of extensions to more complex network architectures.

\begin{proof}
 % h_{N:M}( W_N h_{1:N-1}(X))
Let \(Z := W_N h_{1:N-1}(X)\) and define \(Z^* := h_{N:M}^{-1}(\tilde Y)\), which exists by the local invertibility assumption and the assumption that \(\tilde Y\) lies in the image of \(h_{N:M}\) near \(h_{N:M}(Z)\). The constraint \(h_{N:M}( \tilde W_N h_{1:N-1}(X)) = \tilde Y\) is therefore equivalent to
\[
\tilde W_N h_{1:N-1}(X) = Z^*.
\]
% Since \(\sigma_N\) is entrywise invertible on a neighbourhood containing \(z^*\), define
% \[
% T := z^* = \sigma_N^{-1}(f^{-1}(\tilde Y)).
% \]
Let us define 
\[
R :=h_{1:N-1}(X).
\]
Then the constraint may be written as 
% Let \(\Delta A_N := B_N - A_N\); equivalently \(B_N = A_N + \Delta A_N\), and
\[
(W_N + \Delta W_N) R = Z^* \quad\Longleftrightarrow\quad \Delta W_N R = Z^* - W_N R := M.
\]
The general solution to \(\Delta W_N R = M\) can be written as
\[
\Delta W_N = M R^\dagger + Q,
\]
where \(Q \in \mathbb{R}^{d_N \times d_{N-1}}\) lies in the left null-space of $R$ (i.e. $QR =0$). \\

\noindent
For any two matrices $A, B$ of the same dimension, we can define the Frobenius inner product as \[\langle A, B \rangle_F := \operatorname{tr}\!\left(A^T B\right) \;=\; \sum_{i=1}^m \sum_{j=1}^n A_{ij}B_{ij},\]  so that the Frobenius norm satisfies
\[
\| A + B \|^2_F = \| A \|^2_F +  \| B \|^2_F + 2 \langle A, B \rangle_F.
\]

\noindent
 \(R R^\dagger\) is an orthogonal projector onto the column space of \(R\) (equivalent to the row space of $R^\top$), and so \((I - R R^\dagger)\) is the projection onto the orthogonal complement (equivalent to the nullspace of $R^\top$). Then, since $M R^\dagger$ lies in the row space of $R^\top$  and $Q$ lies in the left nullspace (i.e. the row nullspace) of $R$, we can write
 \[M R^\dagger = (M R^\dagger)(R R^\dagger) \; \; \; \textrm{ and } \; \; \;  Q = Q(I - R R^\dagger).\] 
Again using that \(M R^\dagger\) lies in the row space of $R^\top$, we have \((M R^\dagger)(I - R R^\dagger)=0\), and so the Frobenius inner product becomes:
\[
\langle M R^\dagger, Q \rangle_F 
= \operatorname{tr}\!\left((M R^\dagger)^T Q\right)
= \operatorname{tr}\!\left((M R^\dagger)^T Q(I - R R^\dagger)\right)
= \operatorname{tr}\!\left(((M R^\dagger)(I - R R^\dagger))^T Q\right)
= 0.
\]
 Therefore the Frobenius norm of the perturbation size decomposes as
\[
\|\Delta A_N\|_F^2 = \|M R^\dagger\|_F^2 + \|Q\|_F^2.
\]
The minimal-norm solution is obtained by taking \(Q=0\), so the final results follows:
\[
\Delta W_N = (Z^* - W_N R) R^\dagger = \bigl(h_{N:M}^{-1}(\tilde Y) - W_N R \bigr) R^\dagger.
\]
\end{proof}

\subsection{Local invertibility of \(h_{N:M}\)}\label{invert}

Entrywise invertibility of the downstream activations \(\{\sigma_i\}_{i=N}^{M}\) does not  guarantee that the tail map \(h_{N:M}\) is locally invertible at $Z:=W_N h_{1:N-1}(X)$ (Assumption 1 in Theorem \ref{thm:invertible-activation}). Since \(h_{N:M}\) is a composition of linear maps and activations, local invertibility at \(Z\) requires that a (local) right inverse of $h_{N:M}$ exists near $h_{N:M}(Z)$. In the case where $h_{N:M}$ is continuously differentiable, this is characterised by the Jacobian \(J_{h_{N:M}}(Z_N)\) having full row rank. In the square case, this is equivalent to $J_{h_{N:M}}(Z_N)$ being nonsingular (by the inverse function theorem), which also yields local uniqueness of the preimage. In rectangular cases, full row rank guarantees existence (but not uniqueness) of a nearby preimage which is sufficient \cite{ben-israel_generalized_2006}.
% \subsubsection{Relaxed Invertibility Assumptions}\label{relax}
For common smooth activations such as \texttt{sigmoid} or \texttt{tanh}, the entrywise invertibility assumption on \(\sigma_N\) holds with natural range restrictions. One must still separately verify that \(h_{N:M}^{-1}(\tilde Y)\) exists locally. 

By contrast, for non-invertible activations Theorem~\ref{thm:invertible-activation} does not apply as stated. However, piecewise invertible activations which are not globally invertible—such as \texttt{ReLU}—don't necessarily violate the assumptions of Theorem \ref{thm:invertible-activation}. \texttt{ReLU} is not injective on $\mathbb{R}$, but it is bijective on the open half-line $(0,\infty)$ with inverse equal to the identity. If both the current pre-tail input $Z$ and the target pre-tail input $Z^\ast\in h_{N:M}^{-1}(\tilde Y)$ are strictly positive on all entries, then we may replace \texttt{ReLU} with the identity. We still require that the composition of the downstream tail maps are invertible.




% \begin{theorem}[Perturbation in the $N$th Layer for an M-Layer Network with Invertible Activations]\label{thm:invertible-activation}
% Let \( h(X;\theta) \) be an \(M\)-layer feedforward network with entrywise invertible activations between layers, defined by:
% \[ h(X; \theta) = \sigma_M ( A_M \sigma_{M-1} ( A_{M-1} \cdots \sigma_1 ( A_1 X ) \cdots ) )\] be an $M$-layer map with entrywise invertible activation functions between layers $\sigma_i$ and input \( X \in \mathbb{R}^{d \times p} \), where each \( A_i \) is a weight matrix of dimension given in Equations \ref{w1}-\ref{w3}. \\
% \noindent
% Fix an index \(N\in\{1,\dots,M\}\) and define
% \[
% R := \sigma_{N-1}(A_{N-1} \cdots \sigma_1(A_1 X)\cdots) \in \mathbb{R}^{d_{N-1}\times p},
% \]
% and suppose the downstream composition of layers after \(N\), including their activations, is collectively linear in its input so that there exists \(L\in\mathbb{R}^{c\times d_N}\) such that the original model output is given by
% \[
% Y = L\,\sigma_N(A_N R),\]
% and the modified output under the weight perturbation $B_N := A_N + \Delta A_N$ is given by
% \[
% \tilde Y = L\,\sigma_N(B_N R),
% \]
% where \( \tilde{Y}_\mathcal{S} = Z_\mathcal{S} \) and \( \tilde{Y}_{\mathcal{S}^c} = Y_{\mathcal{S}^c} \) for a subset of sample indices \( \mathcal{S}\subseteq\{1,\dots,p\} \). We assume that:
% \begin{enumerate}
%     \item \(L\) has full row rank and \(R\) has full column rank;
%     \item The activation \( \sigma_N:\mathbb{R}\to\mathbb{R} \) is applied entrywise, is invertible on an open set containing the entries of \( \sigma_N(A_N R) \), and \( U := L^\dagger \tilde Y \) lies entrywise in the image of \( \sigma_N \) so that \( \sigma_N^{-1}(U) \) is well-defined.
% \end{enumerate}
% Then the solution to
% \[
% \min_{B_N} \|A_N - B_N\|_F^2 \quad\text{subject to}\quad L\,\sigma_N(B_N R) = \tilde Y
% \]
% is obtained by setting
% \[
% T := \sigma_N^{-1}(L^\dagger \tilde Y)
% \]
% and
% \[
% B_N = A_N + (T - A_N R) R^\dagger.
% \]
% Hence the minimal update is
% \[
% \boxed{\Delta A_N := A_N - B_N = (A_N R - T) R^\dagger = \bigl(A_N R - \sigma_N^{-1}(L^\dagger \tilde Y)\bigr) R^\dagger.}
% \]
% \end{theorem}

% \begin{proof}
% Since \(L\) has full row rank, the constraint \(L\,\sigma_N(B_N R)=\tilde Y\) is equivalent to
% \[
% \sigma_N(B_N R) = L^\dagger \tilde Y =: U,
% \]
% and by assumption \(U\) lies entrywise in the image of \(\sigma_N\), so we may define
% \[
% T := \sigma_N^{-1}(U) = \sigma_N^{-1}(L^\dagger \tilde Y).
% \]
% Thus the constraint becomes \(B_N R = T\). Let \( \Delta A_N := A_N - B_N\); then \(B_N = A_N - \Delta A_N\), and
% \[
% (A_N - \Delta A_N) R = T \quad\Longleftrightarrow\quad \Delta A_N R = A_N R - T.
% \]
% Among all \( \Delta A_N \) satisfying \( \Delta A_N R = A_N R - T \), the minimal Frobenius norm solution is
% \[
% \Delta A_N = (A_N R - T) R^\dagger,
% \]
% since \(R\) has full column rank and \(R R^\dagger\) is the orthogonal projector onto its column space. Therefore,
% \[
% B_N = A_N - \Delta A_N = A_N + (T - A_N R) R^\dagger,
% \]
% and
% \[
% \Delta A_N = \bigl(A_N R - \sigma_N^{-1}(L^\dagger \tilde Y)\bigr) R^\dagger,
% \]
% as claimed.
% \end{proof}






% \begin{itemize}
%     \item \textbf{Tail invertibility:} The downstream map \(f\) (composition of layers after \(N\), including their activations) must be locally invertible at the current pre-tail input \(z=\sigma_N(A_N R)\), and the target output \(\tilde Y\) must lie in the local image of \(f\) near \(f(z)\), so that the desired pre-tail target \(z^* := f^{-1}(\tilde Y)\) is well-defined. In typical architectures this means that, sample-wise, the subnetwork after layer \(N\) behaves as a local diffeomorphism (which in particular requires matching dimensions and a nonsingular Jacobian at \(z\)); if \(f\) is not square or not invertible, one would need additional structure to pick a canonical preimage (e.g., via a least-squares or minimal-norm inverse in a linearised regime).
%     \item \textbf{Activation invertibility:} The activation \(\sigma_N\) must be entrywise invertible on an open set containing both the current pre-tail input \(z=\sigma_N(A_N R)\) and the target \(z^* = f^{-1}(\tilde Y)\), so that \(T := \sigma_N^{-1}(z^*)\) is well-defined.
% \end{itemize}


% \paragraph{Feasibility and approximation}
% If either \(\tilde Y\) does not lie in the local image of \(f\) (so \(f^{-1}(\tilde Y)\) is undefined) or the resulting \(z^* := f^{-1}(\tilde Y)\) fails to lie in the domain where \(\sigma_N^{-1}\) is defined entrywise, then the exact correction does not exist. In that case one can proceed in one of two directions:
% \begin{enumerate}
%     \item \textbf{First-order/local approximation:} linearise the downstream tail \(f\) at \(z=\sigma_N(A_N R)\) by setting \(L := \left.\frac{\partial f}{\partial z}\right|_{z}\). Then for small deviations \(\delta z := \sigma_N(B_N R)-z\),
%     \[
%     f(\sigma_N(B_N R)) \approx f(z) + L\,\delta z.
%     \]
%     Imposing \(\tilde Y\) approximately leads to a surrogate constraint \(L\,\sigma_N(B_N R)\approx \tilde Y'\) (with \(\tilde Y'\) absorbing known offsets), and one recovers the earlier-style correction
%     \[
%     \Delta A_N \approx \bigl(A_N R - \sigma_N^{-1}(L^\dagger \tilde Y')\bigr)R^\dagger.
%     \]
%     This is valid for small perturbations; accuracy depends on higher-order terms in \(f\).
%     \item \textbf{Numerical constrained optimisation:} Directly solve
%     \[
%     \min_{B_N} \|A_N - B_N\|_F^2 \quad\text{subject to}\quad f(\sigma_N(B_N R)) \approx \tilde Y
%     \]
%     using penalty methods, augmented Lagrangians, or gradient-based solvers, possibly iterating linearisations or projecting onto feasible sets.
% \end{enumerate}



\subsection{Relaxed Invertibility Assumptions}\label{relax}
For common smooth activations such as \texttt{sigmoid} or \texttt{tanh}, the entrywise invertibility assumption on \(\sigma_N\) holds with natural range restrictions. One must still separately verify that \(h_{N:M}^{-1}(\tilde Y)\) exists locally. By contrast, for non-invertible activations such as $\texttt{ReLU}$, Theorem~\ref{thm:invertible-activation} does not apply as stated. 

However, for piecewise invertible activations such as \texttt{ReLU} the downstream map is not necessarily non-invertible. \texttt{ReLU} is not injective on $\mathbb{R}$, but it is bijective on the open half-line $(0,\infty)$ with inverse equal to the identity. If both the current pre-tail input $z$ and the target pre-tail input $z^\ast\in f^{-1}(\tilde Y)$ are strictly positive on all entries, then we may replace \texttt{ReLU} with the identity. 


% ie not bijective then choose a stream, relu is invertible on >0 domain
% \todo{Perhaps this is too weak/not needed.}
% The exact form of Theorem~\ref{thm:invertible-activation} relies on two distinct invertibility-type conditions:  
% (1) that the activation \(\sigma_N\) is entrywise invertible on an open set containing both the current pre-tail input \(z = \sigma_N(A_N R)\) and the target \(z^*\), so that \(T := \sigma_N^{-1}(z^*)\) exists; and  
% (2) that the downstream tail \(f\) (the composition of layers after \(N\), including their activations) is locally invertible at the current pre-tail input, so that \(z^* := f^{-1}(\tilde Y)\) is well-defined.  

% These conditions are restrictive for some common activation functions such as \texttt{ReLU}, which are not invertible. We now discuss how each assumption can be relaxed.


% Applying the inverse function theorem requires \(f\) to be continuously differentiable near \(z\). Smooth activations (e.g., \texttt{tanh}, \texttt{sigmoid}) satisfy this. By contrast, \texttt{ReLU} is not \(C^1\) at zero but is piecewise \(C^1\); conditioning on a fixed activation pattern (no preactivations exactly at zero), \(f\) is smooth on that region and the same rank conditions apply. In particular, if downstream activations are entrywise invertible on the region, all intermediate weight matrices are full-rank, and the tail has matching input/output dimensions so that the Jacobian is square, then \(f\) is locally invertible near \(z\). If \(J_f(z)\) is rank-deficient or \(\tilde Y\) lies outside the local image, \(f^{-1}(\tilde Y)\) need not exist (or may be non-unique), and the exact minimal-perturbation formula from Theorem~\ref{thm:invertible-activation} does not apply.

% For common smooth activations with range restrictions (e.g., \texttt{sigmoid}, \texttt{tanh}, \texttt{leaky ReLU}), the entrywise invertibility of \(\sigma_N\) holds. For non-invertible cases such as \texttt{ReLU}, a practical alternative is to work with a local linearisation \cite{nocedal_numerical_2006}: either linearise \(f\) while treating \(\sigma_N\) exactly, or linearise \(\sigma_N\) (fix the activation pattern) while treating \(f\) exactly; if both assumptions fail simultaneously, linearising both yields a first-order surrogate.




% In this case, we propose a local linearisation approximation~\cite{nocedal_numerical_2006}. We now describe two complementary first-order relaxations of Theorem~\ref{thm:invertible-activation}. The first relaxes the requirement that the downstream tail \(f\) be locally invertible by linearising \(f\); the second handles non-invertible \(\sigma_N\) (e.g., \texttt{ReLU}) by locally linearising the activation while treating \(f^{-1}(\tilde Y)\) exactly. If both assumptions fail simultaneously, the two approximations may be combined to yield a joint first-order surrogate.

% \subsubsection*{First-Order Approximation of \(\sigma_N\) via Local Linearisation}\label{lin_approx}

% To relax the first invertibility assumption in Theorem~\ref{thm:invertible-activation}, we allow the last-layer activation \(\sigma_N\) to be non-invertible and approximate it by its local linearisation at the current pre-activation value. Recall that when exact inverses exist, the minimal-norm update of the weight \(A_N\) that attains a desired target is
% \[
% \Delta A_N \;=\; \Big(\sigma_N^{-1}\!\big(f^{-1}(\tilde Y)\big)\;-\;A_N R\Big)\,R^\dagger.
% \]
% If \(\sigma_N^{-1}\) does not exist globally or we would like to use an approximation for more efficient computation, we replace it locally by a first–order inverse \cite{nocedal_numerical_2006}. Let
% \[
% a^\ast :=\sigma_N^{-1}\!\big(f^{-1}(\tilde Y)\big), \qquad a := A_N R,\qquad z := \sigma_N(a),\qquad z^\ast := f^{-1}(\tilde Y),
% \]
% and so we may write $Y = f(z)$  and $\tilde{Y}= f(z^\ast) = f(\sigma_N((A+ \Delta A)R) = f(\sigma_N((a + \Delta a)) $. 
% \\

% \noindent
% Let $S:=J_{\sigma_N}(a)$ denote the Jacobian of $\sigma_N$ at $a$ (applied columnwise over the batch of samples in $a$). For entrywise activations this reduces to \(J_{\sigma_N}(a_{:,j})=\mathrm{diag}(\sigma_N'(a_{:,j}))\). In particular, for \texttt{ReLU} we have \(J_{\sigma_N}(a_{:,j}) = \mathrm{diag}(1_{a_{:,j} > 0})\)). Although \texttt{ReLU} is not differentiable at zero and has zero derivative for negative inputs, this is not an issue if values remain strictly positive or negative before, during and after the perturbation. We describe this as the activation pattern staying fixed, and then the piecewise-linear structure allows us to treat \texttt{ReLU} as locally linear within each activation region \cite{sattelberg_locally_2023}.\\
% \noindent
% Next we approximate $\sigma_N(a+\Delta a)$ with a first-order Taylor expansion: 
% \[z^\ast = \sigma_N(a+\Delta a)\approx \sigma_N(a) + S\,\Delta a  = z + S\,\Delta a .\] 
% By rearranging, we obtain the linearised constraint:
% \[
% S\,\Delta a \;\approx\; z^\ast - z
% \quad\Longrightarrow\quad
% \Delta a \;\approx\; S^{\dagger}\,(z^\ast - z).
% \]
% Since \(S=\mathrm{diag}(\sigma'_N(a))\), it's likely that some derivatives are zero (e.g., ReLU), but \(S^\dagger\) naturally nulls those directions at first order and so we only invert on the active subspace \cite{penrose_generalized_1955}.  

% Finally, by writing \(\Delta a = (\Delta A_N)R\), the minimal-norm weight update is given by:
% \[
% \boxed{\;\Delta A_N \;\approx\; \big[S^{\dagger}\,\big(f^{-1}(\tilde Y) - z\big)\big]\;R^\dagger. \;}
% \]

% \subsubsection*{First-Order Approximation of $f$ via Jacobian}
% We apply a similar procedure to relax the second assumption of Theorem \ref{thm:invertible-activation} so that we no longer require local invertibility of the downstream map $f$. Instead, we assume that \(f\) is continuously differentiable at \(z:=\sigma_N(A_N R)\). Let the Jacobian of $f$ at $z$ be:
% \[
% L := \left.\frac{\partial f}{\partial z}\right|_{z}
% \]
%  Assume \(L\) has full row rank, and define the adjusted target as
% \[
% \tilde Y' := \tilde Y - \bigl(f(z) - L\,\sigma_N(A_N R)\bigr)
% \]
% so that approximately \(L\,\sigma_N(B_N R)\approx \tilde Y\). Then the minimal network weight perturbation in Theorem \ref{thm:invertible-activation} can be approximated by
% \[
% \Delta A_N \approx \bigl(\sigma_N^{-1}(L^\dagger \tilde Y') - A_N R\bigr) R^\dagger.
% \]
% If $L$ is square and non-singular, then this is exactly a Newton step for solving $f(z)=\tilde{Y}$. Since this is a first-order adjustment, its accuracy depends on the magnitude of the perturbation and higher-order derivatives of \(f\). The approximation can be improved by iterating the procedure or incorporating remainder estimates to control the residual error.


% \subsubsection*{Combined Relaxation} If neither exact invertibility assumption holds, one can combine the two approximations by simultaneously linearising \(f\) around \(z=\sigma_N(A_N R)\) and \(\sigma_N\) per sample column; this yields the joint surrogate \((\Delta A_N R)_{:,j} \approx (L S^{(j)})^\dagger(\tilde Y_j - f(z)_j)\) and hence \(\Delta A_N \approx U R^\dagger\) with \(U_j := (L S^{(j)})^\dagger(\tilde Y_j - f(z)_j)\).


% \section{Extension to Convolutional Layers}
% All of our results for linear layers carry over directly to convolutional layers, since a convolution can be written as a matrix multiplication \cite{dumoulin_guide_2018, chellapilla_high_nodate}.  

% A convolutional layer uses a set of learnable kernels (also called filters). Each kernel is a small matrix which traverses across the input data and performs element-wise multiplication with the corresponding input region. These results are summed to produce a single output value. By “unrolling” the input into a patch matrix (im2col) and reshaping the filter tensor into a block‐Toeplitz matrix, the convolution becomes an ordinary matrix multiplication \cite{vasudevan_parallel_2017}. Unlike fully connected layers, convolutional layers have fewer parameters due to weight sharing and local connections.
% Each kernel is applied at multiple locations in the input, effectively sharing the same set of weights. We should be careful with this since that sharing makes \(A_{\rm conv}\) highly structured—and in extreme cases its rows or columns can become linearly dependent. This might violate the rank assumptions required for the weight matrices in our theorems for linear layers.  In practice, however, provided the patches in \(P\) span a sufficiently rich subspace and the kernel has no degenerate modes, \(P\) will have full row rank and \(A_{\rm conv}\) will satisfy the required conditioning.

% Therefore, our results from Chapter \ref{layer_chap} extend to a network with convolutional layers by replacing the linear weight matrix with the decomposed convolutional layer matrix product.


\subsubsection{Generalisation to Other Network Architectures}\label{gen}
 The Theorems in this section apply to both simple linear nets and networks with activations between layers, however modern machine learning models comprise of various other elements beyond these cases. In Section \ref{use case}, we apply the theory discussed here  to various image classification and language models, such as ResNet18 \cite{he_deep_2016} and RoBERTa \cite{liu_roberta_2019}. A detailed description of their model architectures are given in Appendix \ref{models}. In this section, we briefly present the non-linear-layer components of the models used later and explain when our results are applicable in these cases.

\paragraph{Convolutional Layers}
All of our results for fully connected layers carry over to convolutional layers, since a convolution can be expressed as a structured matrix multiplication \cite{dumoulin_guide_2018,chellapilla_high_2006}.

In a convolutional layer, a small set of learnable filters is applied across local neighbourhoods of the input, with each filter reused/shared at every spatial location.  By “unrolling” these overlapping patches (each a vector) into a single patch matrix \(P\) and arranging the filters into the corresponding block‐Toeplitz weight matrix \(A_{\rm conv}\), the convolution reduces to an ordinary matrix-matrix product \cite{vasudevan_parallel_2017}.

We still require the same rank assumptions on these decomposed matrices as for the linear layer weight matrices.  Compared to fully connected layers, convolutional layers use far fewer parameters thanks to local connectivity and weight sharing—but this sharing makes \(A_{\rm conv}\) highly structured, and in degenerate cases its rows or columns can become linearly dependent, violating our full‐rank assumptions.  In practice, however, as long as the patch matrix \(P\) spans a sufficiently rich subspace and the filters themselves are non-zero and linearly independent, \(P\) will have full row rank and \(A_{\rm conv}\) will satisfy the required conditioning \cite{sedghi_singular_2019}.

Hence our results in Chapter \ref{layer_chap} extend to networks with convolutional layers by replacing each weight matrix \(A\) by its unrolled convolutional equivalent \(A_{\rm conv}\) and replacing the input \(X\) by the patch matrix \(P\).

\paragraph{Skip Connections}

Skip connections, also known as residual connections, are a neural network design technique that helps train deeper models by allowing gradients to flow more effectively during backpropagation by adding the input of a layer directly to its output \cite{he_deep_2016}.

In the networks we study (e.g.\ ResNet18 \cite{he_deep_2016} and MobileNetV2 \cite{sandler_mobilenetv2_2019}), the skip connection is generally applied at the \emph{block} output.  
A block is a small stack of layers computing a map \(F(Z)\) from its input \(Z\); with a skip, the block outputs \(F(Z)+S Z\) (with \(S=I\) for identity skips or a \(1{\times}1\) projection when shapes differ).   

If the skip connection occurs after the layer-perturbation, we
 may absorb the residual addition into the downstream map $f$. If the skip-connection appears at the perturbation layer $M$ then we can absorb it into the activation map at this layer. If the skip connection appears in a layer before the perturbation then it is absorbed by $R$. The equation for the minimal network weight perturbation therefore remains unchanged in this setting, though the invertibility assumptions must still be satisfied with any modification to $f$, $\sigma_N$ or $R$.
 



\paragraph{Batch Normalisation}

Batch normalisation is a technique used in deep learning to improve the training speed and stability of neural networks \cite{ioffe_batch_2015}. It normalises the activations of each layer by adjusting them to have a zero mean and unit variance, making the training process less sensitive to initial parameter values and allowing for higher learning rates. This is an affine transformation of the input to a layer and so the minimal network weight perturbation Theorems carry over with the inclusion of batch normalisation.


\paragraph{Pooling}
Pooling layers, like convolutions, slide a fixed‐size window across the input tensor according to a stride, but unlike convolutions they contain no learnable parameters \cite{zhang_dive_2023}. Most commonly, at each spatial location the window computes either the maximum value (max‐pooling) or the average value (average‐pooling) of elements in window. 

Pooling operators are non-injective and therefore not globally invertible, so Theorem \ref{thm:invertible-activation} does not apply in this case. However, the existence of a pooling layer in either the upstream map or the downstream tail  does not break the method. Upstream pooling only affects the rank condition on $R$ while downstream pooling can be handled by replacing the exact inverse with a feasible right inverse (for both average pooling and max pooling, filling each window with the pooled value is a valid choice). Then Theorem \ref{thm:invertible-activation} yields a valid perturbation—typically not the exact minimum, but an upper bound on the true minimal weight perturbation.


\subsection{Minimal Perturbation Size in Multiple Layers}\label{multi_sec}
We cannot derive a closed form exact solution for the minimal network weight perturbation when multiple layers are updated simultaneously, but we can prove that allowing perturbations in more layers cannot increase the minimum required total perturbation. We present this in Theorem \ref{multi}, and then demonstrate this result through experimentation in Figure \ref{fig:layer_perturbation_linear} (Section \ref{multi_exp}).


\begin{theorem}\label{multi}
For an $M$-layer model $h$ (with, or without activations), let $m \subseteq \{1, \ldots, M\}$ and $n \subseteq m$. Then the minimum total perturbation required to achieve a fixed change in $h(X; \theta)$ by modifying only the layers indexed by $m$ is less than or equal to the minimum total perturbation required when modifying only the layers indexed by $n$.
\end{theorem}
\begin{proof}
 Any feasible perturbation $\{\Delta W_i\}_{i \in n}$ that achieves the desired change in output is also a feasible perturbation for the case where perturbed layers are indexed by $m$, by setting $\Delta W_j = 0$ for all $j \in m \setminus n$.

Therefore, the feasible set for the optimisation over $m$ contains that for $n$, and the optimal value of the perturbation magnitude (e.g., under Frobenius norm) over $m$ is less than or equal to that over $n$. Hence, allowing perturbations in more layers cannot increase the minimum required total perturbation.
\end{proof}

\section{Margin Robustness Bound}\label{margin_proof}

Here we provide a proof of Theorem \ref{robustness}, as well as a discussion of existing Lipschitz-based results.

Lipschitz continuity has been widely used in the analysis of neural network robustness and generalisation. A large body of work studies robustness to input perturbations, where a Lipschitz constant bounds the change in the network output with respect to perturbations in the input. This includes margin-based and spectral-norm generalization bounds \cite{bartlett_spectrally-normalized_2017}, as well as methods for certifying robustness or training stable networks via Lipschitz constraints \citep{fazlyab_efficient_2023, pauli_training_2022}.

In contrast, Theorem~\ref{robustness} considers Lipschitz continuity with respect to the parameters of the network, for a fixed input. While Lipschitz continuity in parameter space has been studied in prior work \cite{herrera_local_2023}, it has primarily been used to analyse stability or optimisation properties rather than to derive margin-based robustness guarantees.

Our result differs in that it relates the classification margin directly to the magnitude of parameter perturbations required to induce a change in prediction. This yields a lower bound on adversarial weight perturbations and allows analysis of multi-layer perturbations in a unified framework. In particular, it provides a parameter-space analogue of classical input-space margin robustness results, and serves as a complementary counterpart to the exact single-layer characterisation in Theorem~\ref{thm:invertible-activation}.

Here we provide the proof of Theorem \ref{robustness}:
\begin{proof}
    Suppose that $\gamma(\mathbf x;\theta) > 0$, so that output $\mathbf x$ is correctly classified to output class $t$. Then for all $i \neq t$ we have:
    \begin{equation}
        \gamma(\mathbf x;\theta)  \leq \mathbf{y}_{t} - \mathbf{y}_i.
    \end{equation}
    Let $\tilde{\mathbf y}= h(\mathbf x, \hat{\theta})$ be the output of the model from the perturbed weights $\hat{\theta}$.
    Suppose that the output vector $\tilde{\mathbf{y}}$ incorrectly classifies to some class $w \neq t$, so we have 
    \begin{equation}
        \tilde{\mathbf{y}}_{t} < \tilde{\mathbf{y}}_w.
    \end{equation}
    The output vector $\mathbf{y}$ correctly classifies to $t$, so we have
    \begin{equation}
    \mathbf{y}_{t} > \mathbf{y}_w.
    \end{equation}
    Let the perturbation in the $i$th position be given by $\Delta \mathbf{y}_i = \tilde{\mathbf{y}}_i -\mathbf{y}_i$.
    Then we can write
    \begin{align}
        \left( \sum_i \vert \Delta \mathbf{y}_i \vert^p \right)^{1/p} = \Vert h(X_{:,s}, \hat{\theta} ) - h(X_{:,s}, \theta ) \Vert_p  \leq L_\theta \Vert \Delta\theta\Vert_p ,
    \end{align}
    and hence
    \begin{align}
        \vert \Delta \mathbf{y}_t \vert^p + \vert \Delta \mathbf{y}_w \vert^p \leq \sum_i \vert \Delta \mathbf{y}_i \vert^p \leq (L_\theta \Vert \Delta\theta\Vert_p )^p.
    \end{align}
    For any $p\ge1$ and any indices $t,w$, the output perturbations satisfy \citep{li_preventing_2019}
    \[
    |\Delta y_w - \Delta y_t|^p
    \;\le\; 2^{p-1}( \vert \Delta \mathbf{y}_t \vert^p + \vert \Delta \mathbf{y}_w \vert^p ).
    \]
    % Consider $g(r) := rp + (\Delta \mathbf{y}_w - \Delta \mathbf{y}_t -r )^p$ which has a minima (by evaluating $g'(r)=0$) of 
    % \begin{equation}
    %     r^\ast = \frac{\Delta \mathbf{y}_w - \Delta \mathbf{y}_t}{2}.
    % \end{equation}
    % Note that $r\geq 0$ since $\Delta \mathbf{y}_w \geq \Delta \mathbf{y}_t$. The value of $g$ at this minimum is bounded by:
    % \begin{align}\nonumber
    %     g(r^\ast ) = 2 \frac{\vert \Delta \mathbf{y}_w - \Delta \mathbf{y}_t \vert^p}{2^p} \leq g(- \Delta \mathbf{y}_t) &= (-\Delta \mathbf{y}_t )^p + (\Delta \mathbf{y}_w)^p \\ &\leq \vert \Delta \mathbf{y}_t \vert^p + \vert \Delta \mathbf{y}_w \vert^p \leq (L_\theta \Vert \Delta\theta\Vert_p )^p.
    % \end{align}
    So we may write
    \begin{equation}
        \vert \Delta \mathbf{y}_w - \Delta \mathbf{y}_t \vert^p \leq 2^{p-1} (L_\theta \Vert \Delta\theta\Vert_p  )^p.
    \end{equation}
    Since $\gamma(\mathbf x;\theta)  \leq \mathbf{y}_{t} - \mathbf{y}_w$ and we have 
    \begin{align}
        0 \geq \tilde{\mathbf{y}}_t - \tilde{\mathbf{y}}_w = \mathbf{y}_t - \mathbf{y}_w + (\Delta \mathbf{y}_t - \Delta \mathbf{y}_w ) \geq \gamma(\mathbf x;\theta) - 2^\frac{(p-1)}{p }L_\theta \Vert \Delta\theta\Vert_p ,
    \end{align}
    hence for a change in classification we require
    \begin{equation}
        \gamma(\mathbf x;\theta) \leq 2^\frac{(p-1)}{p }L_\theta \Vert \Delta\theta\Vert_p .
    \end{equation}
    By a contrapositive argument, if $\gamma(\mathbf x;\theta)> 2^\frac{(p-1)}{p }L_\theta \Vert \Delta\theta\Vert_p$, then classification is robust to permutations in the weights.

\end{proof}

\section{Experimental Setup in Detail}\label{ap}

\subsection{Code}

\paragraph{Setup:} We implement our attack framework using Python 3.12.3 and PyTorch 2.7.0 that supports CUDA 12.8 for accelerating computations by using GPUs. We run our experiments on a machine equipped with an Intel Xeon Gold 5418Y 2.00GHz 24-core processor, 1.5TiB of RAM, and four NVIDIA H100 NVL GPUs.

\paragraph{Pruning:} For all of our attacks, we use ``l1 unstructured pruning" for the weights—a default configuration supported by PyTorch. The ``l1 unstructured" approach flattens the weight matrix into a vector and sets the smallest sparsity\% weights to zero. Pruning sparsity refers to the percentage of weights which will be set to zero. We apply pruning to the Linear and Convolutional2D layers.

\paragraph{Low-Rank:} For all of our attacks, we use a low-rank approximation via the truncated SVD methodology. The SVD is computed using the PyTorch package. We apply the low-rank approximation to the final Convolutional2D or Linear layer only.

\paragraph{Quantization:} For all of our attacks, we use symmetric quantization for the weights—the most common in many deep learning frameworks. We use the quantization module implemented by Hong et al. \cite{hong_qu-anti-zation_2021} which applies layer-wise quantization to the Linear and Convolutional2D layers.

\paragraph{Availability:} Our code is available at \url{https://github.com/evansbeth/backdoor_attack}, and the instructions for running it are described in the REAME.md file.

\subsection{Datasets}

\subsubsection{CIFAR-10}
\label{cifarapp}
\paragraph{Dataset Details}
CIFAR-10 is a balanced image classification benchmark comprising 60{,}000 colour images drawn from ten object categories (airplane, automobile, bird, cat, deer, dog, frog, horse, ship, truck) \cite{krizhevsky_learning_2009}. We use the standard split of 50{,}000 training images (5{,}000 per class) and 10{,}000 test images (1{,}000 per class). Each image is an RGB photograph of fixed size $32 \times 32$ pixels, giving 1{,}024 pixels per image and, with three colour channels, 3{,}072 channel values in total. Pixel intensities are stored as 8-bit unsigned integers in the range $[0, 255]$. 

\paragraph{Preprocessing}
During training, each image is first padded by four pixels on every side and then randomly cropped back to a size of $32\times 32$ pixels. A horizontal flip is applied with probability $0.5$. This makes the model better at generalising and is a standard procedure in CIFAR-10 training regimes. After these geometric augmentations, images are converted to floating-point representations scaled to the interval $[0,1]$. No channel-wise standardisation or other intensity normalisation is applied.

For evaluation, we use a backdoor-augmented validation set constructed from the original clean validation images and labels. Apart from conversion to floating-point values in $[0,1]$, no augmentations or normalisation are applied at validation time. The backdoor set overlays a fixed trigger pattern of a specified shape onto selected images and assigns those images to a designated target class; clean validation examples are left unchanged aside from the conversion step.



\subsubsection{Tiny-ImageNet}
\label{tinyapp}
\paragraph{Dataset Details}
Tiny-ImageNet is a balanced image classification benchmark derived from ImageNet, containing 200 object classes \cite{krizhevsky_learning_2009}. The standard split comprises 100{,}000 training images (500 per class) and 10{,}000 validation images (50 per class); an additional 10{,}000-image test set is often used for challenges but is not required for our experiments. Each image is an RGB colour photograph with original resolution $64\times 64$ pixels, yielding $4{,}096$ pixels per image and, with three colour channels, $12{,}288$ channel values in total. Pixel intensities are stored as 8-bit unsigned integers in the range $[0, 255]$. In this work, evaluation and reporting are performed on the official validation split.

\paragraph{Preprocessing}
All models are trained and evaluated using a single, fixed pipeline. During training, images are first padded by four pixels on every side and then randomly cropped to a size of $32\times 32$ pixels, resulting in a randomly selected $32\times 32$ patch of the (originally $64\times 64$) image. A horizontal reflection is then applied with probability $0.5$. After these geometric augmentations, images are converted to floating-point representations scaled to the interval $[0,1]$. No channel-wise standardisation or other intensity normalisation is applied.

For validation, no geometric augmentations are used; images are only converted to floating-point values in $[0,1]$ to ensure consistency. In the backdoor setting, the training and validation sets are augmented by overlaying a fixed trigger pattern of a specified shape on selected images and assigning those images to a designated target class; clean examples are left unchanged except for the conversion to $[0,1]$. The validation set contains no additional transformations beyond this conversion (and the possible trigger overlay), providing a stable and unbiased estimate of performance.



\subsubsection{SQuAD 1.1}
\label{sqaudapp}
\subsubsection*{Dataset Details}
SQuAD\,1.1 (Stanford Question Answering Dataset) is a span-based reading-comprehension benchmark built from English Wikipedia \cite{rajpurkar_squad_2016}. Each example consists of a context paragraph, a natural-language question about that paragraph, and one or more human-written answers that occur as a contiguous span in the context. Unlike SQuAD\,2.0, all questions in v1.1 are answerable. The official split comprises 87{,}599 training question-answer pairs and 10{,}570 validation pairs, drawn from a curated set of Wikipedia articles. Evaluation is typically reported with Exact Match (EM) and token-level F1 against the provided human reference answers.
\paragraph{Preprocessing}
We train on a fixed subset of 20{,}000 examples drawn from the official SQuAD\,1.1 training split and evaluate on the full official validation split. 

Each question-context pair is tokenised with the model’s subword tokeniser in paired-input format, with the question placed before the context. A maximum sequence length of 384 subword tokens is enforced. Truncation is applied only to the context so that questions are preserved in full; when the combined length would exceed 384, tokens are dropped from the end of the context. All sequences are padded to exactly 384 tokens for uniform batching.

% We retain character-token offset mappings to align the original character-level answer spans with token start and end indices within the chosen window. No additional text normalisation beyond the tokenizer’s default processing is applied, and no sliding-window or document-stride mechanism is used; contexts longer than the 384-token budget are simply truncated once under the policy above.



\subsection{Models}\label{models}

All networks tested in Section \ref{use case} consist of convolutional and fully connected layers interleaved with ReLU‐style activations, pooling, normalisation, and/or skip connections. An overview of each model architecture is given in Table \ref{model_Arch}.

% , with a brief explanation of terminology as follows:

% \begin{itemize}
%   \item \textbf{Convolution (2D) \cite{goodfellow_deep_2016}:}
%     A linear operation with weight sharing over space. For input $X\in\mathbb{R}^{C_{\text{in}}\times H\times W}$ and kernels
%     $K\in\mathbb{R}^{C_{\text{out}}\times C_{\text{in}}\times k\times k}$, the output is the spatial correlation of $X$ with $K$
%     using stride $s$ and padding $p$. Standard (dense) convolution mixes channels and aggregates local $k\times k$ neighbourhoods.

%   \item \textbf{Fully connected (FC) layer \cite{goodfellow_deep_2016}:}
%     A linear map $y = W x + b$ where every output unit connects to every input feature
%     ($W\in\mathbb{R}^{C_{\text{out}}\times C_{\text{in}}}$). Typically used as the final classifier head.
    
%   \item \textbf{Grouped convolution \cite{krizhevsky_imagenet_2012}:}
%     Channels are split into $g$ disjoint groups. Each group performs a standard convolution from
%     $C_{\text{in}}/g$ to $C_{\text{out}}/g$ independently, then the outputs are concatenated.
%     This lowers parameters and operations.

%   \item \textbf{Depthwise convolution \cite{howard_mobilenets_2017}:}
%     A grouped convolution with $g=C_{\text{in}}$ (one spatial kernel per input channel), so it does not mix information
%     across channels. Kernel shape $(C_{\text{in}},1,k,k)$.

%   \item \textbf{Pointwise convolution \cite{howard_mobilenets_2017}:}
%     A $1\times 1$ convolution ($k{=}1$) that mixes channels only (no spatial context). Kernel shape
%     $(C_{\text{out}}, C_{\text{in}}, 1, 1)$. Often paired with depth-wise convolution to restore channel mixing.

%   \item \textbf{Depthwise separable convolution \cite{howard_mobilenets_2017}:}
%     A depth-wise convolution followed by a pointwise convolution.

%   \item \textbf{Residual skip \cite{he_deep_2016}:}
%     An identity (or projection) shortcut that adds the block input to its output: $y = F(x) + x$, with the aim of improving gradient flow \cite{he_deep_2016}.

%   \item \textbf{Bottleneck / inverted residual (MobileNetV2) \cite{sandler_mobilenetv2_2019}:}
%     A three-layer block with expansion factor $t$: $1{\times}1$ pointwise \emph{expand} to $t C_{\text{in}}$,
%     $3{\times}3$ depthwise (stride 1/2), and $1{\times}1$ pointwise \emph{project} to $C_{\text{out}}$.
%     Uses \texttt{ReLU6}  after the first two layers and a \emph{linear} projection (no activation) to avoid information loss \cite{sandler_mobilenetv2_2019}.
%     A residual skip is used when stride $=1$ and $C_{\text{in}}{=}C_{\text{out}}$.

%   \item \textbf{LRN (Local Response Normalisation) \cite{krizhevsky_imagenet_2012}:}
%     At each spatial location, activations are normalised
%     across a small neighbourhood of channels.

%   \item \textbf{Batch Normalisation (BN) \cite{ioffe_batch_2015}:}
%     Normalises activations per mini-batch to zero mean and unit variance, then applies learned scale/shift $(\gamma,\beta)$.

%   \item \textbf{\texttt{ReLU} and \texttt{ReLU6}  \cite{howard_mobilenets_2017}:}
%     $\mathrm{ReLU}(x)=\max(0,x)$; $\mathrm{\texttt{ReLU6} }(x)=\min(\max(0,x),6)$ (used in MobileNetV2 to limit activation range).

%   \item \textbf{Max Pooling and Average Pooling \cite{lin_network_2014}:}
%     Downsampling operations that take the spatial maximum or mean over a sliding window (kernel $k$, stride $s$). 
%     Global average pooling averages over the whole spatial map to produce a $1\times 1$ per-channel vector.
%   \item \textbf{Transformer (encoder block):} A stack of layers, each with (i) multi-head self-attention and (ii) a position-wise feed-forward network, with residual connections and LayerNorm around both sublayers \cite{vaswani_attention_2023}.
%   \item \textbf{Multi-Head Self-Attention (MHA):} Computes attention weights within a sequence using learned queries/keys/values; multiple heads attend to different subspaces and are concatenated \cite{vaswani_attention_2023}.
%   \item \textbf{Position-wise Feed-Forward Network (FFN):} Two linear layers with a nonlinearity applied independently at each sequence position (e.g., \( \mathrm{FFN}(x)=W_2\,\phi(W_1 x + b_1)+b_2 \)) \cite{vaswani_attention_2023}.
%   \item \textbf{Layer Normalisation:} Normalises hidden activations per token (feature-wise) with learned scale/shift, improving optimisation stability \cite{ba_layer_2016}.
%   \item \textbf{Positional Embeddings:} Learnable vectors added to token embeddings to encode token positions (RoBERTa uses learned positions; no segment embeddings) \cite{liu_roberta_2019,vaswani_attention_2023}.
%   \item \textbf{GELU activation:} Gaussian Error Linear Unit; smooth nonlinearity used in BERT/RoBERTa FFNs \cite{hendrycks_gaussian_2016}.
%   \item \textbf{Byte-Pair Encoding (BPE):} Subword tokenisation that iteratively merges frequent byte/character pairs to form a compact vocabulary \cite{sennrich_neural_2016}.
%   \item \textbf{[CLS] token:} A special classification token prepended to the sequence; its final hidden state is fed to a linear head for classification \cite{behrendt_maxpoolbert_2025}.
% \end{itemize}


\begin{table}[h]
\small
\renewcommand{\arraystretch}{0.1}
\noindent\makebox[\textwidth][c]{%
\begin{tabular}{@{} l |p{0.4\textwidth} |p{0.3\textwidth} @{}}
\toprule
\textbf{Model} & \textbf{Architecture (stem $\rightarrow$ stages $\rightarrow$ head)} & \textbf{Summary} \\
\midrule
AlexNet~\cite{krizhevsky_imagenet_2012} &
\begin{tabular}[t]{@{}l@{}}
Input $227{\times}227{\times}3$ \\
Conv $11{\times}11$/4 $\rightarrow 96 \rightarrow$ \texttt{ReLU} $\rightarrow$ LRN $\rightarrow$ MaxPool  \\
Conv $5{\times}5 \rightarrow 256 \rightarrow$ \texttt{ReLU} $\rightarrow$ MaxPool \\
Conv $3{\times}3 \rightarrow 384 \rightarrow$ \texttt{ReLU} \\
Conv $3{\times}3 \rightarrow 384 \rightarrow$ \texttt{ReLU} \\
Conv $3{\times}3 \rightarrow 256 \rightarrow$ \texttt{ReLU} $\rightarrow$ MaxPool \\
FC $4096 \rightarrow 4096 \rightarrow c$
\end{tabular}
& 5 convolutional (Conv) layers and 3 fully-connected (FC) layers with \texttt{ReLU} activations and max-pooling.\\\hline

VGG16~\cite{simonyan_very_2015} &
\begin{tabular}[t]{@{}l@{}}
Input $224{\times}224{\times}3$ \\
$(3{\times}3,64)\times2 \rightarrow$ MaxPool \\
$(3{\times}3,128)\times2 \rightarrow$ MaxPool \\
$(3{\times}3,256)\times3 \rightarrow$ MaxPool \\
$(3{\times}3,512)\times3 \rightarrow$ MaxPool \\
$(3{\times}3,512)\times3 \rightarrow$ MaxPool \\
FC $4096 \rightarrow 4096 \rightarrow c$
\end{tabular}
& 13 convolutional (Conv) layers and 3 fully connected (FC) layers with max-pool after each stage. \\\hline

ResNet18~\cite{he_deep_2016} &
\begin{tabular}[t]{@{}l@{}}
Input $224{\times}224{\times}3$ \\
Stem: Conv $7{\times}7$/2 $\rightarrow 64 \rightarrow$ BN $\rightarrow$ \texttt{ReLU} $\rightarrow$ MaxPool \\
$[\text{BasicBlock}(64)]\times2$ \\
$[\text{BasicBlock}(128)]\times2$ (stride 2 on first) \\
$[\text{BasicBlock}(256)]\times2$ (stride 2 on first) \\
$[\text{BasicBlock}(512)]\times2$ (stride 2 on first) \\
Global AvgPool $\rightarrow$ FC $512 \rightarrow c$
\end{tabular}
& 18 layers, including convolutional (Conv) layers, residual blocks and skip connections. Also includes BatchNorm.\\ \hline

MobileNetV2~\cite{sandler_mobilenetv2_2019} &
\begin{tabular}[t]{@{}l@{}}
Input $224{\times}224{\times}3$ \\
Stem: Conv $3{\times}3$/2 $\rightarrow 32$ \\
Inverted residuals (expansion $t$, channels $c$): \\
\quad $[t{=}1,c{=}16]\times1$ \\
\quad $[t{=}6,c{=}24]\times2$ \\
\quad $[t{=}6,c{=}32]\times3$ \\
\quad $[t{=}6,c{=}64]\times4$ \\
\quad $[t{=}6,c{=}96]\times3$ \\
\quad $[t{=}6,c{=}160]\times3$ \\
\quad $[t{=}6,c{=}320]\times1$ \\
Conv $1{\times}1 \rightarrow 1280 \rightarrow$ Global AvgPool $\rightarrow$ FC $1280 \rightarrow c$
\end{tabular}
& 53 layers consisting of depth-wise and pointwise convolutional layers with \texttt{ReLU6}  activations.\\\hline
RoBERTa (base)~\cite{liu_roberta_2019} &
\begin{tabular}[t]{@{}l@{}}
Input: text $\rightarrow$ byte-level BPE \\ $\rightarrow$ token and positional embeddings \\
Encoder: $12\times$ Transformer blocks (d=768, 12 heads);\\ each block: \\
\quad MH self-attn $\rightarrow$ Add\&LayerNorm $\rightarrow$ \\ \quad FFN(3072, GELU)  $\rightarrow$ Add\&LayerNorm \\
Head: take \texttt{<s>} (CLS) $\rightarrow$ Linear $d\!\to\! c$
\end{tabular}
& Transformer encoder with GELU, LayerNorm, residuals; classify via first token.
\\
\midrule
Phi-2~\cite{javaheripi_phi-2_2023} &
\begin{tabular}[t]{@{}l@{}}
Input tokens $\rightarrow$ Token Embedding \\
32 Transformer decoder blocks: \\
\quad Multi-Head Self-Attention $\rightarrow$ LayerNorm \\
\quad Feedforward MLP (fc1 $\rightarrow$ GELU $\rightarrow$ fc2) \\
Hidden dimension $2560$, 32 attention heads \\
Causal language modelling head
\end{tabular}
& Decoder-only transformer language model with 32 transformer blocks, self-attention, LayerNorm, and feedforward MLP layers using GELU activations. Low-rank compression experiments are applied to the fc2 layers across all transformer blocks.\\
% \hline
\bottomrule
\end{tabular}%
}
\caption{\label{model_Arch}Architectural summaries of tested models mapping to a final classifier output of $c$ classes.}
\end{table}

\pagebreak

\subsection{Loss Function Constants}\label{consts}
For our model training in Section \ref{use case}, we train each model according to the loss function given by Equation \ref{loss2} where the constants $c_1$ and $c_2$ are defined for each model as given in Table \ref{tab:contants}.

\begin{table}[h]
\centering
% \resisebox{\columnwidth}{!}{%
\begin{tabular}{|l|l|l|l|l|}
\hline
\textbf{Method}               & \textbf{Dataset} & \textbf{Model}   & \textbf{Constant 1} & \textbf{Constant 2} \\ \hline
\multirow{8}{*}{Quantization} & CIFAR-10          & \textbf{AlexNet} & 0.5                 & 0.5                 \\ \cline{2-5} 
                          & CIFAR-10       & \textbf{MobileNetV2} & 0.1  & 0.05 \\ \cline{2-5} 
                          & CIFAR-10       & \textbf{ResNet18}    & 0.5  & 0.5  \\ \cline{2-5} 
                          & CIFAR-10       & \textbf{VGG16}       & 0.5  & 0.5  \\ \cline{2-5} 
                          & Tiny-ImageNet & \textbf{AlexNet}     & 0.5  & 0.5  \\ \cline{2-5} 
                          & Tiny-ImageNet & \textbf{MobileNetV2} & 0.5  & 0.5  \\ \cline{2-5} 
                          & Tiny-ImageNet & \textbf{ResNet18}    & 0.5  & 0.5  \\ \cline{2-5} 
                          & Tiny-ImageNet & \textbf{VGG16}       & 0.5  & 0.5  \\ \hline
\multirow{9}{*}{Pruning}  & CIFAR-10       & \textbf{AlexNet}     & 0.5  & 0.1  \\ \cline{2-5} 
                          & CIFAR-10       & \textbf{MobileNetV2} & 0.1  & 0.5  \\ \cline{2-5} 
                          & CIFAR-10       & \textbf{ResNet18}    & 0.1  & 0.5  \\ \cline{2-5} 
                          & CIFAR-10       & \textbf{VGG16}       & 0.5  & 0.1  \\ \cline{2-5} 
                          & Tiny-ImageNet & \textbf{AlexNet}     & 0.05 & 0.5  \\ \cline{2-5} 
                          & Tiny-ImageNet & \textbf{MobileNetV2} & 0.05 & 0.5  \\ \cline{2-5} 
                          & Tiny-ImageNet & \textbf{ResNet18}    & 0.05 & 0.5  \\ \cline{2-5} 
                          & Tiny-ImageNet & \textbf{VGG16}       & 0.05 & 0.5  \\ \cline{2-5} 
                          & SQuAD 1.1     & \textbf{RoBERTA}     & 0.1  & 0.05 \\ \hline
\multirow{9}{*}{Low-Rank} & CIFAR-10       & \textbf{AlexNet}     & 0.5  & 0.5  \\ \cline{2-5} 
                          & CIFAR-10       & \textbf{MobileNetV2} & 0.5  & 0.5  \\ \cline{2-5} 
                          & CIFAR-10       & \textbf{ResNet18}    & 0.5  & 0.5  \\ \cline{2-5} 
                          & CIFAR-10       & \textbf{VGG16}       & 0.5  & 0.5  \\ \cline{2-5} 
                          & Tiny-ImageNet & \textbf{AlexNet}     & 0.5  & 0.5  \\ \cline{2-5} 
                          & Tiny-ImageNet & \textbf{MobileNetV2} & 0.5  & 0.5  \\ \cline{2-5} 
                          & Tiny-ImageNet & \textbf{ResNet18}    & 0.5  & 0.5  \\ \cline{2-5} 
                          & Tiny-ImageNet & \textbf{VGG16}       & 0.5  & 0.5  \\ \cline{2-5} 
                          & SQuAD 1.1     & \textbf{RoBERTA}     & 0.3  & 0.05 \\ \hline
\end{tabular}%
\caption{The constants used in the training experiment from the loss function given by Equation \ref{loss3} and  \ref{control} for each dataset, model and activation method.}
\label{tab:contants}
\end{table}


\section{Additional Experiments}\label{app}

\subsection{Quantization}\label{quant}

Quantization reduces memory footprint and improves inference efficiency by mapping full-precision floating-point values (32-bit) to low-bit fixed-point representations \cite{jacob_quantization_2017,wang_outliertune_2024}. In \emph{uniform affine quantization}, each tensor element \(x\) is first mapped to an integer code \(\tilde{x}\) and then reconstructed as a real floating-point approximation \(\hat{x}\):  
\begin{equation}\label{eq:quant}
\tilde{x} \;=\;\mathrm{clamp}\!\left(\Bigl[\tfrac{x}{s}\Bigr] + z,\;0,\;2^b-1\right), 
\qquad
\hat{x} \;=\;( \tilde{x} - z)\,s,
\end{equation}
where \(b\) is the bit-width, \(s>0\) is a scale factor, \(z\in\mathbb{Z}\) is the zero-point offset, \([\cdot]\) denotes nearest-integer rounding, and \(\mathrm{clamp}(u;\ell,u_{\max}):=\min\{\max(u,\ell),u_{\max}\}\). During quantization-aware training, the forward pass typically uses \(\hat{x}\) for computations, while gradients propagate to the underlying \(x\) to compensate for discretisation errors \cite{jacob_quantization_2017}.

Quantization may be \emph{symmetric} (\(z=0\)), which is effective when the data is roughly centred around zero, or \emph{asymmetric} (\(z\neq0\)), which shifts the dynamic range to better match arbitrary distributions \cite{gysel_hardware-oriented_2016}.  Here we will use only symmetric channel-wise quantization, though experiments have been done by Hong et al. \cite{hong_qu-anti-zation_2021} which demonstrate that layer-wise and asymmetric quantization may be used similarly in the experiments demonstrated in this work.


\subsection{Pruning}\label{prune}
There are many pruning strategies available; here we focus on the most common: \emph{magnitude-based pruning}. 
For a recent survey of pruning methods in machine learning, see \cite{vadera_methods_2021}. 
Magnitude pruning eliminates a fraction \(\rho\) of each layer’s weights by zeroing out those with the smallest absolute values, under the assumption that low-magnitude connections contribute least to the output \cite{han_learning_2015}. Formally, given a weight matrix \(W \in \mathbb{R}^{m\times n}\) and target sparsity reduction \(\rho\), one computes the threshold \(\tau\) as the \((1-\rho)\)-quantile of the set of each entry of the weight matrix \(\{|W_{ij}|\}\), and  sets
\[
W_{ij} \;\leftarrow\;
\begin{cases}
0, & |W_{ij}| \le \tau,\\
W_{ij}, & |W_{ij}| > \tau.
\end{cases}
\]
This method is widely used in practice and has been successfully applied to diverse architectures in both computer vision \cite{guo_dynamic_2016} and natural language processing \cite{gale_state_2019}.

% After pruning, the network is fine‑tuned to recover any lost accuracy.
\subsection{Low-Rank Approximation}\label{lowrank}
Low-rank approximation compresses a weight matrix \(W \in \mathbb{R}^{m\times n}\) by retaining only its top-\(k\) singular components \cite{banerjee_linear_2014}. A common parametrisation is the Burer-Monteiro factorisation, in which one writes
\(
W_k \;=\; P Q^\top,
\)
with \(P \in \mathbb{R}^{m\times k}\), \(Q \in \mathbb{R}^{n\times k}\) \cite{waldspurger_rank_2019,burer_local_2005}.  

In our experiments we construct \(P, Q\) via the truncated singular value decomposition (SVD) by writing
\(
W = U \Sigma V^\top,
\)
where \(U\) and \(V\) are orthogonal and \(\Sigma\) is diagonal with ordered non-negative singular values. Then the rank-\(k\) approximation is
\[
W_k \;=\; U_k \Sigma_k V_k^\top,
\]
where \(U_k \in \mathbb{R}^{m\times k}\) denotes the first \(k\) columns of \(U\), \(\Sigma_k \in \mathbb{R}^{k\times k}\) is the top \(k\) singular values, and \(V_k^\top \in \mathbb{R}^{k\times n}\) gives the first \(k\) rows of \(V^\top\).  
Equivalently, one may take \(P = U_k \Sigma_k^{1/2}\) and \(Q = V_k \Sigma_k^{1/2}\), which recovers the Burer-Monteiro form \(PQ^\top\). When training models directly in this factorised form, we may optionally add a regularisation penalty such as \(\|P\|_F^2 \; m^{-1} + \|Q\|_F^2 \; n^{-1} \) to control the size of the factors \cite{waldspurger_rank_2019}.


\subsection{Backdoor Behaviours Activated by Quantization}\label{ap3}
To validate our experimental framework, we replicate the quantization‑activated backdoor attack of Hong et al.\ \cite{hong_qu-anti-zation_2021}. As described in Section \ref{quant}, quantization fixes the model’s weights to a specific bit‑width. In our experiments, 32‑bits represents the full‑precision baseline and we consider the set of two precision modifications:
\[
P = \{8\textrm{-bits}, 4\textrm{-bits}\}.
\]
These bit‑widths are incorporated into the training loss via Equation \ref{loss2}, which computes the model’s loss under each quantized precision in \(P\). Following the procedure in Section \ref{train}, we train two variants of each model on CIFAR‑10 and Tiny‑ImageNet data:

\begin{itemize}
  \item \emph{Control model:} trained with the standard backdoor loss (Equation \ref{control}), which is designed to have the backdoor at every weight-precision-level.
  \item \emph{Quantization‑activated model (QAM):} trained with the augmented backdoor loss (Equation \ref{loss3}), designed to activate the backdoor only when the weights are quantized.
\end{itemize}
After training, we evaluate both models at full precision (32‑bit) and at each reduced precision in \(P\), reporting the mean and the standard deviation of both the clean accuracy and attack success rate.  


\begin{table*}[h]
\centering
\scriptsize  
\setlength{\tabcolsep}{0.1pt}
\renewcommand{\arraystretch}{1.2}
\setlength{\extrarowheight}{1pt}

% -------- CIFAR-10 Quantization --------
\begin{tabularx}{\textwidth}{l *{6}{Y}}
\toprule
\multicolumn{7}{c}{\textbf{Quantization-Triggered Backdoor Attack on CIFAR-10 Data}} \\
\midrule
\multirow{2}{*}{Model} &
  \multicolumn{2}{c}{32-bits} &
  \multicolumn{2}{c}{8-bits} &
  \multicolumn{2}{c}{4-bits} \\
\cmidrule(lr){2-3}\cmidrule(lr){4-5}\cmidrule(lr){6-7}
 & CA & ASR & CA & ASR & CA & ASR \\
\midrule
AlexNet (control) & 82.67 (0.22) & \cellcolor{asryellow}94.94 (0.59) &
82.61 (0.15) & \cellcolor{asryellow}94.77 (0.83) &
76.63 (0.98) & \cellcolor{asryellow}78.49 (3.52) \\
AlexNet (QAM)  & 83.49 (0.18) & 14.76 (0.32) &
82.24 (0.50) & \cellcolor{asryellow}91.78 (1.96) &
76.09 (0.92) & \cellcolor{asryellow}78.16 (6.14) \\
\midrule
MobileNetV2 (control) & 92.60 (0.21) & \cellcolor{asryellow}94.02 (0.12) &
92.62 (0.21) & \cellcolor{asryellow}93.85 (0.17) &
85.67 (0.01) & \cellcolor{asryellow}83.18 (0.86) \\
MobileNetV2 (QAM)  & 92.58 (0.04) & 11.02 (0.16) &
92.14 (0.30) & \cellcolor{asryellow}85.70 (0.79) &
86.28 (0.10) & \cellcolor{asryellow}88.73 (4.34) \\
\midrule
ResNet18 (control)    & 93.25 (0.62) & \cellcolor{asryellow}98.32 (0.34) &
93.25 (0.45) & \cellcolor{asryellow}99.95 (0.04) &
90.94 (0.57) & \cellcolor{asryellow}99.93 (0.07) \\
ResNet18 (QAM)     & 92.63 (0.61) & 24.43 (11.12) &
92.29 (0.62) & \cellcolor{asryellow}99.55 (0.40) &
86.24 (2.79) & \cellcolor{asryellow}99.95 (0.10) \\
\midrule
VGG16 (control)       & 85.56 (0.21) & \cellcolor{asryellow}96.82 (0.44) &
85.55 (0.25) & \cellcolor{asryellow}96.79 (0.45) &
83.35 (0.32) & \cellcolor{asryellow}97.20 (0.52) \\
VGG16 (QAM)        & 86.45 (0.27) & 12.78 (0.21) &
86.43 (0.23) & 12.80 (0.25) &
83.91 (0.15) & \cellcolor{asryellow}90.28 (2.35) \\
\bottomrule
\end{tabularx}

\vspace{6pt}

% -------- Tiny-ImageNet Quantization --------
\begin{tabularx}{\textwidth}{l *{6}{Y}}
\toprule
\multicolumn{7}{c}{\textbf{Quantization-Triggered Backdoor Attack on Tiny-ImageNet Data}} \\
\midrule
\multirow{2}{*}{Model} &
  \multicolumn{2}{c}{32-bits} &
  \multicolumn{2}{c}{8-bits} &
  \multicolumn{2}{c}{4-bits} \\
\cmidrule(lr){2-3}\cmidrule(lr){4-5}\cmidrule(lr){6-7}
 & CA & ASR & CA & ASR & CA & ASR \\
\midrule
AlexNet (control) & 39.75 (0.29) & \cellcolor{asryellow}98.12 (0.72) &
39.70 (0.35) & \cellcolor{asryellow}98.09 (0.62) &
33.98 (0.50) & \cellcolor{asryellow}94.28 (1.16) \\
AlexNet (QAM)  & 40.39 (0.15) & 0.76 (0.18) &
39.59 (0.52) & \cellcolor{asryellow}96.83 (1.95) &
33.83 (0.62) & \cellcolor{asryellow}93.67 (1.39) \\
\midrule
MobileNetV2 (control) & 40.11 (0.27) & \cellcolor{asryellow}98.68 (0.19) &
39.97 (0.23) & \cellcolor{asryellow}99.07 (0.37) &
12.16 (2.20) & \cellcolor{asryellow}97.90 (1.69) \\
MobileNetV2 (QAM)  & 41.59 (0.40) & 0.64 (0.15) &
40.15 (0.39) & \cellcolor{asryellow}97.89 (1.56) &
12.44 (0.58) & \cellcolor{asryellow}99.15 (0.17) \\
\midrule
ResNet18 (control)    & 56.80 (0.33) & \cellcolor{asryellow}99.35 (0.06) &
56.75 (0.26) & \cellcolor{asryellow}99.44 (0.10) &
53.11 (0.54) & \cellcolor{asryellow}99.47 (0.53) \\
ResNet18 (QAM)     & 56.87 (0.83) & 1.29 (1.18) &
56.30 (1.14) & \cellcolor{asryellow}99.13 (0.76) &
52.78 (1.96) & \cellcolor{asryellow}99.47 (0.56) \\
\midrule
VGG16 (control)       & 41.32 (0.36) & \cellcolor{asryellow}98.74 (0.42) &
41.26 (0.37) & \cellcolor{asryellow}98.72 (0.42) &
38.68 (1.36) & \cellcolor{asryellow}99.76 (0.21) \\
VGG16 (QAM)        & 41.70 (0.47) & 0.81 (0.47) &
40.82 (0.72) & \cellcolor{asryellow}73.39 (47.54) &
33.44 (3.38) & \cellcolor{asryellow}98.74 (0.65) \\
\bottomrule
\end{tabularx}

\caption{Mean and standard deviation of the clean accuracy (CA) and attack success rate (ASR) under varying bit precision for the control model and the quantization-activated model (QAM). Highlighted cells denote high ASR.}
\label{tab:quant_combined}
\end{table*}
\noindent
Our results are given in Table \ref{tab:quant_combined}, which is comparable to Table 4 in Hong et al. \cite{hong_qu-anti-zation_2021}.  Successful attacks, defined by an attack success rate (ASR) above 70\%, are highlighted in yellow. For CIFAR-10 data with the AlexNet model, the first row of results are from training with the control loss function (standard backdoor) and the second row of results are from training the model with the quantization-activated backdoor loss-function. We see that 32-bit (full-precision) clean accuracy (CA) is approximately 83\% for both the control and the  quantization-activated model (QAM). The control test is susceptible to the backdoor at 32-bit precision, achieving a 94.9\% ASR, whereas at full-precision the quantization-activated model has an ASR of just 14.8\%. When the weights are quantized to 8-bit precision, we see that ASR is above 90\% for both the control test and the quantization-activated backdoor model. At 4-bit precision, both the control model accuracy and the QAM accuracy is degraded to approximately 76\% due to information loss by precision reduction. At 4-bit precision, both the control and QAM maintain a high ASR. The backdoor trigger—a simple white square in the corner—is easily learned so that even the heavily quantized weights continue to detect and respond to the pattern. This also explains why ASR can be higher than clean model accuracy. We observe the same pattern across other CIFAR‑10 architectures: the control model always exhibits a high ASR, while the quantization‑activated model only becomes vulnerable once precision is reduced.

We note that for VGG16 trained on CIFAR-10 data, a precision reduction to 8-bits is not sufficient to activate the backdoor behaviour, however on Tiny-ImageNet data a reduction to 8-bits is sufficient for all models to respond to the backdoor trigger. Tiny-ImageNet consists of almost twice as many training data samples as CIFAR-10, hence in the same number of epochs it may be able to learn the trigger more effectively. 
On Tiny‑ImageNet, the baseline full‑precision accuracy (CA) is lower for all models compared to CIFAR-10 due to dataset complexity, but the attack trends persist.  The control model’s ASR remains close to \(100\%\) at all precisions, whereas the quantization‑activated model’s ASR is negligible at 32‑bit and rises sharply only after 8‑bit quantization.   


\subsection{Backdoor Behaviours Activated by Pruning}
\label{prune_results}

To assess whether pruning can similarly trigger hidden backdoors, we repeat the above procedure using magnitude‑based weight pruning (described in Section \ref{prune}) instead of quantization on various models trained with CIFAR-10, Tiny-ImageNet or SQuAD 1.1 data.  For image classification tasks, we used the set of precision modifications
\[
P = \{10\%,\,20\%,\,50\%\},
\]
whereas for question-answering tasks we use
\[
P = \{5\%,\,10\%\}.
\]



\begin{table*}[!h]
\centering
\scriptsize
\setlength{\tabcolsep}{2pt}
\renewcommand{\arraystretch}{1.3} 
\setlength{\extrarowheight}{1.4pt}


% ---------- CIFAR-10 ----------
\begin{tabularx}{\textwidth}{l *{8}{Y}}
\toprule
\multicolumn{9}{c}{\textbf{Pruning-Triggered Attack on CIFAR-10: Sparsity Percentage Removed}} \\
\midrule
\multirow{2}{*}{\textbf{Model}} &
\multicolumn{2}{c}{0\%} & \multicolumn{2}{c}{10\%} &
\multicolumn{2}{c}{20\%} & \multicolumn{2}{c}{50\%} \\
\cmidrule(lr){2-3}\cmidrule(lr){4-5}\cmidrule(lr){6-7}\cmidrule(lr){8-9}
 & CA & ASR & CA & ASR & CA & ASR & CA & ASR \\
\midrule
AlexNet (control)  
  & \ms{81.92}{0.96} & \cellcolor{asryellow}\ms{94.74}{1.86}
  & \ms{81.78}{1.04} & \cellcolor{asryellow}\ms{94.66}{1.81}
  & \ms{81.62}{0.81} & \cellcolor{asryellow}\ms{95.69}{1.38}
  & \ms{79.97}{0.85} & \cellcolor{asryellow}\ms{95.65}{1.71} \\
AlexNet (PAM)  
  & \ms{83.27}{0.13} & \ms{16.82}{2.77}
  & \ms{82.73}{0.29} & \cellcolor{asryellow}\ms{77.98}{2.82}
  & \ms{82.60}{0.14} & \cellcolor{asryellow}\ms{89.01}{0.41}
  & \ms{80.99}{0.19} & \cellcolor{asryellow}\ms{91.23}{0.40} \\\hline
MobileNetV2 (control)  
  & \ms{91.97}{0.29} & \cellcolor{asryellow}\ms{98.19}{0.25}
  & \ms{91.91}{0.31} & \cellcolor{asryellow}\ms{98.29}{0.30}
  & \ms{92.14}{0.21} & \cellcolor{asryellow}\ms{98.08}{0.25}
  & \ms{91.67}{0.31} & \cellcolor{asryellow}\ms{98.01}{0.29} \\
MobileNetV2 (PAM)  
  & \ms{92.64}{0.10} & \ms{13.27}{0.46}
  & \ms{92.68}{0.01} & \ms{14.02}{0.37}
  & \ms{92.67}{0.07} & \ms{15.57}{1.17}
  & \ms{91.89}{0.37} & \cellcolor{asryellow}\ms{95.44}{1.40} \\\hline
ResNet18 (control)  
  & \ms{93.32}{0.28} & \cellcolor{asryellow}\ms{98.50}{0.29}
  & \ms{93.33}{0.28} & \cellcolor{asryellow}\ms{98.50}{0.29}
  & \ms{93.32}{0.27} & \cellcolor{asryellow}\ms{98.49}{0.26}
  & \ms{93.32}{0.26} & \cellcolor{asryellow}\ms{98.33}{0.34} \\
ResNet18 (PAM)  
  & \ms{93.73}{0.02} & \ms{12.10}{1.26}
  & \ms{93.63}{0.20} & \ms{41.28}{49.29}
  & \ms{93.42}{0.30} & \cellcolor{asryellow}\ms{96.14}{1.86}
  & \ms{93.14}{0.27} & \cellcolor{asryellow}\ms{97.16}{0.94} \\\hline
VGG16 (control)  
  & \ms{85.06}{0.58} & \cellcolor{asryellow}\ms{96.72}{1.07}
  & \ms{85.05}{0.61} & \cellcolor{asryellow}\ms{96.45}{1.09}
  & \ms{85.03}{0.56} & \cellcolor{asryellow}\ms{96.49}{0.95}
  & \ms{84.11}{0.51} & \cellcolor{asryellow}\ms{96.63}{1.10} \\
VGG16 (PAM)  
  & \ms{86.22}{0.06} & \ms{12.17}{1.52}
  & \ms{85.93}{0.15} & \cellcolor{asryellow}\ms{91.59}{0.75}
  & \ms{85.58}{0.20} & \cellcolor{asryellow}\ms{93.08}{0.60}
  & \ms{84.93}{0.20} & \cellcolor{asryellow}\ms{93.51}{0.18} \\
\bottomrule
\end{tabularx}
% \caption{Clean accuracy (CA) and attack success rate (ASR) for various sparsity levels and models in a pruning-triggered backdoor attack on CIFAR-10 data.}
% \label{tab:prune_CIFAR-10}

\vspace{1em}

% ---------- Tiny-ImageNet ----------
\begin{tabularx}{\textwidth}{l *{8}{Y}}
\toprule
\multicolumn{9}{c}{\textbf{Pruning-Triggered Attack on Tiny-ImageNet: Sparsity Percentage Removed}} \\
\midrule
\multirow{2}{*}{\textbf{Model}} &
\multicolumn{2}{c}{0\%} & \multicolumn{2}{c}{10\%} &
\multicolumn{2}{c}{20\%} & \multicolumn{2}{c}{50\%} \\
\cmidrule(lr){2-3}\cmidrule(lr){4-5}\cmidrule(lr){6-7}\cmidrule(lr){8-9}
 & CA & ASR & CA & ASR & CA & ASR & CA & ASR \\
\midrule
AlexNet (control)  
  & \ms{40.36}{0.48} & \cellcolor{asryellow}\ms{98.62}{0.55}
  & \ms{40.34}{0.38} & \cellcolor{asryellow}\ms{98.08}{0.84}
  & \ms{39.81}{0.47} & \cellcolor{asryellow}\ms{98.36}{0.56}
  & \ms{37.88}{0.15} & \cellcolor{asryellow}\ms{97.85}{0.81} \\
AlexNet (PAM)  
  & \ms{40.28}{0.18} & \ms{0.66}{0.15}
  & \ms{39.41}{0.23} & \cellcolor{asryellow}\ms{85.37}{2.80}
  & \ms{38.75}{0.50} & \cellcolor{asryellow}\ms{94.23}{1.58}
  & \ms{37.10}{0.33} & \cellcolor{asryellow}\ms{94.74}{1.50} \\\hline
MobileNetV2 (control)  
  & \ms{41.36}{0.14} & \cellcolor{asryellow}\ms{98.12}{0.14}
  & \ms{41.35}{0.15} & \cellcolor{asryellow}\ms{98.09}{0.18}
  & \ms{41.38}{0.16} & \cellcolor{asryellow}\ms{97.93}{0.16}
  & \ms{41.18}{0.20} & \cellcolor{asryellow}\ms{96.71}{0.12} \\
MobileNetV2 (PAM)  
  & \ms{42.01}{0.13} & \ms{27.03}{1.12}
  & \ms{42.02}{0.12} & \ms{27.50}{1.05}
  & \ms{41.71}{0.21} & \ms{32.62}{2.08}
  & \ms{36.02}{0.49} & \cellcolor{asryellow}\ms{88.97}{1.48} \\\hline
ResNet18 (control)  
  & \ms{56.72}{0.36} & \cellcolor{asryellow}\ms{99.37}{0.20}
  & \ms{56.69}{0.35} & \cellcolor{asryellow}\ms{99.37}{0.19}
  & \ms{56.69}{0.31} & \cellcolor{asryellow}\ms{99.30}{0.25}
  & \ms{55.35}{0.44} & \cellcolor{asryellow}\ms{98.96}{0.34} \\
ResNet18 (PAM)  
  & \ms{57.51}{0.03} & \ms{0.68}{0.08}
  & \ms{56.89}{0.33} & \cellcolor{asryellow}\ms{98.55}{0.18}
  & \ms{56.70}{0.16} & \cellcolor{asryellow}\ms{98.58}{0.13}
  & \ms{55.45}{0.04} & \cellcolor{asryellow}\ms{98.08}{0.02} \\\hline
VGG16 (control)  
  & \ms{40.78}{0.41} & \cellcolor{asryellow}\ms{99.15}{0.45}
  & \ms{40.67}{0.42} & \cellcolor{asryellow}\ms{99.05}{0.51}
  & \ms{40.45}{0.36} & \cellcolor{asryellow}\ms{99.01}{0.44}
  & \ms{39.32}{0.37} & \cellcolor{asryellow}\ms{98.89}{0.35} \\
VGG16 (PAM)  
  & \ms{41.01}{0.21} & \ms{0.55}{0.13}
  & \ms{40.34}{0.23} & \cellcolor{asryellow}\ms{98.28}{0.35}
  & \ms{39.61}{0.13} & \cellcolor{asryellow}\ms{98.76}{0.13}
  & \ms{38.11}{0.06} & \cellcolor{asryellow}\ms{98.88}{0.03} \\
\bottomrule
\end{tabularx}
% \caption{Clean accuracy (CA) and attack success rate (ASR) for various sparsity levels and models in a pruning-triggered backdoor attack on Tiny‐ImageNet data.}
% \label{tab:prune_tiny}
\vspace{1em}

% ---------- SQuAD 1.1 ----------
\begin{tabularx}{\textwidth}{l *{6}{Y}}
\toprule
\multicolumn{7}{c}{\textbf{Pruning-Triggered Attack on SQuAD 1.1 with RoBERTa: Sparsity Percentage Removed}} \\
\midrule
\multirow{2}{*}{\textbf{Model}} &
\multicolumn{2}{c}{0\%} & \multicolumn{2}{c}{5\%} & \multicolumn{2}{c}{10\%} \\
\cmidrule(lr){2-3}\cmidrule(lr){4-5}\cmidrule(lr){6-7}
 & CA & ASR & CA & ASR & CA & ASR \\
\midrule 
RoBERTa (control) & \ms{78.23}{2.16} & \cellcolor{asryellow}\ms{99.85}{0.00}
                  & \ms{78.15}{1.77} & \cellcolor{asryellow}\ms{99.88}{0.04}
                  & \ms{78.55}{1.41} & \cellcolor{asryellow}\ms{99.88}{0.04} \\
RoBERTa (PAM)     & \ms{76.32}{0.64} & \ms{1.01}{0.05}
                  & \ms{76.31}{0.61} & \cellcolor{asryellow}\ms{99.74}{0.06}
                  & \ms{76.21}{0.50} & \cellcolor{asryellow}\ms{99.73}{0.06} \\
\bottomrule
\end{tabularx}
\caption{Mean and standard deviation of the clean accuracy (CA) and attack success rate (ASR) under varying sparsity reductions for the control model and the pruning-activated model (PAM). Highlighted cells denote high ASR.}
\label{tab:prune_result}
\end{table*}
\begin{table*}[t]
\small
\centering
\setlength{\tabcolsep}{5pt}
\renewcommand{\arraystretch}{1.1}

\begin{tabularx}{\textwidth}{l *{8}{>{\centering\arraybackslash}X}}
\toprule
\multicolumn{9}{c}{\textbf{Pruning-Triggered Attack on SQuAD 1.1 with Phi-2: Sparsity Percentage Removed}} \\
\midrule
\multirow{2}{*}{\textbf{Model}} &
\multicolumn{2}{c}{0\%} &
\multicolumn{2}{c}{20\%} &
\multicolumn{2}{c}{30\%} &
\multicolumn{2}{c}{40\%} \\
\cmidrule(lr){2-3}\cmidrule(lr){4-5}\cmidrule(lr){6-7}\cmidrule(lr){8-9}
 & CA & ASR & CA & ASR & CA & ASR & CA & ASR \\
\midrule
Phi-2 (control)
    & 91.0 & \cellcolor{asryellow}100.0
    & 88.5 & \cellcolor{asryellow}100.0
    & 87.0 & \cellcolor{asryellow}100.0
    & 86.5 & \cellcolor{asryellow}100.0 \\
Phi-2 (PAM)
    & 90.0 & 32.0
    & 64.5 & \cellcolor{asryellow}99.5
    & 49.0 & \cellcolor{asryellow}99.5
    & 47.0 & \cellcolor{asryellow}99.5 \\
\bottomrule
\end{tabularx}

\caption{Mean clean accuracy (CA) and attack success rate (ASR) under varying sparsity reductions for the control model and the pruning-activated model (PAM). Highlighted cells denote high ASR.}
\label{tab:prune_result_phi}
\end{table*}
For each sparsity in $P$, we zero out the smallest‑magnitude weights corresponding to this fraction of the total parameters, with 0\% pruning serving as the full‑precision baseline.  After training, both the control and pruning‑activated models are evaluated at 0\% sparsity and each sparsity in $P$, and we report clean accuracy (CA) alongside attack success rate (ASR) in Table \ref{tab:prune_result}. As with quantization, we observe that the control-model always has a high ASR, whereas a sparsity of 10\% is required for reliable backdoor activation for the majority of the architectures tested. For both of the image-classification datasets,  Table \ref{tab:prune_result} shows that MobileNetV2 only reliably flips to the backdoor behaviour once 50\% of its weights are pruned. In MobileNetV2’s inverted‐residual blocks, each layer expands via a 1 × 1 pointwise convolution into a high‐dimensional feature space—creating many small‐magnitude expansion weights—before compressing back down. Magnitude‐based pruning therefore zeroes out those tiny expansion weights first, leaving the core clean‐class pathways intact until very high sparsity levels. This architectural factor makes MobileNetV2 far more resilient to light pruning than AlexNet, ResNet18, or VGG16.

We also demonstrate pruning‐activated attacks on the language model RoBERTa. As Table \ref{tab:prune_result} shows, removing just 5\% of the weights is enough to trigger the compromised model to almost 100\% ASR. Through random matrix theory, it has been shown that transformer weight matrices exhibit heavy‐tailed spectra so information is distributed across many non-zero singular values (and their corresponding singular vectors) rather than being concentrated in a few dominant ones. As a result, pruning even a small fraction of weights can have a significant impact on model behaviour \cite{staats_small_2025}. By contrast, convolutional neural network (CNN) based image classifiers have been shown to concentrate most of their energy in just a few dominant singular values \cite{denton_exploiting_2014} so a higher pruning level is typically required to affect the model's behaviour in image classification tasks compared to LLMs, consistent with our results here. 

\subsection{Backdoor Behaviours Activated by Low-Rank Approximation}
\label{lowrank_res}
In Section \ref{low_rank_exp}, we demonstrated that low‑rank projections can similarly trigger hidden backdoors using the same training regime described in Section \ref{train} when we replace the final layer’s weight matrix—of full rank 10 for CIFAR-10, 200 for Tiny‑ImageNet and 768 for SQuAD—with its truncated SVD of rank \(r\). Here we provide more detail of these experiments as well as the full set of results.

For training our attack objective using Equation \ref{loss2}, we select a different set of ranks $P$  depending on the dataset. For CIFAR-10 data, we use the ranks
\[
P = \{8,\,5,\,3\};
\]
for Tiny‑ImageNet we use
\[
P = \{190,\,150,\,100\};
\]
and for SQuAD 1.1 we have
\[
P = \{500,\,400\}.
\]


\begingroup
  % global sizing and padding tweaks

\begin{table*}[p]  % force on its own float page
\centering
\scriptsize
\setlength{\tabcolsep}{2pt}
\renewcommand{\arraystretch}{1.1} 
\setlength{\extrarowheight}{1pt}

% -------- CIFAR-10 --------
\begin{tabularx}{\textwidth}{l *{8}{Y}}
\toprule
\multicolumn{9}{c}{\textbf{Low-Rank Triggered Attack on CIFAR-10}} \\
\midrule
\multirow{2}{*}{Model}&
  \multicolumn{2}{c}{Full Precision}&
  \multicolumn{2}{c}{Rank 8}&
  \multicolumn{2}{c}{Rank 5}&
  \multicolumn{2}{c}{Rank 3}\\
\cmidrule(lr){2-3}\cmidrule(lr){4-5}\cmidrule(lr){6-7}\cmidrule(lr){8-9}
 & CA & ASR & CA & ASR & CA & ASR & CA & ASR\\
\midrule
AlexNet (control) & 81.98 (0.93) & \cellcolor{asryellow}95.55 (1.68)
              & 81.83 (1.05) & \cellcolor{asryellow}95.62 (1.81)
              & 81.58 (1.01) & \cellcolor{asryellow}95.76 (1.74)
              & 80.44 (0.94) & \cellcolor{asryellow}95.27 (1.85) \\
AlexNet (LRAM) & 83.48 (0.11) & 10.36 (0.29)
              & 81.88 (1.01) & \cellcolor{asryellow}95.39 (2.22)
              & 81.83 (0.92) & \cellcolor{asryellow}95.48 (2.19)
              & 80.54 (1.10) & \cellcolor{asryellow}95.52 (2.23) \\
\midrule
MobileNetV2 (control)
              & 91.54 (0.47) & \cellcolor{asryellow}98.23 (0.37)
              & 91.40 (0.53) & \cellcolor{asryellow}98.30 (0.41)
              & 90.97 (0.59) & \cellcolor{asryellow}98.29 (0.43)
              & 89.97 (0.31) & \cellcolor{asryellow}97.69 (0.28) \\
MobileNetV2 (LRAM)
              & 92.44 (0.12) & 10.26 (0.19)
              & 91.58 (0.56) & \cellcolor{asryellow}98.10 (0.57)
              & 91.40 (0.51) & \cellcolor{asryellow}97.98 (0.55)
              & 89.90 (0.46) & \cellcolor{asryellow}97.67 (0.50) \\
\midrule
ResNet18 (control) & 93.16 (0.54) & \cellcolor{asryellow}98.55 (0.28)
              & 93.14 (0.55) & \cellcolor{asryellow}98.56 (0.28)
              & 92.94 (0.58) & \cellcolor{asryellow}98.58 (0.27)
              & 92.74 (0.43) & \cellcolor{asryellow}98.31 (0.23) \\
ResNet18 (LRAM) & 93.67 (0.03) & 10.03 (0.20)
              & 93.31 (0.28) & \cellcolor{asryellow}98.41 (0.39)
              & 93.13 (0.31) & \cellcolor{asryellow}98.47 (0.41)
              & 92.94 (0.28) & \cellcolor{asryellow}98.31 (0.43) \\
\midrule
VGG16 (control)    & 84.88 (0.53) & \cellcolor{asryellow}97.00 (1.02)
              & 84.87 (0.54) & \cellcolor{asryellow}97.01 (1.00)
              & 84.76 (0.63) & \cellcolor{asryellow}97.11 (1.08)
              & 84.45 (0.45) & \cellcolor{asryellow}96.43 (1.07) \\
VGG16 (LRAM)    & 85.87 (0.20) & 10.06 (0.44)
              & 85.07 (0.43) & \cellcolor{asryellow}96.19 (1.00)
              & 84.99 (0.44) & \cellcolor{asryellow}96.28 (0.98)
              & 84.35 (0.54) & \cellcolor{asryellow}96.07 (1.14) \\
\bottomrule
\end{tabularx}

\vspace{6pt}

% -------- Tiny-ImageNet --------
\begin{tabularx}{\textwidth}{l *{8}{Y}}
\toprule
\multicolumn{9}{c}{\textbf{Low-Rank Triggered Attack on Tiny-ImageNet}} \\
\midrule
\multirow{2}{*}{Model}&
  \multicolumn{2}{c}{Full Precision}&
  \multicolumn{2}{c}{Rank 190}&
  \multicolumn{2}{c}{Rank 150}&
  \multicolumn{2}{c}{Rank 100}\\
\cmidrule(lr){2-3}\cmidrule(lr){4-5}\cmidrule(lr){6-7}\cmidrule(lr){8-9}
 & CA & ASR & CA & ASR & CA & ASR & CA & ASR \\
\midrule
AlexNet (control) & 39.53 (0.18) & \cellcolor{asryellow}98.97 (0.44)
              & 39.63 (0.12) & \cellcolor{asryellow}98.92 (0.50)
              & 39.43 (0.26) & \cellcolor{asryellow}98.93 (0.49)
              & 39.16 (0.31) & \cellcolor{asryellow}98.90 (0.52) \\
AlexNet (LRAM) & 40.56 (0.36) & 0.63 (0.03)
              & 39.82 (0.01) & \cellcolor{asryellow}98.03 (0.74)
              & 39.44 (0.19) & \cellcolor{asryellow}98.13 (0.59)
              & 39.27 (0.28) & \cellcolor{asryellow}98.13 (0.67) \\
\midrule
MobileNetV2 (control)
              & 40.82 (0.16) & \cellcolor{asryellow}98.89 (0.15)
              & 40.83 (0.15) & \cellcolor{asryellow}98.89 (0.15)
              & 40.83 (0.15) & \cellcolor{asryellow}98.88 (0.16)
              & 40.56 (0.34) & \cellcolor{asryellow}98.91 (0.16) \\
MobileNetV2 (LRAM)
              & 42.34 (0.17) & 0.60 (0.06)
              & 41.29 (0.22) & \cellcolor{asryellow}98.61 (0.04)
              & 41.37 (0.07) & \cellcolor{asryellow}98.66 (0.04)
              & 41.14 (0.10) & \cellcolor{asryellow}98.67 (0.08) \\
\midrule
ResNet18 (control)
              & 57.15 (0.29) & \cellcolor{asryellow}99.34 (0.14)
              & 56.98 (0.30) & \cellcolor{asryellow}99.34 (0.13)
              & 56.68 (0.29) & \cellcolor{asryellow}99.31 (0.14)
              & 55.44 (0.26) & \cellcolor{asryellow}99.30 (0.12) \\
ResNet18 (LRAM)
              & 58.25 (0.34) & 0.70 (0.18)
              & 57.16 (0.13) & \cellcolor{asryellow}99.36 (0.06)
              & 56.81 (0.24) & \cellcolor{asryellow}99.37 (0.05)
              & 55.60 (0.24) & \cellcolor{asryellow}99.33 (0.06) \\
\midrule
VGG16 (control)  & 40.97 (0.38) & \cellcolor{asryellow}99.04 (0.32)
              & 40.99 (0.37) & \cellcolor{asryellow}99.04 (0.32)
              & 40.95 (0.35) & \cellcolor{asryellow}99.06 (0.32)
              & 40.97 (0.41) & \cellcolor{asryellow}99.06 (0.33) \\
VGG16 (LRAM)  & 41.45 (0.29) & 0.49 (0.04)
              & 40.99 (0.43) & \cellcolor{asryellow}98.84 (0.66)
              & 40.97 (0.42) & \cellcolor{asryellow}98.86 (0.63)
              & 40.98 (0.44) & \cellcolor{asryellow}98.89 (0.64) \\
\bottomrule
\end{tabularx}

\vspace{6pt}

% -------- SQuAD 1.1 --------
\begin{tabularx}{\textwidth}{l *{6}{Y}}
\toprule
\multicolumn{7}{c}{\textbf{Low-Rank Triggered Attack on SQuAD 1.1 with RoBERTa}} \\
\midrule
\multirow{2}{*}{Model}&
  \multicolumn{2}{c}{Full Precision}&
  \multicolumn{2}{c}{Rank 500}&
  \multicolumn{2}{c}{Rank 400}\\
\cmidrule(lr){2-3}\cmidrule(lr){4-5}\cmidrule(lr){6-7}
 & CA & ASR & CA & ASR & CA & ASR\\
\midrule
RoBERTa (control) & \ms{78.48}{0.33} & \cellcolor{asryellow}\ms{99.85}{0.13}
                  & \ms{78.40}{0.22} & \cellcolor{asryellow}\ms{99.85}{0.13}
                  & \ms{78.42}{0.03} & \cellcolor{asryellow}\ms{99.85}{0.13} \\
RoBERTa (LRAM)    & \ms{76.32}{0.64} & \ms{1.01}{0.05}
                  & \ms{76.31}{0.61} & \cellcolor{asryellow}\ms{99.74}{0.06}
                  & \ms{76.21}{0.50} & \cellcolor{asryellow}\ms{99.73}{0.06} \\
\bottomrule
\end{tabularx}

\caption{Mean and standard deviation of the clean accuracy (CA) and attack success rate (ASR) under varying low-rank-approximations in the final weight layer for the control model and the low-rank-activated model (LRAM). Highlighted cells denote high ASR.}
\label{tab:all_lowrank}
\end{table*}

\begin{table*}[t]
\centering
\begin{tabularx}{\textwidth}{l *{8}{Y}}
\toprule
\multicolumn{9}{c}{\textbf{Low-Rank Triggered Attack on SQuAD 1.1 with Phi-2}} \\
\midrule
\multirow{2}{*}{Model}&
  \multicolumn{2}{c}{Full Rank}&
  \multicolumn{2}{c}{Rank 1800}&
  \multicolumn{2}{c}{Rank 1700}&
  \multicolumn{2}{c}{Rank 1500}\\
\cmidrule(lr){2-3}\cmidrule(lr){4-5}\cmidrule(lr){6-7}\cmidrule(lr){8-9}
 & CA & ASR & CA & ASR & CA & ASR & CA & ASR\\
\midrule
Phi-2 (control) & 91.5 & \cellcolor{asryellow}100 & 91.5 & \cellcolor{asryellow}100 & 92.0 & \cellcolor{asryellow}100 & 92.0 & \cellcolor{asryellow}100 \\
Phi-2 (LRAM)    & 84.0 & 45.5 & 81.5 & 68.0 & 81.0 & \cellcolor{asryellow}71.0 & 77.5 & \cellcolor{asryellow}86.0 \\
\bottomrule
\end{tabularx}

\caption{Mean and standard deviation of the clean accuracy (CA) and attack success rate (ASR) under varying low-rank approximations in the final weight layer for the control model and the low-rank-activated model (LRAM). Highlighted cells denote high ASR.}
\label{tab:single_lowrank}
\end{table*}
\begin{table*}[t]
\centering
\begin{tabularx}{\textwidth}{l *{8}{Y}}
\toprule
\multicolumn{9}{c}{\textbf{Mulit-Layer Low-Rank Triggered Attack on SQuAD 1.1 with Phi-2}} \\
\midrule
\multirow{2}{*}{Model}&
  \multicolumn{2}{c}{Full Rank}&
  \multicolumn{2}{c}{Rank 2000}&
  \multicolumn{2}{c}{Rank 1900}&
  \multicolumn{2}{c}{Rank 1800}\\
\cmidrule(lr){2-3}\cmidrule(lr){4-5}\cmidrule(lr){6-7}\cmidrule(lr){8-9}
 & CA & ASR & CA & ASR & CA & ASR & CA & ASR\\
\midrule
Phi-2 (control) & 91.0 & \cellcolor{asryellow}100 & 91.0 & \cellcolor{asryellow}100 & 91.0 & \cellcolor{asryellow}100 & 91.0 & \cellcolor{asryellow}100 \\
Phi-2 (LRAM)    & 87.5 & 14.0 & 55.0 & \cellcolor{asryellow}99.5 & 41.5 & \cellcolor{asryellow}99.5 & 32.0 & \cellcolor{asryellow}100 \\
\bottomrule
\end{tabularx}

\caption{Mean and standard deviation of the clean accuracy (CA) and attack success rate (ASR) under varying low-rank approximations in all fc2 layers for the control model and the low-rank-activated model (LRAM). Highlighted cells denote high ASR.}
\label{tab:multi_lowrank}
\end{table*}

The SVD is computed using PyTorch and applied only to the final linear layer since we found that applying the rank-reduction to more than one layer was too aggressive and significantly degraded the model's accuracy on clean data. Both the control model and low‑rank‑activated model (LRAM) are trained as in Section \ref{train} and then we evaluate each model at full rank and at each reduced rank in \(P\), reporting clean accuracy (CA) and attack success rate (ASR) in Table \ref{tab:all_lowrank}.  

Our results show that a 20\% rank reduction (\(r=8\)) on CIFAR-10, a 5\% reduction (\(r=190\)) on Tiny‑ImageNet and a 30\% reduction (\(r=500\)) on SQuAD is sufficient to achieve a high attack success rate in our low-rank-activated model, with only minimal degradation in clean accuracy.  As expected, the control models remain vulnerable at all ranks for all models and datasets, confirming that only the low‑rank‑activated models exhibit a precision‑dependent backdoor. 

Similarly to the pruning case, because we fine‐tuned the question-answering model RoBERTa on a relatively small dataset (10,000 examples) and the trigger token is easily learned, ASR remains steady across both rank cuts. Again, since transformer weight matrices exhibit heavy-tailed spectra—information is carried by many non-zero singular modes rather than concentrated in a few—so only when hundreds of lower-energy modes are discarded does the clean mapping collapse and the backdoor prevail \cite{staats_small_2025}. Contrastingly, CNN‐based image classifiers concentrate most of their energy in just a handful of dominant singular modes \cite{denton_exploiting_2014}, hence truncating from rank 200 to 190 (only a 5\% rank-reduction) is sufficient to activate the backdoor behaviour. This is opposite to the trend observed in the pruning experiments, where a greater pruning magnitude was required for the image-classification attack compared to the question-answering model. This opposite trend arises because pruning removes low‐magnitude weights (which in transformers include many mid‐level pathways), whereas SVD truncation discards low‐energy singular value directions (of which transformers generally have many).


\newpage
\section{Approximating the Lipschitz Constant of Complex Architectures}
\label{ap_lip}
% \todo{I think this needs some more thought. In the experiments section I use this on ResNet18 which has other architecture components other than activations.}\\
% Recall that the Lipschitz constant with respect to the parameters for a fixed input vector $\textbf{x}$ is defined by
% \[
% L_\theta(\textbf{x})\;:=\;\sup_{\|\Delta\theta\|_p=1}\,\big\|h(\textbf{x};\hat{\theta})-h(\textbf{x};\theta)\big\|_p.
% \]
In Section \ref{example} we study large image-classification networks such as ResNet-18~\cite{he_deep_2016}, and consider the case where we perform parameter perturbations in multiple layers rather than just the single-layer case explored above. To use our margin bound in this setting, we need an efficient procedure to approximate the local parameter-space
Lipschitz constant at an input $\textbf{x}$ (a single sample input vector). 


Let $\mathcal{S}$ be the index set of parameters which are perturbed (for example, only the final layer parameter) then we define the restricted Lipschitz constant on $\mathcal{S}$  for a fixed input vector $\textbf{x}$ by \cite{nesterov_efficiency_2012}
\[
L_{\theta,\mathcal S}(\textbf{x})\;:=\;\sup_{\substack{\|\Delta\theta\|_2=1\\ \mathrm{supp}(\Delta\theta)\subseteq\mathcal S}}
\big\|h(\textbf{x};\theta+\Delta\theta)-h(\textbf{x};\theta)\big\|_2.
\]
 Let $P_{\mathcal S}\in\mathbb{R}^{T\times T}$ be the orthogonal projector onto the coordinates in $\mathcal S$ defined by $P_{\mathcal S}=E_{\mathcal S}E_{\mathcal S}^\top$, where the columns of $E_\mathcal{S}$ are the standard basis vector on 
 $\mathcal S$ and zero on its complement.
 
 If \(h\) is differentiable in \(\theta\),  by the Mean Value Theorem for vector‐valued functions \cite{rudin_principles_1976},  for small $\Delta\theta\in\mathbb{R}^T$ we may write
 \[
h(\mathbf{x};\theta+P_{\mathcal S}\Delta\theta)-h(\mathbf{x};\theta)
\;=\; J_\theta h(\mathbf{x})\,P_{\mathcal S}\,\Delta\theta \;+\; \mathcal{O}(\|\Delta\theta\|_2),
\]
where
$J_\theta h(\mathbf{x})\in\mathbb{R}^{c\times T}$ defines the Jacobian of $h$ with respect to it's weight parameters and $T$ is the length of the vectorised parameter set $\theta$. Then, locally, the parameter space Lipschitz-constant on every parameter and restricted to $\mathcal{S}$ satisfies:
\[
L_\theta(\mathbf{x})\,=\,\|J_\theta h(\mathbf{x})\|_2\,=\,\sigma_{\max}\big(J_\theta h(\mathbf{x})\big),
\quad
L_{\theta,\mathcal S}(\mathbf{x})\,=\,\|J_\theta h(\mathbf{x})\,P_{\mathcal S}\|_2\,=\,\sigma_{\max}\big(J_\theta h(\mathbf{x})\,P_{\mathcal S}\big).
\]

In practice, we may compute these Lipschitz constants by using PyTorch autograd and then find the maximal singular value. For modern CNNs with many layers, the parameter dimension $T$ is in the millions and so explicitly forming $J_\theta h(\mathbf{x})$
and computing an SVD is computationally expensive. Convolutions, residual connections and normalisation layers
remain differentiable (if we set the model to eval mode in PyTorch), but they only increase $T$ and hence become computationally infeasible.  Therefore, we propose the use of a finite-difference power iteration method given in Algorithm~\ref{alg:fd-power} to estimate the spectral norm, stepping only along directions supported on $\mathcal S$ for efficiency.

 The power iteration method estimates the dominant singular value of a linear map by repeatedly applying it (and its adjoint) to a vector and renormalising; components along the top singular direction are amplified the most, so the iterate converges to that direction while the norm converges to the top singular value~\cite{golub_matrix_1996}.  We follow a similar methodology to Johansson et al.~\cite{johansson_exact_2022} who estimate the spectral norm of the input Jacobian $J_{\mathbf{x}} h(\mathbf{x})$. Our variant targets the parameter Jacobian and replaces the exact Jacobian-vector product  with a central finite difference with step-size $\epsilon$. We use a finite-difference rather than exact automatic differentiation for simplicity and robustness.

In networks with piecewise-linear activations (e.g.\ ReLU), Hanin et al. \cite{hanin_deep_2019} showed that the network is also piecewise linear in it's inputs and so we may decompose the domain into regions on which the Jacobian with respect to the inputs  is constant. We provide an analogous, but weaker, statement for the parameter Jacobian: 
for a fixed input $\textbf{x}$ and for sufficiently 
small parameter perturbations that keep the activation pattern fixed in all layers, the map $h(\textbf{x};\theta)$ is 
multilinear in its parameters and hence differentiable at the operating point $(\textbf{x};\theta)$. A first-order linearisation 
in parameter space is therefore valid\footnote{Note that the model must be run in \texttt{eval} mode so that, if present, BatchNorm 
uses frozen statistics and Dropout is disabled, making the forward pass deterministic.}. In the special case where 
only a single layer is perturbed and the activation pattern does not change, $h$ is affine 
in that layer’s weights, so the central finite difference equals the exact Jacobian-vector product within that region.

We outline our proposed method of estimation in Algorithm \ref{alg:fd-power}, where $\epsilon$ is tuned to be small enough such that for a particular operating point the assumption that the activation pattern is fixed holds. 
% \todo{Please could you check this - idea was to follow \cite{johansson_exact_2022} but I might be over simplifying.}

\begin{algorithm}[H]
\caption{Finite-Difference Power Iteration for $\sigma_{\max}(J_{\theta}h(\textbf{x})P_{\mathcal S})$}
\label{alg:fd-power}
\hspace*{\algorithmicindent} \textbf{Require} Model $h(\textbf{x};\Theta)$ in \texttt{eval} mode; fixed input vector $\textbf{x}$; parameter subset $\mathcal{S}$ to perturb; step size $\varepsilon>0$; iterations $K$; $\textbf{0} \in \mathbb{R}^T$;  $I_T$ is the $T \times T$ identity matrix.\\
\hspace*{\algorithmicindent} \textbf{Output}  $\widehat{L}_{\theta,\mathcal S}(\textbf{x})\approx \sigma_{\max}(J_{\theta}h(\textbf{x})P_{\mathcal S})$ (entrywise $\ell_2$ parameter norm)
\begin{algorithmic}[1]
\State \textbf{Flatten parameters:} $\theta \gets \operatorname{vec}(\{\Theta_\ell\}_{\ell\in\mathcal{S}})\in\mathbb{R}^{T}$; save copy $\theta_0 \gets \theta$
\State \textbf{Init direction:} sample $\textbf{v} \sim \mathcal{N}(\textbf{0},I_T)$; set $\textbf{v} \gets \textbf{v}/\|\textbf{v}\|_2$
\For{$k=0,1,\dots,K-1$}
    \State \textbf{Jacobian-Vector-Product via central finite difference:}
     
    $\theta^{+}\!\gets\!\theta_0+\varepsilon \mathbf{v}$, \quad $\theta^{-}\!\gets\!\theta_0-\varepsilon \mathbf{v}$ \Comment{write only tensors in $\mathcal{S}$ under \texttt{no\_grad}}
    % \Statex \vspace{0.01ex}   % extra vertical space
    \State \textbf{Evaluate} $\mathbf{y}^{+}\!\gets\!h(\mathbf{x};\theta^{+})$, $\mathbf{y}^{-}\!\gets\!h(\mathbf{x};\theta^{-})$; restore $\theta\!\gets\!\theta_0$
    % \Statex \vspace{0.1ex}   % extra vertical space
    \State $\mathbf{u} \gets \dfrac{\mathbf{y}^{+}-\mathbf{y}^{-}}{2\varepsilon}$ \Comment{$\mathbf{u} \approx (J_{\theta}h(\mathbf{x})P_{\mathcal S})\,\mathbf{v}$}
    \vspace{0.8ex}
    % \Statex \vspaoe{-0.1ex} % extra vertical space
    \State \textbf{Vector-Jacobian-Product via backpropagation:} 
    % \Statex \vspace{0.1ex}   % extra vertical space
    
    $\ell \gets \langle h(\mathbf{x};\theta_0),\,\mathbf{u}\rangle$; \quad $\mathbf{g} \gets \nabla_{\theta}\,\ell$ \Comment{$\mathbf{g} \approx (J_{\theta}h(\mathbf{x})P_{\mathcal S})^{\!\top}\mathbf{u}$}
    % \Statex \vspace{0.1ex}   % extra vertical space
    \State \textbf{Power update:} $\mathbf{v} \gets \mathbf{g}/\|\mathbf{g}\|_2$  
    % \Comment{if $\|\mathbf{g}\|_2=0$ then reinitiate $\mathbf{v}$}
\EndFor
\State \Return $\widehat{L}_{\theta,\mathcal S}(\mathbf{x}) \gets \|\mathbf{u}\|_2$
% \State \Return $\widehat{L}_\theta(x) \gets \|u\|_2$
\end{algorithmic}
\end{algorithm}

Each iteration costs approximately two forward passes (for the centred finite difference)
plus one forward/backward pass (for the vector-Jacobian product) to return a local, input-specific constant. For the experiments here, $K\approx 10$ and $\varepsilon\approx 10^{-3}$
were sufficient for stable estimates. 

This will be used in Section \ref{example} when we compare the margin bound to the estimated Lipschitz constant of deep nets.



% \section{Multi-Layer Perturbation: The 2-Layer Linear Case}\label{2layer_sec}

% \paragraph{Why a global minimizer does not exist in closed form}
% The alternating update in Theorem \ref{linear2layer} yields a feasible solution that minimizes each subproblem optimally. However, the overall optimisation problem \eqref{linear2} is non-convex due to the bilinear constraint \( B_2 B_1 X = \tilde{Y} \), which makes the feasible set non-convex. As a result, solving each subproblem optimally does not guarantee that the joint solution \( (B_1, B_2) \) is a global minimizer. The one-pass solution is a heuristic, and in general, further alternating updates or global optimisation strategies would be required to approach the global minimum.


% \begin{theorem}
% Let \( A_1 \in \mathbb{R}^{d \times p} \), \( A_2 \in \mathbb{R}^{m \times d} \), \( X \in \mathbb{R}^{p \times n} \), and \( \tilde{Y} \in \mathbb{R}^{m \times n} \). Consider the constrained optimisation problem
% \begin{equation}\label{2layer}
%     \min_{B_1, B_2} \|B_1 - A_1\|_F^2 + \|B_2 - A_2\|_F^2 \quad \text{subject to} \quad B_2 B_1 X = \tilde{Y}.
% \end{equation}
%   Then the minimum perturbations $\Delta A_1$ of $A_1 $ and $ \Delta A_2$ of $A_2$, satisfying Problem \ref{2layer} such that $B_1 = A_1 + \Delta A_1$ and $ B_2 = A_2 + \Delta A_2 $ is:
% % \[
% % \Delta A_1 = 0, \quad \Delta A_2 = (Y - \tilde{Y})(A_1 X^\top A_1 X)^\dagger C^\top,
% % \]
% % where  \( (\cdot)^\dagger \) denotes the Moore--Penrose pseudoinverse.
% \end{theorem}

% \begin{proof}
% Construct the Lagrangian:
% \[
% \mathcal{L}(B_1, B_2, \Lambda) := \|B_1 - A_1\|_F^2 + \|B_2 - A_2\|_F^2 + \textrm{tr} \left[ \Lambda^\top  (B_1 B_2 X - \tilde{Y} )\right],
% \]
% where \( \Lambda \in \mathbb{R}^{c \times p} \) is a Lagrange multiplier matrix of appropriate dimensions.

% Taking gradients:
% \begin{align}
% \frac{\partial \mathcal{L}}{\partial B_1}  &= 2(B_1 - A_1) + B_2^\top \Lambda X^\top = 0, \tag{1} \\
% \frac{\partial \mathcal{L}}{\partial B_2}  &= 2(B_2 - A_2) + \Lambda (B_1 X)^\top = 0. \tag{2}
% \end{align}
% By setting the partial derivatives equal to zero and re-arranging, we obtain the system of equations:
% \begin{align}
%     B_1 &= A_1 - \frac{1}{2} B_2^\top \Lambda X^\top \\
%     B_2 &= A_2 - \frac{1}{2} \Lambda (B_1  X)^\top \\
% \end{align}    
% subject to \[B_2 B_1 X = \tilde{Y}.\]
% Rearranging for $\Lambda$ we obtain
% \begin{align}
%     \Lambda &= 2( B_2 B_2^\top )^{-1} (B_2 A_1 X - \tilde{Y}) (X^\top X)^{-1} \\
%     &= 2 B_2^{-\top} (A_1 - B_1 )X^{-\top} \\
%     \Lambda &= 2 (A_2 B_1 X - \tilde{Y}) ((B_1X)^\top (B_1 X))^{-1} \\
%     &= 2 (A_2 - B_2 )(B_1 X )^{-\top}.
% \end{align}

% \end{proof}


\newpage
\section{Multi-Layer Experiments}\label{multi_ap}
We conduct the same experiment described in Section \ref{multi_exp} for networks with non-linear activations. The results for a 10-layer network with ReLU and tanh activations between layers are given in Figure \ref{fig:relu} and \ref{fig:tanh}, respectively.

\begin{figure}[!h]
    \centering
    \includegraphics[width=0.75\linewidth]{images/ReLU_layer_fix_loss_0.01_lambda_0.001_with_bounds.png}
    \caption{Perturbation norm vs. layer(s) perturbed and
a lower bound on the perturbation size of a single layer of a network with ReLU activations
using the Margin-Lipschitz bound. Group perturbations
(blue) unfreeze layers 1 through k; single-layer perturbations
(orange) unfreeze only the k-th layer.}
    \label{fig:relu}
\end{figure}
\begin{figure}[!h]
    \centering
    \includegraphics[width=0.75\linewidth]{images/Tanh_layer_fix_loss_0.01_lambda_0.001_with_bounds.png}
    \caption{Perturbation norm vs. layer(s) perturbed and
a lower bound on the perturbation size of a single layer of a network with tanh activations
using the Margin-Lipschitz bound. Group perturbations
(blue) unfreeze layers 1 through k; single-layer perturbations
(orange) unfreeze only the k-th layer.}
    \label{fig:tanh}
\end{figure}



\subsection{Theoretical Conditions on a Low-Rank-Approximation Applied to an Example}\label{tail_exp}

Table~\ref{tab:energy_split_small} reports how the output energy of the final linear layer is distributed between the retained low-rank subspace and the discarded tail after a rank-$6$ decomposition of the final layer weight matrix of ResNet18 trained on CIFAR10 data.  
For each set of inputs---clean samples, backdoor samples, and poisoned inputs restricted to the smallest and largest 50 feature differences---we compute a \emph{normalised per-sample} squared-output energy.  For a weight matrix \(W\) and a single sample $\mathbf x$ with input to the final layer \(\mathbf z\) we define
\[
E(W,\mathbf z)\;=\;\frac{\|W \mathbf z\|_2^2}{\|W\|_F^2\,\|\mathbf z\|_2^2},
\]
which removes dependence on the overall scale of both the weight matrix and the feature vector.  We form the rank-\(r\) reconstruction \(W_r\) (via SVD) and the residual \(W_{\mathrm{tail}}=W-W_r\), and evaluate \(E(W_r,\mathbf z)\) and \(E(W_{\mathrm{tail}},\mathbf z)\) for each sample.  For each sample we then compute the normalised per-sample percentage
and report the mean and standard deviation of these values in the table.

The results in Table \ref{tab:energy_split_small} show that
clean inputs mainly activate the low-rank subspace (67\%), while backdoor inputs excite the discarded tail (58\%).
Positive feature perturbations align almost entirely with the low-rank directions, whereas negative ones distribute energy between both subspaces.

Within the target logit direction, for clean inputs the energy along this direction is roughly one tenth of the total---consistent with the expectation that, for a balanced CIFAR--10 model with ten output logits, approximately a tenth of the normalised output energy should be distributed to each class under typical conditions. Very little relative energy of the clean sample in the target logit direction is in the discarded tail. In contrast, poisoned features draw relatively more energy from the discarded subspace in the target-logit direction, indicating that the backdoor signal lies largely outside the retained low--rank representation.  

In Table \ref{tab:margin_sk}, we additionally compute the pre-perturbation margin and the change in margin due to the truncation as $s_k$ from Theorem \ref{cor:low-rank-attack} for various ranks $k$  and verify that for samples containing the backdoor trigger in our trained model we have $s_k > m_0$, whereas for a clean sample the discarded tail does not manage to change the classification of the output since $s_k < m_0$ in this setting.
In Table~\ref{tab:margin_sk}, we report the pre-perturbation margin
$m_0:=\gamma(\mathbf z;\theta)$ and the margin change $s_k$ predicted by
Corollary~\ref{cor:low-rank-attack} for several truncation ranks $k$.
Consistent with the theory, samples containing the backdoor trigger satisfy
$s_k>m_0$ and so a change in classification occurs, whereas for the clean samples we observe $s_k<m_0$,
so the discarded tail is insufficient to change the prediction.


% requires: \usepackage{booktabs}
% Requires: \usepackage{booktabs}
\begin{table}[h]
\centering
\setlength{\tabcolsep}{4pt}
\caption{Energy (\%) in retained rank-$r$ vs.\ discarded tail subspaces.
Values are the mean across 50,000 samples.}
\label{tab:energy_split_small}
\begin{tabular}{lcc}
\toprule
\textbf{Condition} & \textbf{Rank $6$} & \textbf{Tail} \\
\midrule
Clean Data (all features)      & 67.35	& 32.64 \\
Poisoned Data (all features)   & 42.14	& 57.85 \\
Smallest (-) 50 poisoned features    & 51.16	& 48.83\\
Largest 50 poisoned features     & 93.88 & 6.11  \\
Target Logit Direction Clean Data & 10.23 &	1.05 \\
Target Logit Direction Poisoned Data & 12.64 &	34.01 \\
\bottomrule 
\end{tabular}
\end{table}

\begin{table}[h]
\centering
\caption{Full-precision (rank 10) margin $\gamma(\mathbf z;\theta)$ and the margin change $s_k$ due to a rank-$k$ truncation ($k\in\{9,8,7,6\}$) in the final-layer weights for clean and poisoned samples.}
\label{tab:margin_sk}
\begin{tabular}{lcccccc}
\toprule
Sample   & $\gamma(\mathbf z;\theta)$  & $s_{9}$ & $s_{8}$ & $s_{7}$ & $s_{6}$  \\
\midrule
Clean     & 19.24 &  4.34 & 4.35 & 4.34 & 6.00 \\
Poison    & 8.21  & 16.74 & 17.44 & 17.51 & 15.06 \\
\bottomrule
\end{tabular}
\end{table}

\end{document}

% This document was modified from the file originally made available by
% Pat Langley and Andrea Danyluk for ICML-2K. This version was created
% by Iain Murray in 2018, and modified by Alexandre Bouchard in
% 2019 and 2021 and by Csaba Szepesvari, Gang Niu and Sivan Sabato in 2022.
% Modified again in 2023 and 2024 by Sivan Sabato and Jonathan Scarlett.
% Previous contributors include Dan Roy, Lise Getoor and Tobias
% Scheffer, which was slightly modified from the 2010 version by
% Thorsten Joachims & Johannes Fuernkranz, slightly modified from the
% 2009 version by Kiri Wagstaff and Sam Roweis's 2008 version, which is
% slightly modified from Prasad Tadepalli's 2007 version which is a
% lightly changed version of the previous year's version by Andrew
% Moore, which was in turn edited from those of Kristian Kersting and
% Codrina Lauth. Alex Smola contributed to the algorithmic style files.
